% This file is an experimental variant of dynamic-algebra-potentials.tex
% where the tensor convention in Theta_z is reversed:
%   Theta_z(c^-/c^+) = [c^-, z] tensor c^+    (function-from-input on the left)
% instead of the chosen
%   Theta_z(c^-/c^+) = c^+ tensor [c^-, z]    (output on the left).
%
% The output-first convention was retained in the main document because
% it minimizes braidings: with parameter actions V . x = J(V) tensor x acting
% from the left, output-on-the-left composes through Theta_z without symmetries
% at Para-action sites (e.g. cor.parameterized_representation). The
% input-first variant requires roughly 3 to 5 extra braidings across the
% document.
%
% Kept as a reference; not maintained alongside the main document.

\documentclass[11pt, oneside, article]{memoir}


\settrims{0pt}{0pt} % page and stock same size
\settypeblocksize{*}{35.9pc}{*} % {height}{width}{ratio}
\setlrmargins{*}{*}{1} % {spine}{edge}{ratio}
\setulmarginsandblock{.98in}{.98in}{*} % height of typeblock computed
\setheadfoot{\onelineskip}{2\onelineskip} % {headheight}{footskip}
\setheaderspaces{*}{1.5\onelineskip}{*} % {headdrop}{headsep}{ratio}
\checkandfixthelayout


\usepackage{amsthm}
\usepackage{mathtools}

\usepackage[inline]{enumitem}
\usepackage{ifthen}
\usepackage[utf8]{inputenc} %allows non-ascii in bib file
\usepackage{xcolor}

\usepackage[backend=biber, backref=true, maxbibnames = 10, style = alphabetic]{biblatex}
\usepackage[bookmarks=true, colorlinks=true, linkcolor=blue!50!black,
citecolor=orange!50!black, urlcolor=orange!50!black, pdfencoding=unicode]{hyperref}
\usepackage[capitalize]{cleveref}

\usepackage{tikz}
\usepackage{pdflscape}

\usepackage{amssymb}
\usepackage{newpxtext}
\usepackage[varg,bigdelims]{newpxmath}
\usepackage{mathrsfs}
\usepackage{dutchcal}
\usepackage[scaled=0.95]{helvet}
\DeclareMathAlphabet{\mathsf}{T1}{phv}{m}{n}
\usepackage{fontawesome}
\usepackage{ebproof}
\usepackage{stmaryrd}

% cleveref %
  \newcommand{\creflastconjunction}{, and\nobreakspace} % serial comma
  \crefformat{enumi}{\card#2#1#3}
  \crefalias{chapter}{section}


% biblatex %
  \addbibresource{Library20260419.bib}

% hyperref %
  \hypersetup{final}

% enumitem %
  \setlist{nosep}
  \setlistdepth{6}



% tikz %

  \usetikzlibrary{
  	cd,
  	math,
  	decorations.markings,
		decorations.pathreplacing,
  	positioning,
  	arrows.meta,
  	shapes,
		shadows,
		shadings,
  	calc,
  	fit,
  	quotes,
  	intersections,
  	circuits.logic.US,
  }

  \tikzset{
biml/.tip={Glyph[glyph math command=triangleleft, glyph length=.95ex]},
bimr/.tip={Glyph[glyph math command=triangleright, glyph length=.95ex]},
}

\tikzset{
	tick/.style={postaction={
  	decorate,
    decoration={markings, mark=at position 0.5 with
    	{\draw[-] (0,.4ex) -- (0,-.4ex);}}}
  }
}
\tikzset{
	slash/.style={postaction={
  	decorate,
    decoration={markings, mark=at position 0.5 with
    	{\draw[-] (.3ex,.3ex) -- (-.3ex,-.3ex);}}}
  }
}

% oriented WD style (Schultz--Spivak wiring diagrams)
\tikzset{
	oriented WD/.style={
		every to/.style={out=0,in=180,draw},
    label/.style={
    	font=\everymath\expandafter{\the\everymath\scriptstyle},
      inner sep=0pt,
      node distance=2pt and -2pt},
    semithick,
    node distance=1 and 1,
    decoration={markings, mark=at position \stringdecpos with \stringdec},
    ar/.style={postaction={decorate}},
    execute at begin picture={\tikzset{
    	x=\bbx, y=\bby,
      every fit/.style={inner xsep=\bbx, inner ysep=\bby}}}
    },
    string decoration/.store in=\stringdec,
    string decoration={\arrow{stealth};},
    string decoration pos/.store in=\stringdecpos,
    string decoration pos=.7,
    bbx/.store in=\bbx,
    bbx = 1.5cm,
    bby/.store in=\bby,
    bby = 1.5ex,
    bb port sep/.store in=\bbportsep,
    bb port sep=1.5,
    bb port length/.store in=\bbportlen,
    bb port length=4pt,
    bb penetrate/.store in=\bbpenetrate,
    bb penetrate=0,
    bb min width/.store in=\bbminwidth,
    bb min width=1cm,
    bb rounded corners/.store in=\bbcorners,
    bb rounded corners=2pt,
    bb small/.style={
    	bb port sep=1, bb port length=2.5pt, bbx=.4cm, bb min width=.4cm, bby=.7ex},
    bb/.code 2 args={
    	\pgfmathsetlengthmacro{\bbheight}{\bbportsep * (max(#1,#2)+1) * \bby}
      \pgfkeysalso{draw,minimum height=\bbheight,minimum
       width=\bbminwidth,outer sep=0pt,
         rounded corners=\bbcorners,thick,
         prefix after command={\pgfextra{\let\fixname\tikzlastnode}},
         append after command={\pgfextra{\draw
            \ifnum #1=0{} \else foreach \i in {1,...,#1} {
            	($(\fixname.north west)!{\i/(#1+1)}!(\fixname.south west)$) +(-\bbportlen,0) coordinate (\fixname_in\i) -- +(\bbpenetrate,0) coordinate (\fixname_in\i')}\fi
            \ifnum #2=0{} \else foreach \i in {1,...,#2} {
            	($(\fixname.north east)!{\i/(#2+1)}!(\fixname.south east)$) +(-\bbpenetrate,0) coordinate (\fixname_out\i') -- +(\bbportlen,0) coordinate (\fixname_out\i)}\fi;
           }}}
		},
		bb name/.style={
     	append after command={
				\pgfextra{\node[anchor=north] at (\fixname.north) {#1};}
			}
		}
}

\newcommand{\upp}{\begin{tikzcd}[row sep=6pt]~\\~\ar[u, bend left=50pt, looseness=1.3, start anchor=east, end anchor=east]\end{tikzcd}}

\newcommand{\bito}[1][]{
	\begin{tikzcd}[ampersand replacement=\&, cramped]\ar[r, biml-bimr, "#1"]\&~\end{tikzcd}
}
\newcommand{\bifrom}[1][]{
	\begin{tikzcd}[ampersand replacement=\&, cramped]\ar[r, bimr-biml, "{#1}"]\&~\end{tikzcd}
}
\newcommand{\bifromlong}[2][]{
	\begin{tikzcd}[ampersand replacement=\&, column sep=#2, cramped]\ar[r, bimr-biml, "#1"]\&~\end{tikzcd}
}

% Adjunctions
\newcommand{\adj}[5][30pt]{%[size] Cat L, Left, Right, Cat R.
\begin{tikzcd}[ampersand replacement=\&, column sep=#1]
  #2\ar[r, shift left=6pt, "#3"]
  \ar[r, phantom, "\scriptstyle\Rightarrow"]\&
  #5\ar[l, shift left=6pt, "#4"]
\end{tikzcd}
}

\newcommand{\adjr}[5][30pt]{%[size] Cat R, Right, Left, Cat L.
\begin{tikzcd}[ampersand replacement=\&, column sep=#1]
  #2\ar[r, shift left=7pt, "#3"]\&
  #5\ar[l, shift left=7pt, "#4"]
  \ar[l, phantom, "\scriptstyle\Leftarrow"]
\end{tikzcd}
}

\newcommand{\xtickar}[1]{\begin{tikzcd}[baseline=-0.5ex,cramped,sep=small,ampersand
replacement=\&]{}\ar[r,tick, "{#1}"]\&{}\end{tikzcd}}
\newcommand{\xslashar}[1]{\begin{tikzcd}[baseline=-0.5ex,cramped,sep=small,ampersand
replacement=\&]{}\ar[r,slash, "{#1}"]\&{}\end{tikzcd}}



  % amsthm + aliascnt (so cleveref prints the correct environment name) %
\usepackage{aliascnt}

\theoremstyle{definition}
\newtheorem{definitionx}{Definition}[section]

\newaliascnt{examplex}{definitionx}
\newtheorem{examplex}[examplex]{Example}
\aliascntresetthe{examplex}
\crefname{examplex}{Example}{Examples}

\newaliascnt{remarkx}{definitionx}
\newtheorem{remarkx}[remarkx]{Remark}
\aliascntresetthe{remarkx}
\crefname{remarkx}{Remark}{Remarks}

\newaliascnt{notation}{definitionx}
\newtheorem{notation}[notation]{Notation}
\aliascntresetthe{notation}
\crefname{notation}{Notation}{Notations}

\newaliascnt{conventionx}{definitionx}
\newtheorem{conventionx}[conventionx]{Convention}
\aliascntresetthe{conventionx}
\crefname{conventionx}{Convention}{Conventions}

\theoremstyle{plain}

\newaliascnt{theorem}{definitionx}
\newtheorem{theorem}[theorem]{Theorem}
\aliascntresetthe{theorem}
\crefname{theorem}{Theorem}{Theorems}

\newaliascnt{proposition}{definitionx}
\newtheorem{proposition}[proposition]{Proposition}
\aliascntresetthe{proposition}
\crefname{proposition}{Proposition}{Propositions}

\newaliascnt{corollary}{definitionx}
\newtheorem{corollary}[corollary]{Corollary}
\aliascntresetthe{corollary}
\crefname{corollary}{Corollary}{Corollaries}

\newaliascnt{lemma}{definitionx}
\newtheorem{lemma}[lemma]{Lemma}
\aliascntresetthe{lemma}
\crefname{lemma}{Lemma}{Lemmas}

\newaliascnt{warning}{definitionx}
\newtheorem{warning}[warning]{Warning}
\aliascntresetthe{warning}
\crefname{warning}{Warning}{Warnings}

\newaliascnt{fact}{definitionx}
\newtheorem{fact}[fact]{Fact}
\aliascntresetthe{fact}
\crefname{fact}{Fact}{Facts}

\newtheorem*{theorem*}{Theorem}
\newtheorem*{proposition*}{Proposition}
\newtheorem*{corollary*}{Corollary}
\newtheorem*{lemma*}{Lemma}
\newtheorem*{warning*}{Warning}

\newenvironment{example}
  {\pushQED{\qed}\renewcommand{\qedsymbol}{$\lozenge$}\examplex}
  {\popQED\endexamplex}

 \newenvironment{remark}
  {\pushQED{\qed}\renewcommand{\qedsymbol}{$\lozenge$}\remarkx}
  {\popQED\endremarkx}

  \newenvironment{definition}
  {\pushQED{\qed}\renewcommand{\qedsymbol}{$\lozenge$}\definitionx}
  {\popQED\enddefinitionx}

  \newenvironment{convention}
  {\pushQED{\qed}\renewcommand{\qedsymbol}{$\lozenge$}\conventionx}
  {\popQED\endconventionx}


%-------- Single symbols --------%

\DeclareSymbolFont{stmry}{U}{stmry}{m}{n}
\DeclareMathSymbol\fatsemi\mathop{stmry}{"23}

\DeclareFontFamily{U}{mathx}{\hyphenchar\font45}
\DeclareFontShape{U}{mathx}{m}{n}{
      <5> <6> <7> <8> <9> <10>
      <10.95> <12> <14.4> <17.28> <20.74> <24.88>
      mathx10
      }{}
\DeclareSymbolFont{mathx}{U}{mathx}{m}{n}
\DeclareFontSubstitution{U}{mathx}{m}{n}
\DeclareMathAccent{\widecheck}{0}{mathx}{"71}


%-------- Renewed commands --------%

\renewcommand{\ss}{\subseteq}

%-------- Other Macros --------%


\DeclarePairedDelimiter{\present}{\langle}{\rangle}
\DeclarePairedDelimiter{\pairing}{\langle}{\rangle}
\DeclarePairedDelimiter{\copair}{[}{]}
\DeclarePairedDelimiter{\floor}{\lfloor}{\rfloor}
\DeclarePairedDelimiter{\ceil}{\lceil}{\rceil}
\DeclarePairedDelimiter{\corners}{\ulcorner}{\urcorner}
\DeclarePairedDelimiter{\ihom}{[}{]}

\DeclareMathOperator{\Hom}{Hom}
\DeclareMathOperator{\Aut}{Aut}
\DeclareMathOperator{\End}{End}
\DeclareMathOperator{\Mor}{Mor}
\DeclareMathOperator{\dom}{dom}
\newcommand{\Int}{\mathop{\Fun{Int}}\nolimits}
\newcommand{\Lie}{\mathop{\Fun{Lie}}\nolimits}
\DeclareMathOperator{\cod}{cod}
\DeclareMathOperator{\idy}{idy}
\DeclareMathOperator*{\colim}{colim}
\DeclareMathOperator{\im}{im}
\DeclareMathOperator{\ob}{Ob}
\DeclareMathOperator{\Tr}{Tr}
\DeclareMathOperator{\el}{El}




\newcommand{\const}[1]{\texttt{#1}}%a constant, or named element of a set
\newcommand{\Set}[1]{\mathsf{#1}}%a named set
\newcommand{\ord}[1]{\mathsf{#1}}%an ordinal
\newcommand{\cat}[1]{\mathcal{#1}}%a generic category
\newcommand{\Cat}[1]{\mathbf{#1}}%a named category
\newcommand{\fun}[1]{\mathrm{#1}}%a function
\newcommand{\Fun}[1]{\mathsf{#1}}%a named functor




\newcommand{\id}{\mathrm{id}}
\newcommand{\then}{\mathbin{\fatsemi}}

\newcommand{\cocolon}{:\!}


\newcommand{\iso}{\cong}
\newcommand{\too}{\longrightarrow}
\newcommand{\tto}{\rightrightarrows}
\newcommand{\To}[2][]{\xrightarrow[#1]{#2}}
\renewcommand{\Mapsto}[1]{\xmapsto{#1}}
\newcommand{\Tto}[3][13pt]{\begin{tikzcd}[sep=#1, cramped, ampersand replacement=\&, text height=1ex, text depth=.3ex]\ar[r, shift left=2pt, "#2"]\ar[r, shift right=2pt, "#3"']\&{}\end{tikzcd}}
\newcommand{\Too}[1]{\xrightarrow{\;\;#1\;\;}}
\newcommand{\from}{\leftarrow}
\newcommand{\fromm}{\longleftarrow}
\newcommand{\ffrom}{\leftleftarrows}
\newcommand{\From}[1]{\xleftarrow{#1}}
\newcommand{\Fromm}[1]{\xleftarrow{\;\;#1\;\;}}
\newcommand{\surj}{\twoheadrightarrow}
\newcommand{\inj}{\rightarrowtail}
\newcommand{\wavyto}{\rightsquigarrow}
\newcommand{\lollipop}{\multimap}
\newcommand{\imp}{\Rightarrow}
\renewcommand{\iff}{\Leftrightarrow}
\newcommand{\down}{\mathbin{\downarrow}}
\newcommand{\fromto}{\leftrightarrows}
\newcommand{\tickar}{\xtickar{}}
\newcommand{\slashar}{\xslashar{}}
\newcommand{\card}{\,^{\#}}
\newcommand{\bk}[1]{\overset{\scriptscriptstyle\leftarrow}{#1}}


\newcommand{\inv}{^{-1}}
\newcommand{\op}{^\tn{op}}

\newcommand{\tn}[1]{\textnormal{#1}}
\newcommand{\ol}[1]{\overline{#1}}
\newcommand{\ul}[1]{\underline{#1}}
\newcommand{\wt}[1]{\widetilde{#1}}
\newcommand{\wh}[1]{\widehat{#1}}
\newcommand{\wc}[1]{\widecheck{#1}}
\newcommand{\ubar}[1]{\underaccent{\bar}{#1}}



\newcommand{\bb}{\mathbb{B}}
\newcommand{\cc}{\mathbb{C}}
\newcommand{\nn}{\mathbb{N}}
\newcommand{\pp}{\mathbb{P}}
\newcommand{\qq}{\mathbb{Q}}
\newcommand{\zz}{\mathbb{Z}}
\newcommand{\rr}{\mathbb{R}}


\newcommand{\finset}{\Cat{Fin}}
\newcommand{\smset}{\Cat{Set}}
\newcommand{\smcat}{\Cat{Cat}}
\newcommand{\catsharp}{\Cat{Cat}^{\sharp}}
\newcommand{\en}{\Cat{End}}
\newcommand{\vect}{\Cat{Vect}}

\newcommand{\List}{\Fun{List}}
\newcommand{\set}{\tn{-}\Cat{Set}}

\newcommand{\act}{\tn{act}}
\newcommand{\upd}{\tn{upd}}


\newcommand{\yon}{\mathcal{y}}
\newcommand{\poly}{\Cat{Poly}}
\newcommand{\polycart}{\poly^{\tn{cart}}}
\newcommand{\tri}{\mathbin{\triangleleft}}

% lenses
\newcommand{\biglens}[2]{
     \begin{bmatrix}{\vphantom{f_f^f}#2} \\ {\vphantom{f_f^f}#1} \end{bmatrix}
}
\newcommand{\littlelens}[2]{
     \begin{bsmallmatrix}{\vphantom{f}#2} \\ {\vphantom{f}#1} \end{bsmallmatrix}
}
\newcommand{\lens}[2]{
  \relax\if@display
     \biglens{#1}{#2}
  \else
     \littlelens{#1}{#2}
  \fi
}


\newcommand{\slogan}[1]{\begin{center}\textit{#1}\end{center}}

\newcommand{\qand}{\quad\text{and}\quad}
\newcommand{\qqand}{\qquad\text{and}\qquad}
\newcommand{\where}{\qquad\text{where}\quad}

\newcommand{\cofun}{{\raisebox{2pt}{\resizebox{2.5pt}{2.5pt}{$\setminus$}}}}

\newcommand{\coalg}{\tn{-}\Cat{Coalg}}

\newcommand{\org}{{\mathbb{O}\Cat{rg}}}

\newcommand{\idcoalg}[1]{{\{\id_{#1}\}}}

\newcommand{\dnote}[1]{{\quad \color{blue}$\lozenge$\;David says:}~#1\;{\color{blue}$\lozenge$}\quad}

%-------- Paper-specific macros --------%

\newcommand{\mfd}{\Cat{Mfd}}
\newcommand{\plenspot}{\Cat{PLensPot}}
\newcommand{\boxx}{\Fun{Box}}
\RenewDocumentCommand{\cot}{g}{\IfValueTF{#1}{\fun{cot}(#1)}{\fun{cot}}}
\newcommand{\ev}{\fun{ev}}
\newcommand{\coev}{\fun{coev}}
\newcommand{\pr}{\fun{pr}}
\newcommand{\Part}{\bullet}
\newcommand{\boxob}{\mathrm{box}}
\newcommand{\Ad}{\fun{Ad}}
\newcommand{\ad}{\fun{ad}}
\newcommand{\lie}[1]{\mathfrak{#1}}
\newcommand{\tpow}[1]{^{\otimes #1}}
\newcommand{\potlens}{\Cat{PLens}}
\newcommand{\Para}{\Cat{Para}}
\newcommand{\para}[2]{\Para_{#1}(#2)}
\newcommand{\pnla}{\Cat{PNLA}}
\newcommand{\pvect}{\Cat{PVect}}
\newcommand{\Lens}[1]{\Cat{Lens}_{#1}}
\newcommand{\lmfd}{\Lens{\mfd}}
\newcommand{\plm}{\Cat{PLM}}
\newcommand{\inpt}[1]{#1^{-}}
\newcommand{\outp}[1]{#1^{+}}
\newcommand{\lensob}[1]{\binom{\inpt{#1}}{\outp{#1}}}% vertical lens-object notation; to revert, change to (\inpt{#1},\outp{#1})
\newcommand{\lensobs}[2]{\binom{\inpt{#1}\times\inpt{#2}}{\outp{#1}\times\outp{#2}}}% monoidal product of two lens objects
\newcommand{\potd}{\fun{d}}
\newcommand{\cint}{\fun{cint}}
\newcommand{\leg}{\fun{leg}}
\newcommand{\dyn}{\fun{dyn}}

\font\MnSyCfont=MnSymbolC10
\newcommand{\smallsquare}{\mbox{\MnSyCfont\symbol{"68}}}
\newcommand{\smallbullet}{\mbox{\MnSyCfont\symbol{"59}}}
\newcommand{\blank}{\mathord{\smallsquare}}

% ---- Changeable document parameters ---- %

\linespread{1.15}
%\allowdisplaybreaks
\setsecnumdepth{subsection}
\settocdepth{subsection}
\setlength{\parindent}{15pt}
\setcounter{tocdepth}{1}



%--------------- Document ---------------%
\begin{document}

\title{Nested Potentials}

\author{David I. Spivak}

\date{\vspace{-.2in}}

\maketitle

\begin{abstract}

\end{abstract}

\tableofcontents

%--------------------------------------------------------------------
\chapter{Introduction}\label{ch.intro}
%--------------------------------------------------------------------

% TODO: write introduction
\emph{[Introduction to be written.]}


%--------------------------------------------------------------------
\chapter{Background}\label{ch.background}
%--------------------------------------------------------------------

\section{Notation and preliminaries}\label{sec.prelim}

We fix some typographic conventions. A generic category is written in calligraphic script, $\cat{C},\cat{M},\cat{A}$, and a specifically-named category in boldface, $\Cat{Set},\Cat{Poly},\Cat{Mfd}$. A function or operation is written in roman, $\fun{ev},\fun{pr}$, and a named functor in sans-serif, $\Fun{T},\Fun{Box}$. Lie algebras use fraktur, $\lie{g},\lie{p}$; number systems use blackboard bold, $\nn,\rr$, as does the only bicategory we work with, $\org$. 

In a symmetric monoidal category $(\cat{C},I,\otimes)$, a \emph{$(\otimes)$-comonoid} is an object $c$ equipped with a coassociative comultiplication $\delta_c\colon c\to c\otimes c$ and a counit $\varepsilon_c\colon c\to I$ satisfying the counit law; a \emph{$(\otimes)$-monoid} is the dual. A \emph{comonoid homomorphism} $f\colon c\to c'$ satisfies $f\then\varepsilon_{c'}=\varepsilon_c$ and $f\then\delta_{c'}=\delta_c\then(f\otimes f)$. 

A \emph{supply of comonoids}~\cite{fong2019supplying} is a choice, compatible with $\otimes$, of a cocommutative such structure on every object of $\cat{C}$. A supply is \emph{homomorphic} if every morphism of $\cat C$ is a comonoid homomorphism for the supplied structures.

\begin{lemma}\label{lem.cartesian_homomorphic_supply}
Any cartesian monoidal category canonically has a homomorphic supply of commutative comonoids, given by diagonals and terminal maps.
\end{lemma}

\begin{proof}
This is the standard cartesian comonoid structure of \cite{fox1976coalgebras}.
\end{proof}

\section{Manifolds and paired vector spaces}\label{sec.lie_pnla_bg}

\subsection{Manifolds and notation}\label{sec.manifolds_notation}

We work in the cartesian category $(\mfd,\rr^0,\times)$ of finite-dimensional smooth real manifolds and smooth maps. For a manifold $M$ and point $m\colon\rr^0\to M$, write $T_mM$ for the tangent space at $m$ and $TM$ for the tangent bundle; $T^*_mM$ and $T^*M$ for the cotangent space and bundle. The tangent functor preserves products, so $T(M\times N)\cong TM\times TN$; in particular $T_{(m,n)}(M\times N)\cong T_mM\oplus T_nN$ and dually $T^*_{(m,n)}(M\times N)\cong T^*_mM\oplus T^*_nN$.

For a smooth map $f\colon M\to N$, write
\[
Tf\colon TM\to TN
\]
for the \emph{tangent map} (also called the \emph{global differential}), and
\[
T_m f\colon T_mM\to T_{f(m)}N
\qqand
(T_m f)^\top\colon T^*_{f(m)}N\to T^*_mM
\]
for the \emph{differential at $m$}---the restriction of $Tf$ to the fiber over $m\in M$--- and its transpose.%
\footnote{What we call $T_m f$ is often denoted $Df_m$.}

When the codomain is $\rr$, the tangent map of $U\colon M\to\rr$ can be combined with the standard constant covector field $r\mapsto +1$ on $\rr$ to give a covector field on $M$, also called the \emph{differential} of $U$:
\begin{equation}\label{eqn.differential}
dU\colon M\to T^*M,\qquad dU|_m\coloneqq(T_m U)^\top(+1)\in T^*_mM.
\end{equation}
That is, $dU|_m$ is the covector at $m$ whose pairing with $\dot m\in T_mM$ is $+1(T_m U(\dot m))$, identified with a real number using $T_{U(m)}\rr\cong\rr$. We generally write $\xi$ for a covector, and use dotted point notation such as $\dot m$ for a tangent vector based at $m$.

We often elide the inclusion $\vect\to\mfd$, treating finite-dimensional real vector spaces as manifolds. The direct sum $(0,\oplus)$ structure in $\vect$ is (a coproduct and) a product, and the inclusion is product preserving. By standard manifold theory, the canonical global chart on a finite-dimensional real vector space $V$ provides linear isomorphisms
\begin{equation}\label{eqn.tangent_vec}
T_vV \;\cong\; V \qqand T^*_vV \;\cong\; V^*
\end{equation}
at every point $v\in V$; we use these identifications throughout, often without comment. In particular, the differential at $v$ of a smooth function $f\colon V\to\rr$ is an element $df|_v\in V^*$, also called \emph{the differential of $f$ at $v$}.
We also use the canonical finite-dimensional dual-sum identification
\begin{equation}\label{eqn.canonical_dual_sum}
(V\oplus V^*)^*\cong V^*\oplus V,\qquad
\lambda\leftrightarrow(\xi,v)\quad\text{where}\quad
\lambda(v',\xi')=\xi(v')+\xi'(v).
\end{equation}

\subsection{Paired vector spaces}\label{sec.pnla}

\begin{definition}\label{def.pnla}
A \emph{paired vector space} is a finite-dimensional real vector space $V$ together with a linear isomorphism $\sharp_V\colon V^*\To{\cong}V$, the \emph{sharp} map, with inverse $\flat_V\coloneqq\sharp_V\inv\colon V\to V^*$, the \emph{flat} map. The corresponding nondegenerate bilinear form on $V$, the \emph{pairing}, is
\begin{equation}\label{eqn.pairing}
\pairing{v,v'}_V\coloneqq\flat_V(v)(v')\colon V\otimes V\to\rr
\end{equation}
The three forms $\sharp_V$, $\flat_V$, and $\pairing{-,-}_V$ package the same data; we use whichever is most convenient.

We define the category $\pvect$ to have paired vector spaces as objects, and as morphisms $(V,\sharp_V)\to(W,\sharp_W)$ the linear isomorphisms $\phi\colon V\To{\cong}W$ preserving the pairing: $\pairing{\phi v,\phi v'}_W=\pairing{v,v'}_V$ for all $v,v'\in V$.%
\footnote{
Restricting to isomorphisms (making $\pvect$ a groupoid) is essential for our development below: each $\pvect$-morphism is used both as $\phi$ on vectors and as $(\phi^*)\inv$ on covectors, e.g.\ in $T^*\phi$ \eqref{eqn.TT_phi} below.
}
Equivalently, $\phi$ preserves the pairing iff the \emph{pairing-preservation equation} holds:
\begin{equation}\label{eqn.pairing_triangle}
\phi\circ\sharp_V\circ\phi^*=\sharp_W\colon W^*\to W
\qedhere
\end{equation}
\end{definition}

\begin{proposition}\label{prop.pnla_monoidal}
Direct sum gives $\pvect$ a symmetric monoidal structure $(0,\oplus)$, with the zero space (vacuous pairing) as unit and direct-sum pairing $\sharp_{V\oplus W}\coloneqq\sharp_V\oplus\sharp_W$ on $V\oplus W$, equivalently the bilinear form
\begin{equation}\label{eqn.directsum_pairing}
\pairing{(v,w),(v',w')}_{V\oplus W} = \pairing{v,v'}_V+\pairing{w,w'}_W.
\end{equation}
\end{proposition}

\begin{proof}
A direct sum of linear isomorphisms is a linear isomorphism, and \eqref{eqn.directsum_pairing} is direct from $\flat_{V\oplus W}=\flat_V\oplus\flat_W$. Coherence is inherited from $(\vect,0,\oplus)$.\qedhere
\end{proof}

\begin{proposition}\label{prop.pnla_smc_action}
The composite $\pvect\to\vect\hookrightarrow\mfd$ is strong symmetric monoidal, so $\pvect$ acts on $\mfd$ via $((V,\sharp_V),M)\mapsto V\times M$.
\end{proposition}

\begin{proof}
The forgetful $\pvect\to\vect$ preserves direct sums and identities; the inclusion $\vect\hookrightarrow\mfd$ preserves products.\qedhere
\end{proof}

\begin{proposition}\label{prop.canonical_symplectic_pairing}
For any finite-dimensional vector space $V$, the direct sum $V\oplus V^*$ is a paired vector space under the \emph{canonical symplectic pairing}
\begin{equation}\label{eqn.canonical_form}
\pairing{(v,\xi),(v',\xi')}_{V\oplus V^*} \;=\; \xi'(v)-\xi(v').
\end{equation}
\end{proposition}

\begin{proof}
The form in \eqref{eqn.canonical_form} is bilinear. To see that it is nondegenerate, suppose $(v,\xi)$ pairs to $0$ with every $(v',\xi')$. Taking $v'=0$ gives $\xi'(v)=0$ for all $\xi'\in V^*$, hence $v=0$; taking $\xi'=0$ gives $\xi(v')=0$ for all $v'\in V$, hence $\xi=0$.\qedhere
\end{proof}

For the canonical symplectic pairing, under \eqref{eqn.canonical_dual_sum}, the associated sharp map is
\begin{equation}\label{eqn.canonical_sharp}
\sharp_{V\oplus V^*}\colon V^*\oplus V\to V\oplus V^*,\qquad(\xi,v)\mapsto(v,-\xi).
\end{equation}

\begin{remark}[Sign convention]\label{rmk.swap_pairing}
The bilinear form
\[
\pairing{(v,\xi),(v',\xi')}=\xi'(v)+\xi(v')
\]
also makes $V\oplus V^*$ a paired vector space; under \eqref{eqn.canonical_dual_sum}, its sharp map is the swap $(\xi,v)\mapsto(v,\xi)$ whereas \eqref{eqn.canonical_form} has sharp map \eqref{eqn.canonical_sharp}. We use \eqref{eqn.canonical_form} because its antisymmetry is required for Hamiltonian dynamics (\cref{prop.hamilton,rmk.antisymmetry_substantive}).\qedhere
\end{remark}

\subsection{Tangent and cotangent endofunctors}\label{sec.TT}

For a paired vector space $(V,\sharp_V):\pvect$, define
\begin{equation}\label{eqn.TT_def}
TV\coloneqq V\oplus V \qqand T^*V\coloneqq V\oplus V^*.
\end{equation}
Under the inclusion functor $\vect\to\mfd$ and the global-chart identifications \eqref{eqn.tangent_vec}, these definitions agree with the tangent and cotangent bundles of $V$ as a manifold.

Each is naturally a paired vector space: $TV$ via $\sharp_V\oplus\sharp_V$ (\cref{prop.pnla_monoidal}), $T^*V$ via the canonical symplectic pairing, whose sharp map is \eqref{eqn.canonical_sharp}.
For a $\pvect$-morphism $\phi\colon V\to W$, define
\begin{equation}\label{eqn.TT_phi}
T\phi\coloneqq\phi\oplus\phi
\qqand
 T^*\phi\coloneqq\phi\oplus(\phi^*)\inv.
\end{equation}

\begin{proposition}\label{prop.TT_endofunctors}
The assignments \eqref{eqn.TT_def}, \eqref{eqn.TT_phi} define endofunctors $T,T^*\colon\pvect\to\pvect$.
\end{proposition}

\begin{proof}
If $\phi\colon V\to W$ is a linear isomorphism so are $T\phi$ and $T^*\phi$. Functoriality is obvious. We first check that $T\phi$ preserves pairings.
\begin{align*}
\pairing{T\phi(v_1,v_1'),T\phi(v_2,v_2')}_{TW}
&= \pairing{(\phi v_1,\phi v_1'),(\phi v_2,\phi v_2')}_{TW} && \text{[\eqref{eqn.TT_phi}]}\\
&= \pairing{\phi v_1,\phi v_2}_W+\pairing{\phi v_1',\phi v_2'}_W && \text{[\eqref{eqn.directsum_pairing}]}\\
&= \pairing{v_1,v_2}_V+\pairing{v_1',v_2'}_V && \text{[$\phi$ preserves pairings]}\\
&= \pairing{(v_1,v_1'),(v_2,v_2')}_{TV}.
\end{align*}

Next we check that $T^*\phi$ preserves pairings. Writing $\psi\coloneqq(\phi^*)\inv\colon V^*\to W^*$, the defining property of $\phi^*$ states that $(\phi^*\xi')(v) = \xi'(\phi v)$ for all $\xi'\in W^*\text{ and }v\in V.$ Setting $\xi'\coloneqq\psi\xi$ and using $\phi^*\circ\psi=\id_{V^*}$, we note that
\begin{equation}\label{eqn.dualeval}
(\psi\xi)(\phi v) = (\phi^*(\psi\xi))(v) = \xi(v),
\end{equation}
and we can conclude that $T^*\phi$ preserves pairings by calculating
\begin{align*}
\pairing{T^*\phi(v_1,\xi_1),T^*\phi(v_2,\xi_2)}_{T^*W}
&= (\psi\xi_2)(\phi v_1)-(\psi\xi_1)(\phi v_2) && \text{[\eqref{eqn.TT_phi}, \eqref{eqn.canonical_form}]}\\
&= \xi_2(v_1)-\xi_1(v_2) && \text{[\eqref{eqn.dualeval}]}\\
&= \pairing{(v_1,\xi_1),(v_2,\xi_2)}_{T^*V} && \text{[\eqref{eqn.canonical_form}]}.
\end{align*}
\end{proof}

\begin{proposition}\label{prop.TT_monoidal}
The endofunctors $T,T^*\colon\pvect\to\pvect$ are strong symmetric monoidal with respect to $(0,\oplus)$.
\end{proposition}

\begin{proof}
The unit comparisons are identities since $T0=0\oplus 0=0$ and $T^*0=0\oplus 0^*=0$. The product comparisons are given by interchange:
\begin{align}
\sigma\colon T(V_1\oplus V_2)&\To{\cong}TV_1\oplus TV_2,\quad ((v_1,v_2),(v_1',v_2'))\mapsto((v_1,v_1'),(v_2,v_2')),\label{eqn.T_swap}\\
\sigma^*\colon T^*(V_1\oplus V_2)&\To{\cong}T^*V_1\oplus T^*V_2,\quad ((v_1,v_2),(\xi_1,\xi_2))\mapsto((v_1,\xi_1),(v_2,\xi_2))\label{eqn.Tstar_swap}
\end{align}
(using $(V_1\oplus V_2)^*\cong V_1^*\oplus V_2^*$ for $\sigma^*$). Both are linear isomorphisms. To check preservation of pairings, compare with a second tangent element $((w_1,w_2),(w_1',w_2'))$ and a second cotangent element $((w_1,w_2),(\xi'_1,\xi'_2))$. The two sides for $\sigma$ expand to
\[
\pairing{v_1,w_1}_{V_1}+\pairing{v_1',w_1'}_{V_1}+\pairing{v_2,w_2}_{V_2}+\pairing{v_2',w_2'}_{V_2},
\]
and the two sides for $\sigma^*$ expand via \eqref{eqn.canonical_form} to
\[
\xi'_1(v_1)+\xi'_2(v_2)-\xi_1(w_1)-\xi_2(w_2).
\]
\qedhere
\end{proof}

\begin{proposition}\label{prop.legendre}
There is a monoidal natural isomorphism $\sharp\colon T^*\;\xRightarrow{\cong}\;T$ whose components are given by
\[
\id_V\oplus\sharp_V\colon T^*V\to TV,\qquad (v,\xi)\mapsto(v,\sharp_V(\xi)).
\]
\end{proposition}

\begin{proof}
Each component is a linear isomorphism by \cref{def.pnla}; naturality is the pairing-preservation axiom of \cref{def.pnla}; monoidality follows from the direct-sum pairing convention $\sharp_{V_1\oplus V_2}=\sharp_{V_1}\oplus\sharp_{V_2}$ (\cref{prop.pnla_monoidal}).\qedhere
\end{proof}

\begin{remark}[Legendre transform]\label{rmk.legendre}
\Cref{prop.legendre} is the (infinitesimal) \emph{Legendre transform} of classical mechanics: $\flat_V$ sends a velocity to its conjugate momentum, and $\sharp_V$ sends a momentum to the corresponding velocity.\qedhere
\end{remark}

The \emph{exponential} of a paired vector space $(V,\sharp_V)$ is the time-1 flow of the constant left-invariant vector field $\sharp_V(\xi)$ from base point $v$:
\begin{equation}\label{eqn.flow_cot}
\exp_V\colon T^*V\to V,\qquad (v,\xi)\longmapsto v+\sharp_V(\xi).
\end{equation}
It is natural in $\pvect$-isomorphisms by linearity and the pairing-preservation equation \eqref{eqn.pairing_triangle}. We use \eqref{eqn.flow_cot} pointwise as the on-directions bijection of $V\yon^V\cong\cot{V}$ in \cref{prop.pnla_polynomial}.

\begin{remark}[Generalization to nilpotent Lie algebras]\label{rmk.pnla_generalization}
\Cref{sec.pnla,sec.TT} generalize from paired vector spaces (abelian $\lie g$) to \emph{paired nilpotent Lie algebras} (PN Lie Algebras): pairs $(\lie g,\pairing{-,-})$ of a finite-dimensional nilpotent real Lie algebra and a pairing on its underlying vector space. Replacing $\pvect$ by $\pnla$ throughout, \cref{prop.pnla_monoidal,prop.canonical_symplectic_pairing,prop.TT_monoidal,prop.legendre} hold verbatim; the rest pick up Lie-group ingredients:
\begin{itemize}
\item \Cref{prop.pnla_smc_action}: the action becomes $((\lie g,\sharp),M)\mapsto\Int\lie g\times M$, where $\Int\lie g$ is the simply-connected Lie group integrating $\lie g$ (Lie's third theorem~\cite[App.~B]{knapp2002lie}; $\Int V\cong V$ for vector spaces).
\item \Cref{prop.TT_endofunctors}: $TV=V\oplus V$ and $T^*V=V\oplus V^*$ in \eqref{eqn.TT_def} are unchanged as vector spaces but now carry brackets mixing $\ad$ and $\ad^*$ (nilpotent of class $\le 2n-1$ when $\lie g$ has class $n$). Morphism formulas \eqref{eqn.TT_phi} are unchanged; bracket-preservation needs $(\phi^*)\inv\circ\ad^*_X=\ad^*_{\phi X}\circ(\phi^*)\inv$, a transpose of $\phi$'s Lie-algebra-homomorphism property.
\item Exponential \eqref{eqn.flow_cot}: extends to $T^*G\to G$, $(g,\xi)\mapsto g\cdot\exp_G(\sharp_G\xi)$ for $G\coloneqq\Int\lie g$, via Malcev's diffeomorphism $\exp_G\colon\lie g\To{\cong}\Int\lie g$~\cite[Thm.~1.127]{knapp2002lie}. The vector-space case is abelian: $\Int\lie g=\lie g$, $\exp_G=\id$, multiplication is addition.\qedhere
\end{itemize}
\end{remark}

\subsection{Hamiltonian vector fields}\label{sec.hamiltonian_vector_fields}

\cref{sec.hamiltonian_vector_fields} is not necessary and can be skimmed on a first reading; it will only be used in discussing the wave equation \cref{sec.spring}. We include it because it allows physicists and category theorists to exchange basic intuitions using the basic concepts and results we've developed.

\begin{definition}[Hamiltonian vector field]\label{def.hamiltonian_vector_field}
Let $V$ be a paired vector space and $H\colon T^*V\to\rr$ a smooth function (the \emph{Hamiltonian}). The \emph{Hamiltonian vector field of $H$} is the composite $X_H$ obtained as follows:
\[
\begin{tikzcd}[column sep=1.8em, row sep=2.4em, ampersand replacement=\&]
T^*V
  \ar[d, dashed, "{X_H}"']
  \ar[r, "{(\id,dH)}"]
\&[15pt] T^*V\oplus(T^*V)^*
  \ar[r, equal, "{\eqref{eqn.TT_def},\,\eqref{eqn.canonical_dual_sum}}"]
\&[15pt] (V\oplus V^*)\oplus(V^*\oplus V)
  \ar[d, "{\id\oplus\sharp_{V\oplus V^*}}"]
\\
T(T^*V)
\& {}
\& (V\oplus V^*)\oplus(V\oplus V^*)
  \ar[ll, equal, "{\eqref{eqn.TT_def}^{-1}}"]
\end{tikzcd}
\]
The first arrow sends $(v,\xi)$ to $((v,\xi),dH|_{(v,\xi)})$. The arrow labeled $\id\oplus\sharp_{V\oplus V^*}$ is the component at $T^*V$ of the Legendre transform $\sharp\colon T^*\Rightarrow T$ (\cref{prop.legendre}); explicitly, it leaves the basepoint fixed and applies the canonical symplectic sharp map $\sharp_{V\oplus V^*}=\sharp_{T^*V}$ of \eqref{eqn.canonical_sharp} to the cotangent-fiber component:
\[
((v,\xi),(\xi',w))\longmapsto((v,\xi),(w,-\xi')).
\]
The bottom $\eqref{eqn.TT_def}^{-1}$ arrow repackages the explicit direct sum as $T(T^*V)$. Under \eqref{eqn.canonical_dual_sum}, we write
\begin{equation}\label{eqn.hamilton_differential}
dH|_{(v,\xi)}=(\partial_v H,\partial_\xi H)\in V^*\oplus V,
\qquad
dH|_{(v,\xi)}(v',\xi')=(\partial_v H)(v')+\xi'(\partial_\xi H).
\end{equation}
Thus the composite is
\[
X_H(v,\xi)=((v,\xi),(\partial_\xi H,-\partial_v H))\in T(T^*V).\qedhere
\]
\end{definition}

\begin{proposition}[Hamilton's equations]\label{prop.hamilton}
The integral curves $(v(t),\xi(t))$ of the Hamiltonian vector field of \cref{def.hamiltonian_vector_field} — those whose velocity is the fiber component of $X_H(v,\xi)$ — satisfy \emph{Hamilton's equations}
\[
\dot v=\partial_\xi H,\qquad \dot\xi=-\partial_v H,
\]
and $H$ is conserved along them: $dH/dt=0$.
\end{proposition}

\begin{proof}
By \cref{def.hamiltonian_vector_field}, the fiber component of $X_H(v,\xi)$ is $(\partial_\xi H,-\partial_v H)$. Thus the integral-curve condition is
\begin{equation}\label{eqn.hamilton_integral_curve}
(\dot v,\dot\xi)=(\partial_\xi H,-\partial_v H),
\end{equation}
which is Hamilton's equations.

For conservation, the chain rule together with \eqref{eqn.hamilton_differential} and \eqref{eqn.hamilton_integral_curve} gives
\[
\begin{aligned}
\frac{dH}{dt}
&= dH|_{(v,\xi)}(\dot v,\dot\xi) \\
&= (\partial_v H)(\dot v)+\dot\xi(\partial_\xi H) \\
&= (\partial_v H)(\partial_\xi H)+(-\partial_v H)(\partial_\xi H)=0.
\end{aligned}
\]
\end{proof}

\begin{remark}[Antisymmetry is substantive]\label{rmk.antisymmetry_substantive}
Had we used the swap pairing of \cref{rmk.swap_pairing} in place of \eqref{eqn.canonical_form}, the vertical map $\id\oplus\sharp_{V\oplus V^*}$ in \cref{def.hamiltonian_vector_field} would use the swap sharp $(\xi,v)\mapsto(v,\xi)$ instead of \eqref{eqn.canonical_sharp}. This would give $\dot\xi=+\partial_v H$ and $dH/dt=2(\partial_v H)(\partial_\xi H)$, which is generically nonzero: it represents a gradient flow rather than a Hamiltonian flow. Conservation is precisely what the antisymmetry of \eqref{eqn.canonical_form} buys.\qedhere
\end{remark}

\begin{remark}[Coordinate invariance of $X_H$]\label{rmk.XH_natural}
A $\pvect$-isomorphism $\phi\colon V\To{\cong}W$ induces a linear symplectomorphism $T^*\phi\colon T^*V\To{\cong}T^*W$ (\cref{prop.TT_endofunctors}). For any $H\colon T^*W\to\rr$, the chain rule for $H\circ T^*\phi$ together with naturality of the Legendre transform $\sharp$ (\cref{prop.legendre}) at the morphism $T^*\phi$ give
\[
X_H\circ T^*\phi=T(T^*\phi)\circ X_{H\circ T^*\phi},
\]
so $T^*\phi$ takes integral curves of $X_{H\circ T^*\phi}$ on $T^*V$ to those of $X_H$ on $T^*W$. Hamilton's equations transform covariantly under linear canonical transformations.\qedhere
\end{remark}

\begin{remark}[Decoupling of separable Hamiltonians]\label{rmk.XH_separable}
For Hamiltonians $H_i\colon T^*V_i\to\rr$, the separable Hamiltonian on $T^*(V_1\oplus V_2)$ defined in coordinates by
\[
H((v_1,v_2),(\xi_1,\xi_2))\coloneqq H_1(v_1,\xi_1)+H_2(v_2,\xi_2)
\]
satisfies $\partial_{v_i}H=\partial_{v_i}H_i$ and $\partial_{\xi_i}H=\partial_{\xi_i}H_i$, so by \cref{prop.hamilton} the dynamics decouple:
\[
\dot v_i=\partial_{\xi_i}H_i,\qquad \dot\xi_i=-\partial_{v_i}H_i\qquad(i=1,2).
\]
At the level of vector fields, the swaps \eqref{eqn.T_swap}, \eqref{eqn.Tstar_swap} of \cref{prop.TT_monoidal} together with the monoidality of $\sharp$ (\cref{prop.legendre}) factor $X_H$ as $X_{H_1}\oplus X_{H_2}$; non-interacting systems evolve independently on their own factors.\qedhere
\end{remark}

\section{Polynomial functors, comonads, and coalgebras}\label{sec.poly_dynamic_bg}

\subsection{The symmetric monoidal closed category $\poly$}\label{sec.poly}

We briefly review the category $\poly$ of polynomial functors; see~\cite{niu2025polynomial} for a comprehensive treatment.

A \emph{polynomial functor} is an endofunctor $p\colon\smset\to\smset$ of the form $\sum_{i: I}\yon^{A_i}$, where $I$ is the set of \emph{positions} and $A_i$ is the set of \emph{directions} at position $i$. Since $p(1)=\sum_{i:I}1=I$, we use the compact notation $p(1)$ for the position set and $p[i]\coloneqq A_i$ for the direction set at $i: p(1)$, so that
\begin{equation}\label{eqn.poly}
p=\sum_{i: p(1)}\yon^{p[i]}=\sum_{i: p(1)}\;\prod_{d: p[i]}\yon.
\end{equation}
We use juxtaposition for scalar multiplication: $pq\coloneqq p\times q$, and we write $1=\yon^0$. A \emph{monomial} is a polynomial of the form $M\yon^N=\sum_{m\in M}\yon^N$; we refer to $\yon^N$ as a \emph{representable} and to $M=M\yon^0$ as a \emph{constant}.

One can derive from the universal property of coproducts and the Yoneda lemma that natural transformations $\varphi\colon p\to q$ are given by:
\begin{equation}\label{eqn.poly_map}
\poly(p,q)=\prod_{i: p(1)}\;\sum_{j: q(1)}\;\prod_{e: q[j]}\; \sum_{d: p[i]}1.
\end{equation}
The outer $\prod_i\sum_j$ yields an ``forward on-positions'' function $\varphi_1\colon p(1)\to q(1)$, which is the component of $\varphi$ at $1:\smset$, and the inner $\prod_e\sum_d$ yields, for each $i:p(1)$, a ``backward on-directions'' function from the directions of $q$ at $\varphi_1(i)$ to the directions of $p$ at $i$; we denote this function
\[\bk{\varphi}_i\colon q[\varphi_1(i)]\to p[i],\]
writing the overline-bracket $\bk{(\,\cdot\,)}$ for the backward component throughout. We take the morphisms $p\to q$ in $\poly$ to be the natural transformations between the underlying functors.

For polynomials $p,q$, their \emph{Dirichlet product} is
\begin{equation}\label{eqn.dirichlet}
p\otimes q\coloneqq\sum_{(i,j):p(1)\times q(1)}\yon^{p[i]\times q[j]}\;.
\end{equation}
The unit is $\yon=\yon^1$. The Dirichlet product is symmetric monoidal and closed: the \emph{internal hom} is
\begin{equation}\label{eqn.dirichlet_hom}
\ihom{p,q}=\prod_{i: p(1)}\;\sum_{j: q(1)}\;\prod_{e: q[j]}\;\sum_{d: p[i]}\yon
\end{equation}
satisfying $\poly(r\otimes p,\, q)\cong\poly(r,\,\ihom{p,q})$ naturally in $r$. Note that $\ihom{p,q}(1)=\poly(p,q)$, as is immediate from comparing \eqref{eqn.dirichlet_hom} with \eqref{eqn.poly_map}. We can rewrite it as
\[
\ihom{p,q}\cong\sum_{\varphi: \poly(p,q)}\yon^{\sum_{i: p(1)}q[\varphi_1(i)]}
\]
A direction at $\varphi:\ihom{p,q}(1)$ consists of a pair $(i,e)$, where $i:p(1)$ is a $p$-position and $e:q[\varphi_1(i)]$ is a $q$-direction.

\subsection{Polynomial comonads as categories}\label{sec.comonads}

The \emph{substitution} (or \emph{composition}) \emph{product} on $\poly$ is defined by $(p\tri q)(S)\coloneqq p(q(S))$, with unit $\yon$. Using the expanded form of \eqref{eqn.poly}:
\begin{equation}\label{eqn.substitution}
p\tri q=\sum_{i: p(1)}\;\prod_{d: p[i]}\;\sum_{j: q(1)}\;\prod_{e: q[j]}\yon.
\end{equation}
A position of $p\tri q$ is a pair $(i,j)$ where $i: p(1)$ and $j\colon p[i]\to q(1)$; the direction set is $(p\tri q)[(i,j)]=\sum_{d:p[i]}q[j_d]$. In particular, viewing a set $S:\smset$ as the constant polynomial $S\yon^0$, substitution recovers evaluation: $p\tri S=p(S)$.

A \emph{$\tri$-comonoid} in $\poly$ is a comonoid for the substitution product: a polynomial $c$ equipped with a counit $\epsilon\colon c\to\yon$ and comultiplication $\delta\colon c\to c\tri c$ satisfying coassociativity and counitality. Since the substitution product is composition of endofunctors on $\smset$, a $\tri$-comonoid is equivalently a comonad on $\smset$ whose underlying endofunctor is polynomial; we call it a \emph{polynomial comonad}.

Polynomial comonads are---up to isomorphism---small categories. Writing $c=\sum_{i:c(1)}\yon^{c[i]}$, the positions $c(1)$ are the objects of some category $\cat{C}$, and for each object $i:c(1)$, the directions $c[i]$ are the morphisms of $\cat{C}$ with domain $i$:
\[
c[i]\coloneqq\sum_{j:c(1)}\cat{C}(i,j).
\]
The counit $\epsilon\colon c\to\yon$ picks out an identity morphism $\id_i:c[i]$ at each object $i$. The comultiplication $\delta\colon c\to c\tri c$ assigns codomains and composition: given an object $i$ and a morphism $f:c[i]$, the on-positions component of $\delta$ returns $\cod(f):c(1)$, and given a further morphism $g:c[\cod(f)]$, the on-directions component returns the composite $f\then g:c[i]$. The comonoid laws guarantee that $\cod(\id_i)=i$, that composition is unital, that $\cod(f\then g)=\cod(g)$, and that composition is associative.

\begin{proposition}[Polynomial comonads are categories {\cite{ahman2016directed}; see also \cite{niu2025polynomial}}]\label{prop.comonads_cats}
The above correspondence defines an isomorphism between polynomial comonads and small categories.
\end{proposition}

\begin{example}[The store comonad and the pair groupoid]\label{ex.store}
For any set $S$, the monomial $S\yon^S$ arises as the store comonad from the adjunction $S\times\blank\;\dashv\;(\blank)^S$ on $\smset$, i.e.\ the comonad with underlying endofunctor $X\mapsto S\times X^S$. We can also write
\begin{equation}\label{eqn.store_decomp}
S\yon^S\;\cong\;\ihom{S\yon,\yon}\otimes S\yon
\end{equation}
since $\ihom{S\yon,\yon}=\yon^S$, and $\yon^S\otimes S\yon\cong S\yon^{S\times 1}\cong S\yon^S$. The counit $\epsilon\colon S\yon^S\to\yon$ sends each position $s:S$ to the identity $\id_s\coloneqq s:S=c[s]$. The comultiplication $\delta\colon S\yon^S\to S\yon^S\tri S\yon^S$ acts as follows: on positions, $\delta$ sends $s:S$ to the pair $(s,\cod)$ where $\cod\colon S\to S$ is the identity---so $\cod(t)=t$ for all $t:S$. On directions, given $s:S$ and $t:c[s]=S$, the codomain is $\cod(t)=t$, and given a further $u:c[t]=S$, the composite is $t\then u\coloneqq u$.

By \cref{prop.comonads_cats} this is a category: it has object set $S$ and, for each pair $(s,t):S\times S$, exactly one morphism $s\to t$. The composition $t\then u=u$ says that composites are determined by the codomain of the second factor, as expected when every hom-set is a singleton. Every arrow $s\to t$ is an isomorphism with inverse $t\to s$, so this is the \emph{pair groupoid} (or \emph{codiscrete category}) on $S$.
\end{example}



\subsection{Polynomial coalgebras as dynamical systems}\label{sec.coalgebras}

For a polynomial $p$, a \emph{$p$-coalgebra} is a pair $(S,\beta)$ where $S:\smset$ is a set and $\beta\colon S\to p(S)$ is a function. We call elements of $S$ states; given a state $s:S$, the map assigns a \emph{readout} position $r_s\coloneqq\beta_1(s):p(1)$ and an \emph{update} function $\bk{\beta}_s\colon p[r_s]\to S$ sending each direction at $r_s$ to a new state. For example when $p=A\yon^B$ then a $p$-coalgebra $S\to AS^B$ is equivalently a readout function $S\to A$ and an update function $S\times B\to S$.

A $p$-coalgebra is thus a deterministic dynamical system whose interface is $p$: the state exposes an output in $p(1)$ and transitions in response to inputs in $p[\beta_1(s)]$, capturing Moore machines, discretized differential equations, Markov decision processes, and other mode-dependent dynamics in a single framework~\cite{niu2025polynomial}. It is a special case of the general theory of coalgebras for endofunctors; see~\cite{jacobs2017introduction} for a comprehensive treatment.

\begin{proposition}[Coalgebras as polynomial maps]\label{prop.coalg_as_poly_map}
For $p:\poly$ and $S:\smset$, there is a bijection
\[
\smset(S,\,p\tri S)\;\cong\;\poly(S\yon^S,\,p),
\]
natural in $p:\poly$ and in $S:\smset_\cong$.
\end{proposition}

\begin{proof}
A coalgebra $\beta\colon S\to p\tri S$ unpacks as a pair $\beta=(\beta_1,\beta^\sharp)\colon S\to\sum_{i:p(1)}\prod_{d:p[i]}S$: for each $s:S$, a position $\beta_1(s):p(1)$ and an update $\beta^\sharp_s\colon p[\beta_1(s)]\to S$. A polynomial map $\hat\beta\colon S\yon^S\to p$ unpacks as a pair $(\hat\beta_1,\hat\beta^\sharp)$ of identical signature by \eqref{eqn.poly_map}. The bijection sets $\hat\beta_1\coloneqq\beta_1$ and $\hat\beta^\sharp\coloneqq\beta^\sharp$.\footnote{
Equivalently, $\poly$ has a left Kan operation $\lens{p}{q}$ with $\poly(p,q'\tri q)\cong\poly(\lens{p}{q},q')$ given by $\lens{p}{q}=\sum_{i:p(1)}\yon^{q\tri p[i]}$, so $\lens{S}{S}=S\yon^S$.
}
Naturality in $p:\poly$ is just composition. For naturality in $S$, a bijection $f\colon S\To{\cong}T$ acts as $f\yon^{f\inv}\colon S\yon^S\to T\yon^T$ (positions $f$, directions $f\inv$), and we need to check
\[
\begin{tikzcd}
	S\ar[r,"\beta"]\ar[d, "f"']&p\tri S\ar[d, "p\tri f"]\\
	T\ar[r,"\beta'"']&p\tri T
\end{tikzcd}
\quad\leftrightsquigarrow\quad
\begin{tikzcd}
	S\yon^S\ar[r,"\hat\beta"]\ar[d, "f\yon^{f\inv}"']&
	p\ar[d, equal]\\
	T\yon^T\ar[r,"\hat{\beta'}"']&
	p
\end{tikzcd}
\]
The first commutes iff $\beta'_1\circ f=\beta_1$ and $f\circ\beta^\sharp_s=\beta'^\sharp_{f(s)}$ for all $s:S$; the second iff $\hat{\beta'}_1\circ f=\hat\beta_1$ and $\hat\beta^\sharp_s=f\inv\circ\hat{\beta'}^\sharp_{f(s)}$. Under the bijection $\hat\beta_\bullet=\beta_\bullet$, these coincide ($f$ being a bijection).\qedhere
\end{proof}

The Dirichlet product formula \eqref{eqn.dirichlet} gives $1\yon^1=\yon$ and $S\yon^S\otimes T\yon^T=(S\times T)\yon^{S\times T}$, so the assignment $S\mapsto S\yon^S$ extends to a strong symmetric monoidal functor $\smset_\cong\to\poly$, sending a bijection $f\colon S\To{\cong}T$ to $f\yon^{f\inv}$ (\cref{prop.coalg_as_poly_map}). Composing with the Dirichlet product gives a strong monoidal action of $\smset_\cong$ on $\poly$.

\begin{definition}[$\Fun{Store}$ and the store action]\label{def.store_action}
Write
\[
\Fun{Store}\colon(\smset_\cong,1,\times)\to(\poly,\yon,\otimes),\qquad S\longmapsto S\yon^S
\]
for the above functor, and define the \emph{store action} by
\[
(\cdot)\colon\smset_\cong\times\poly\to\poly
\where
S\cdot p\coloneqq S\yon^S\otimes p.
\qedhere
\]
\end{definition}


\subsection{Cotangent bundles as polynomials}\label{sec.cot}

The category of smooth manifolds provides the source of polynomials relevant to physics. Each manifold has a cotangent bundle, and each smooth map has a transpose on cotangent fibers; combining these gives a functor into $\poly$.

\begin{definition}[Cotangent functor]\label{def.cot}
The \emph{cotangent functor} $\cot\colon\mfd\to\poly$ sends a smooth manifold $M$ to
\begin{equation}\label{eqn.cot_def}
\cot{M}\coloneqq\sum_{m\in M}\yon^{T^*_mM}
\end{equation}
where $T^*_mM$ is the cotangent space at $m$. A smooth map $f\colon M\to N$ is sent to the $\poly$ map $\cot{f}\colon\cot{M}\to\cot{N}$ with forward part $\cot{f}_1=f$ and backward part $\bk{\cot{f}}_m\coloneqq(T_m f)^\top\colon T^*_{f(m)}N\to T^*_mM$.
\end{definition}

This is indeed functorial: the identity map $\id_M\colon M\to M$ has $T_m(\id_M)=\id_{T_mM}$, whose transpose is $\id_{T^*_mM}$, so $\cot{\id_M}=\id_{\cot{M}}$. For composable smooth maps $f\colon M\to N$ and $g\colon N\to P$, the chain rule gives $T_m(g\circ f)=T_{f(m)}g\circ T_m f$, so
\[
(T_m(g\circ f))^\top=(T_m f)^\top\circ(T_{f(m)}g)^\top=\bk{\cot{f}}_m\circ\bk{\cot{g}}_{f(m)},
\]
which is the backward part of $\cot{f}\then\cot{g}$.

\begin{proposition}\label{prop.cot_monoidal}
The cotangent functor $\cot\colon(\mfd,\mathbb{R}^0,\times)\to(\poly,\yon,\otimes)$ is strong symmetric monoidal:
\[
\cot{\mathbb{R}^0}\cong\yon
\qqand
\cot{M\times N}\cong\cot{M}\otimes\cot{N}.
\]
\end{proposition}

\begin{proof}
\textit{Unit.} The one-point manifold $\mathbb{R}^0=\{*\}$ has $T^*_*\mathbb{R}^0=0$, so $\cot{\mathbb{R}^0}=\sum_{*:1}\yon^0=\yon$.

\textit{Products.} Since $T$ preserves products, the cotangent space of a product splits as $T^*_{(m,n)}(M\times N)\cong T^*_mM\oplus T^*_nN$, hence as an underlying set is canonically $T^*_mM\times T^*_nN$. Thus
\begin{equation}\label{eqn.cot_times}
\cot{M\times N}=\sum_{(m,n)\in M\times N}\yon^{T^*_mM\times T^*_nN}=\cot{M}\otimes\cot{N}
\end{equation}
using the definition of the Dirichlet product \eqref{eqn.dirichlet}.

\textit{Naturality.} For a product map $f\times g\colon M\times N\to M'\times N'$, the fiber map is block-diagonal: $T_{(m,n)}(f\times g)=T_m f\oplus T_n g$. Hence $(T_{(m,n)}(f\times g))^\top=(T_m f)^\top\times(T_n g)^\top$, so the isomorphism $\cot{M\times N}\cong\cot{M}\otimes\cot{N}$ is natural in $M$ and~$N$.
%
%\textit{Coherence.} The associator and unitor for the Dirichlet product correspond to the canonical isomorphisms $(M\times N)\times P\cong M\times(N\times P)$ and $M\times 1\cong M$ in $\mfd$, so the coherence diagrams for a strong monoidal functor commute.
\end{proof}

\Cref{prop.cot_monoidal} uses only the underlying sets of the cotangent fibers, e.g.\ in \eqref{eqn.cot_times}. The linear structure is not lost, however: \cref{lem.cot_comonoid} shows that the diagonal and terminal maps in $\mfd$ induce fiberwise addition and zero on $\cot{M}$.

\begin{remark}
In fact, $\cot$ also preserves coproducts,
\[
\cot{M}+\cot{N}\cong\cot{M+N}.
\]
Hence $\cot\colon(\mfd,0,+,1,\times)\to(\poly,0,+,\yon,\otimes)$ is a distributive monoidal functor.
\end{remark}

\begin{lemma}\label{lem.cot_comonoid}
For any smooth manifold $M$, the polynomial $\cot{M}$ is a cocommutative comonoid in $(\poly,\yon,\otimes)$. The comultiplication and counit
\[
\cot{\Delta_M}\colon\cot{M}\to\cot{M}\otimes\cot{M}
\qqand
 \cot{!_M}\colon\cot{M}\to\yon
\]
sends each position $m\in M$ to $(m,m)$. On directions, each cotangent fibre $T^*_mM$ carries the induced cocommutative comonoid structure
\[
T^*_mM\times T^*_mM\to T^*_mM,\quad (\xi,\xi')\mapsto\xi+\xi'
\qqand 1\to T^*_mM,\quad *\mapsto 0,
\]
i.e.\ fiberwise addition and zero. Thus the $\otimes$-comonoid structure on $\cot{M}$ recovers the additive structure of each cotangent space.%
\footnote{
One can recover scalar structure on $T^*_mM$ as well. For any $M:\mfd$, there is a natural map
\[
(\cdot)\colon\cot{M}\to\yon^\rr\otimes\cot{M}
\]
Axioms relating the comonoid structure of $\cot{M}$ and various structures on $\rr$ enforce that the forward map is $x\mapsto (!,x)$ and that the backwards map $\bk{(\cdot)}_x\colon (r,\xi)\mapsto r\cdot\xi$ encodes the following laws:
\[
	1\cdot\xi=\xi,\qquad
	(r_1r_2)\cdot\xi=r_1\cdot(r_2\cdot\xi),\qquad
	0\cdot\xi=0,\qand
	r\cdot(\xi_1+\xi_2)=(r\cdot\xi_1)+(r\cdot\xi_2).
\]
Notably missing is the axiom $(r_1+r_2)\cdot\xi=(r_1\cdot\xi)+(r_2\cdot\xi)$, which seemingly cannot be encoded this way because there is seemingly no way to copy a cotangent vector ($\xi$).
}
\end{lemma}

\begin{proof}
Since $(\mfd,\mathbb{R}^0,\times)$ is cartesian, every manifold $M$ is a cocommutative comonoid via the diagonal $\Delta_M\colon M\to M\times M$ and $!_M\colon M\to\mathbb{R}^0$. Since $\cot$ is strong symmetric monoidal (\cref{prop.cot_monoidal}), it preserves comonoid structures. On positions this gives $m\mapsto(m,m)$ and $m\mapsto *$. For the backward maps: $T_m\Delta_M$ is the diagonal $T_mM\to T_mM\times T_mM$, $v\mapsto(v,v)$; its transpose is addition $T^*_mM\times T^*_mM\to T^*_mM$. The map $T_m(!_M)$ is $T_mM\to 0$, whose transpose is the map $1\to T^*_mM$ that selects the zero covector.
\end{proof}


\begin{proposition}[Hopf monoid structure on cotangent bundles of Lie groups]\label{prop.cot_hopf}
Let $G$ be a Lie group with unit $\eta\colon\mathbb{R}^0\to G$, multiplication $\mu\colon G\times G\to G$, and inversion $\iota\colon G\to G$. Then $\cot{G}$ is a Hopf monoid in $(\poly,\yon,\otimes)$: the monoid structure is
\[
\cot{\eta}\colon\yon\to\cot{G},\qquad \cot{\mu}\colon\cot{G}\otimes\cot{G}\to\cot{G},
\]
the comonoid structure is as in \cref{lem.cot_comonoid}, and the antipode is $\cot{\iota}\colon\cot{G}\to\cot{G}$.
\end{proposition}

\begin{proof}
The Lie group axioms are equations in $\mfd$; applying the strong monoidal functor $\cot$ preserves them. In particular, the monoid axioms for $(\eta,\mu)$ become monoid axioms for $(\cot{\eta},\cot{\mu})$, the compatibility between multiplication and comultiplication (which holds because $\Delta$ is natural), and the antipode condition $\mu\circ(\iota\times\id)\circ\Delta=\eta\circ!=\mu\circ(\id\times\iota)\circ\Delta$ maps to the Hopf monoid axiom in $\poly$. See also~\cite{spivak2023hopf}.
\end{proof}

Recall the inclusion $\pvect\to\vect\hookrightarrow\mfd$ (\cref{prop.pnla_smc_action}) and $\Fun{Store}(S)\coloneqq S\yon^S$ from \cref{def.store_action}.

\begin{proposition}\label{prop.pnla_polynomial}
The following square, in which all maps are strong monoidal, commutes up to isomorphism:
\begin{equation}\label{eqn.pnla_polynomial_square}
\begin{tikzcd}[ampersand replacement=\&]
\pvect\ar[r,"\Fun{inc}"]\ar[d,"|{\blank}|"']
\ar[dr,phantom,"\cong^\theta"description] \& \mfd\ar[d,"\cot"]\\
\smset_\cong\ar[r,"\Fun{Store}"'] \& \poly
\end{tikzcd}
\end{equation}
The 2-cell $\theta_V\colon V\yon^V\To{\cong}\cot{V}$ at each $V:\pvect$ is identity on positions and, on directions at $v\in V$, the cotangent flow \eqref{eqn.flow_cot}
\begin{equation}\label{eqn.directions_sharp}
T^*_vV\To{\eqref{eqn.tangent_vec}}V^*\To{\sharp_V}V\To{v+\blank}V,
\qquad
\xi\longmapsto v+\sharp_V(\xi).
\end{equation}
\end{proposition}

\begin{proof}
On positions, both functors send $V$ to its underlying set, so the on-positions component is the identity. At each $v\in V$, \eqref{eqn.directions_sharp} is the composite of two linear isomorphisms with translation by $v$, hence a bijection. So $\Fun{Store}(|V|)=V\yon^V\cong\cot{V}$ in $\poly$ for each $V$. For naturality at a $\pvect$-isomorphism $\phi\colon V\to W$, the on-positions square commutes by definition. On directions at $v$, we must check that the two composites $T^*_{\phi(v)}W\to V$ agree:
\begin{align*}
\xi\;&\Mapsto{\phi^*}\;\phi^*\xi\;\Mapsto{\sharp_V}\;\sharp_V(\phi^*\xi)\;\Mapsto{v+\blank}\;v+\sharp_V(\phi^*\xi), && \text{[$F^V_v\circ\cot{\phi}$]}\\
\xi\;&\Mapsto{\sharp_W}\;\sharp_W(\xi)\;\Mapsto{\phi(v)+\blank}\;\phi(v)+\sharp_W(\xi)\;\Mapsto{\phi^{-1}}\;v+\phi^{-1}(\sharp_W\xi) && \text{[$\phi\yon^\phi\circ F^W_{\phi(v)}$]}.
\end{align*}
The two are equal iff $\sharp_V\circ\phi^*=\phi^{-1}\circ\sharp_W$, which is \eqref{eqn.pairing_triangle}.\qedhere
\end{proof}

\begin{remark}\label{rmk.pnla_fiberwise_addition}
By \cref{prop.comonads_cats}, a $\tri$-comonoid structure on $V\yon^V$ is a category whose object set is $V$ and whose arrows out of each object are indexed by $V$. By \cref{ex.store}, the canonical such structure is the codiscrete category on $V$: there is exactly one morphism $v\to v'$ for each pair, and the direction-set $V$ at position $v$ is interpreted as the set of codomains.

Transporting this $\tri$-comonoid structure to $\cot{V}$ along the iso of \cref{prop.pnla_polynomial}, a covector $\xi\in T^*_vV$ at $v$ corresponds to the unique morphism $v\to v+\sharp_V(\xi)$ in the codiscrete category on $V$. The identity at $v$ corresponds to the zero covector $0\in T^*_vV$, since $v+\sharp_V(0)=v$. Composition of $\xi$ at $v$ with $\xi'$ at $v+\sharp_V(\xi)$ corresponds to $\xi+\xi'$ at $v$, since both produce the morphism with codomain $v+\sharp_V(\xi)+\sharp_V(\xi')=v+\sharp_V(\xi+\xi')$. Thus the pairing $\sharp_V$ identifies $\cot{V}$ with the codiscrete category on $V$, with codomain map $(v,\xi)\mapsto v+\sharp_V(\xi)$ and composition by fiberwise addition of covectors.\qedhere
\end{remark}

For any paired vector space $(V,\sharp_V)$, $T^*V\cong V\oplus V^*$ by \eqref{eqn.tangent_vec}. We write $(v,\xi)$ for a \emph{position} $v\in V$ and \emph{momentum} $\xi\in V^*$. The cotangent space decomposes as
\[
T^*_{(v,\xi)}(T^*V)\cong T^*_vV\oplus T^*_\xi(V^*)\cong V^*\oplus V,
\]
again by \eqref{eqn.tangent_vec}, using $(V^*)^*\cong V$ on the $V^*$-factor.

\begin{definition}[Legendre projection]\label{def.legendre_projection}
The \emph{Legendre projection} $\rho_V\colon\cot{T^*V}\to\cot{V}$ has position-action $(v,\xi)\mapsto v$ and direction-action at $(v,\xi)$
\[
T^*_vV\cong V^*\To{\xi_V\mapsto(\xi_V,\,\sharp_V(\xi))}V^*\oplus V\cong T^*_{(v,\xi)}(T^*V).
\qedhere
\]
\end{definition}

\begin{lemma}\label{lem.rho_natural}
The components $\rho_V$ assemble into a monoidal natural transformation
\[
\begin{tikzcd}[ampersand replacement=\&, column sep=50pt]
\pvect\ar[r,bend left=25,"\cot{T^*\blank}", ""{name=U,below}]\ar[r,bend right=25,"\cot{\blank}"', ""{name=D,above}]\& \poly.
\arrow[Rightarrow, from=U, to=D-|U, pos=.4, "\rho"]
\end{tikzcd}
\]
\end{lemma}

\begin{proof}
Naturality in $\pvect$-isomorphisms and monoidality both follow from $\sharp_{V_1\oplus V_2}=\sharp_{V_1}\oplus\sharp_{V_2}$ (\cref{prop.TT_monoidal}).
\end{proof}

\section{Dynamic category-like structures}\label{sec.dynamic_structures}

\subsection{Wiring-diagram operads}\label{sec.wd_operads}

Categories, monoidal categories, traced categories, operads, etc.\ can be understood in terms of the sorts of compositions they allow. For example:
\begin{equation}\label{eqn.wd_examples}
\begin{tikzpicture}[baseline=(Catg)]
\begin{scope}[font=\footnotesize, text height=1.5ex, text depth=.5ex]
  \begin{scope}[oriented WD, bb port sep=1, bb port length=0, bb min width=.4cm, bby=.2cm, inner xsep=.2cm, x=.5cm, y=.3cm]
  	\node[bb={1}{1}] (Catf) {$f_1$};
  	\node[bb={1}{1}, right=1 of Catf] (Catg) {$f_2$};
  	\node[bb={0}{0}, fit=(Catf) (Catg)] (Cat) {};
  	\node[coordinate] at (Cat.west|-Catf_in1) (Cat_in1) {};
  	\node[coordinate] at (Cat.east|-Catg_out1) (Cat_out1) {};
  	\draw (Cat_in1) -- (Catf_in1);
  	\draw (Catf_out1) -- (Catg_in1);
  	\draw (Catg_out1) -- (Cat_out1);
  	\draw[label]
  		node[left=3pt of Cat_in1] {$A$}
  		node at ($(Catf_out1)!.5!(Catg_in1)+(0,.8)$) {$B$}
  		node[right=3pt of Cat_out1] {$C$}
  	;
  %
  	\node[bb={1}{2}, above right=-1 and 6.5 of Catf] (Monf) {$f_1$};
  	\node[bb={2}{1}, below right=-.5 and 1 of Monf] (Mong) {$f_2$};
  	\node[bb={0}{0}, fit=(Monf) (Mong)] (Mon) {};
  	\node[coordinate] at (Mon.west|-Monf_in1) (Mon_in1) {};
  	\node[coordinate] at (Mon.west|-Mong_in2) (Mon_in2) {};
  	\node[coordinate] at (Mon.east|-Monf_out1) (Mon_out1) {};
  	\node[coordinate] at (Mon.east|-Mong_out1) (Mon_out2) {};
  	\draw  (Mon_in1) -- (Monf_in1);
  	\draw (Monf_out1) -- (Mon_out1);
  	\draw (Monf_out2) to (Mong_in1);
  	\draw (Mon_in2) -- (Mong_in2);
  	\draw (Mong_out1) -- (Mon_out2);
  	\draw[label]
  		node[left=3pt of Mon_in1] {$A$}
  		node[left=3pt of Mon_in2] {$B$}
  		node[right=3pt of Mon_out1] {$C$}
  		node at ($(Monf_out2)!.5!(Mong_in1)+(.2,.4)$) {$D$}
  		node[right=3pt of Mon_out2] {$E$}
  	;
  %
  	\node[bb={2}{1}, below right= -1.75 and 6.5 of Monf] (Trf) {$f_1$};
  	\node[bb={2}{2}, below right=-1 and 1 of Trf] (Trg) {$f_2$};
  	\node[bb={0}{0}, fit={($(Trf.north west)+(-.5,2)$) ($(Trg.south east)+(.5,0)$)}] (Tr) {};
  	\node[coordinate] at (Tr.west|-Trf_in2) (Tr_in1) {};
  	\node[coordinate] at (Tr.west|-Trg_in2) (Tr_in2) {};
  	\node[coordinate] at (Tr.east|-Trf_out1) (Tr_out1) {};
  	\node[coordinate] at (Tr.east|-Trg_out2) (Tr_out2) {};
  	\draw (Tr_in1) -- (Trf_in2);
  	\draw (Trf_out1) to (Trg_in1);
  	\draw (Tr_in2) -- (Trg_in2);
  	\draw (Trg_out2) -- (Tr_out2);
  	\draw let \p1=(Trg.east), \p2=(Trf.north west), \n1=\bbportlen, \n2=\bby in
  		(Trg_out1) to[in=0] (\x1+\n1,\y2+\n2) -- (\x2-\n1,\y2+\n2) to[out=180] (Trf_in1);
  	\draw[label]
  		node[left=3pt of Tr_in1] {$A$}
  		node[left=3pt of Tr_in2] {$B$}
  		node at ($(Trf_out1)!.4!(Trg_in1)+(0,1)$) {$D$}
  		node[right=3pt of Tr_out2] {$E$}
  		node at ($(Trf.north)+(0,1.5)$) {$C$}
  	;
  \end{scope}
  \begin{scope}[circuit logic US, thick, every to/.style={out=0,in=180}]
  	\node[and gate, draw, above right=-.2 and 3 of Trf] (Opdf) {$f_1$};
  	\node[and gate, draw, below right=-.1 and 0.5 of Opdf] (Opdg) {$f_2$};
  	\node[and gate, inner ysep=-4pt, draw, fit={($(Opdf.north west)+(0,-.2)$) ($(Opdg.south east)+(0,.2)$)}] (Opd) {};
  	\node[coordinate] at (Opd.input 1|-Opdf.west) (Opd_in1) {};
  	\node[coordinate] at (Opd.input 1|-Opdg.input 2) (Opd_in2) {};
  	\node[coordinate] at (Opd.output) (Opd_out) {};
  	\draw (Opd_in1) to (Opdf.west);
  	\draw (Opd_in2) to (Opdg.input 2);
  	\draw (Opdf.output) to (Opdg.input 1);
  	\draw (Opdg.output) to (Opd_out);
  	\draw[label]
  		node[left=3pt of Opd_in1] {$A$}
  		node[left=3pt of Opd_in2] {$B$}
  		node at ($(Opdf.output)!.5!(Opdg.input 1)+(.1,.25)$) {$C$}
  		node[right=3pt of Opd_out] {$D$}
  	;
  \end{scope}
	\node[below=.3 of Cat.south] (Cat name) {category};
	\node[text width=1.5cm, align=center] at (Cat name-|Mon) {monoidal\\category};
	\node[text width=2.5cm, align=center] at (Cat name-|Tr) {traced monoidal\\category};
	\node at (Cat name-|Opd) {operad};
\end{scope}
\end{tikzpicture}
\end{equation}
\smallskip

\noindent In each case there is an associated operad $\cat{W}$ whose objects are the legal interfaces (boxes-with-objet-labeled-ports) and whose morphisms $\phi\colon b_1,\ldots,b_K\to c$ are the legal arrangements of boxes in a box. That is, there are operads $\cat{W}_{\tn{O-Cat}}$, $\cat{W}_{\tn{O-MnCat}}$, $\cat{W}_{\tn{O-TrCat}}$, and $\cat{W}_{\tn{O-Opd}}$ each with a 2-ary morphism that can be drawn as above.%
\footnote{
There is a potentially-confusing level-shift in that there is an operad $\cat{W}_{\tn{O-Opd}}$ for which algebras are operads with object set $O$. In case it's helpful, the object set for $\cat{W}_{\tn{O-Opd}}$ would be $\List(O)\times O$, because every ``box'' has a list of inputs and an output.
}

A $\cat{W}$-algebra $\cat{C}\colon\cat{W}\to\smset$ assigns to each box $b$ the set $\cat{C}(b)$ of ``fillers'' and to each arrangement $\phi$ the function $\cat{C}(\phi)\colon \cat{C}(b_1)\times\cdots\times \cat{C}(b_K)\to \cat{C}(c)$ sending fillers for little boxes to a filler of the big box.%
\footnote{
By a $\cat{W}$-algebra, we mean an operad functor $\cat{W}\to\smset$, where $\smset$ is the operad (multicategory) underlying the monoidal category $(\smset,1,\times)$. In general, the \emph{underlying operad} of a symmetric monoidal category $(\cat{C},I,\otimes)$ has the same objects as $\cat{C}$ and multimorphisms $\cat{C}(a_1,\ldots,a_n;b)\coloneqq\cat{C}(a_1\otimes\cdots\otimes a_n,b)$, with operadic composition by tensoring and composing in $\cat{C}$.
}

\begin{example}
For a set $O$, let $\ob\cat{W}_{\tn{O-Cat}}\coloneqq O\times O$ and
\[\cat{W}_{\tn{O-Cat}}((o_0',o_1),(o_1',o_2),\ldots(o_{K-1}',o_K);(o_0,o_K'))\coloneqq
\begin{cases}
\{1\}&\tn{ if }o_k=o_k'\quad \forall\; 0\leq k\leq K\\
\varnothing&\tn{ otherwise}
\end{cases}
\]
Thus the first diagram shown in \eqref{eqn.wd_examples} represents a map $\phi\colon (A,B),(B,C)\to(A,C)$ in $\cat{W}_{\tn{O-Cat}}$ for some $O$ with $A,B,C:O$. If $\cat{C}\colon\cat{W}_{\tn{O-Cat}}\to\smset$ is an algebra, then the above diagram shows $f_1:\cat{C}(A,B)$, $f_2:\cat{C}(B,C)$ and $\phi(f_1,f_2):\cat{C}(A,C)$. In other words the diagram $\phi$ represents composition in a category. The category of small categories with object set $O$ is equivalent to the category of $\cat{W}_{\tn{O-Cat}}$-algebras.
\end{example}

A variety of examples of these ideas---wiring-diagram operads whose algebras are structured categories---can be found in the literature. \cite{Spivak.Schultz.Rupel:2016a} treats the traced and compact cases, showing in particular that the operad $\cat{W}_{\tn{TrCat}}$ for traced categories is the operad $1\tn{-}\Cat{Cob}$ of oriented 1-dimensional cobordisms. \cite{fong2019hypergraph} treats hypergraph categories, whose associated operad $\cat{W}_{\tn{HypCat}}$ is given by cospans of finite sets. The symmetric monoidal case~\cite{patterson2021wiring} is in some sense the most basic of the three, but its operad $\cat{W}_{\tn{MnCat}}$ is less easy to describe directly.

An advantage of the operadic perspective is that it allows us to more uniformly treat various sorts of category, including ones that may not have an existing name. For example, 
consider the notion of ``traced monoidal categories without identity''; one might balk at this because they wouldn't know whether the ``yanking'' axiom could so easily just be discarded. But as an operad it is fine: it is the suboperad comprising the left class of a certain factorization system on $1\tn{-}\Cat{Cob}$.

Another advantage is that wiring diagram algebras also easily represent $\cat{V}$-enriched categories for monoidal categories $\cat{V}$, namely by using algebras $\cat{W}\to\cat{V}$ valued in $\cat{V}$. For example, a $\vect$-enriched category with object set $O$ is an opeard functor $\cat{W}_{\tn{O-Cat}}\to\vect$.


One can build \emph{dynamic} versions of the structures in \eqref{eqn.wd_examples}, in which the ``fillers'' of a box are ``morphisms that change in time'', i.e.\ coalgebras.  In \cite{shapiro2022dynamic}, a dynamic category is a category enriched in $\org$; similarly for dynamic monoidal categories or dynamic operads. Given the ideas above, these are just operad functors
\[
\cat{W}_{\tn{O-Cat}}\to\org
\qquad
\cat{W}_{\tn{O-MnCat}}\to\org
\qquad
\cat{W}_{\tn{O-Opd}}\to\org
\] 
for any $O:\smset$. Our next goal then is to recall $\org$; in \cref{ch.potentials} we will talk about our choice of $\cat{W}$.

\subsection{The monoidal bicategory $\org$}\label{sec.org}

We recall the monoidal bicategory $\org$ from~\cite{spivak2021learners}; see also~\cite{shapiro2022dynamic}. The objects of $\org$ are polynomials, and for polynomials $p,q$ the hom-category is
\[\org(p,q)\coloneqq\ihom{p,q}\tn{-}\Cat{Coalg}.\]
Its monoidal structure is induced by the Dirichlet product on polynomials. A morphism $p\to q$ in $\org$ is an $\ihom{p,q}$-coalgebra $(S,\,\beta\colon S\to\ihom{p,q}(S))$, specifying for each $s:S$ an action $\act(s):\poly(p,q)$ and an update $\upd_s\colon\ihom{p,q}[\act(s)]\to S$ assigning a successor state to each direction.

\begin{definition}\label{def.pq_coalg}
The category $\ihom{p,q}\coalg$ has as objects pairs $(S,\beta)$ where $S$ is a set and $\beta\colon S\to\ihom{p,q}\tri S$ is a function, and where a morphism from $(S,\beta)$ to $(S',\beta')$ is a function $f\colon S\to S'$ such that $(\ihom{p,q}\tri f)\circ\beta=\beta'\circ f$. We refer to $S$ as the \emph{state set}.
\end{definition}

Unwinding this using \eqref{eqn.dirichlet}, $\beta=(\act^\beta,\upd^\beta)$ decomposes as an \emph{action} function
\[\act^\beta\colon S\to\poly(p,q)=\ihom{p,q}(1)\]
and, for each state $s:S$, an \emph{update} function
\[\upd^\beta_s\colon\sum_{I:p(1)}q\left[\act^\beta_s(I)\right]\to S.\]
For a state $s:S$ and position $I:p(1)$ we often write $\act^\beta_s\colon p\to q$ and $\upd^\beta_s(I)\colon q[\act^\beta_s(I)]\to S$. We may suppress $\beta$ when clear from context.

\begin{proposition}[\cite{spivak2021learners,shapiro2022dynamic}]\label{prop.org}
The hom-categories $\ihom{p,q}\coalg$ assemble into a monoidal bicategory $\org$ with polynomials as objects, with monoidal structure induced by the Dirichlet product on polynomials.
\end{proposition}

This structure arises from the fact that $\coalg\colon(\poly,\yon,\otimes)\to(\smcat,1,\times)$ is lax monoidal:
\[
1\to\yon\coalg
\qqand
p\coalg\times q\coalg\to(p\otimes q)\coalg.
\]
Thus the structure maps $\yon\to[p,p]$ and $[p,q]\otimes[q,r]\to[p,r]$ induce the identity and composition functors
\[
1\to[p,p]\coalg
\qqand
[p,q]\coalg\times[q,r]\coalg\to[p,r]\coalg.
\]

\subsection{The Para construction}\label{sec.para_general}

We use the following form of the Para construction. If a symmetric monoidal category $\cat{A}$ acts on a category $\cat{D}$, the category $\para{\cat{A}}{\cat{D}}$ has the same objects as $\cat{D}$, and a morphism $x\to y$ is a pair $(a,f)$ with $a:\cat{A}$ and
\[
f\colon a\cdot x\to y
\]
in $\cat{D}$. We call such morphisms \emph{parameterized maps} in $\cat D$. The identity on $x$ has parameter $I$ and the action unitor $I\cdot x\cong x$. The composite of $(a,f)\colon x\to y$ and $(b,g)\colon y\to z$ has parameter $b\otimes a$, or $a\otimes b$ after the symmetry, and uses
\begin{equation}\label{eqn.para_map}
(b\otimes a)\cdot x\cong b\cdot(a\cdot x)\To{b\cdot f}b\cdot y\To{g}z.
\end{equation}

When $\cat A$ and $\cat D$ are both symmetric monoidal and the action is monoidal, $\para{\cat A}{\cat D}$ inherits a symmetric monoidal structure: on objects, $x\otimes y$ is the $\cat D$-tensor, and the tensor of $(a,f)\colon x\to y$ and $(b,g)\colon x'\to y'$ has parameter $a\otimes b$ and underlying map
\begin{equation}\label{eqn.para_tensor}
(a\otimes b)\cdot(x\otimes x')\To{\cong}(a\cdot x)\otimes(b\cdot x')\To{f\otimes g}y\otimes y',
\end{equation}
where the first isomorphism comes from the action being monoidal.

The Para construction extends to a bicategory: a 2-cell $(a,f)\Rightarrow(a',f')$ in $\para{\cat A}{\cat D}(x,y)$ is a morphism $\alpha\colon a\to a'$ in $\cat A$ such that $f'\circ(\alpha\cdot x)=f$. When $\cat A$ is a groupoid, all 2-cells are invertible, and each hom-category $\para{\cat A}{\cat D}(x,y)$ is a groupoid.

\begin{proposition}[Para by strong monoidal functors]\label{prop.para_strong_induced}
A strong monoidal functor $J\colon\cat{A}\to\cat{D}$ induces an action of $\cat{A}$ on $\cat{D}$ by left multiplication, and a monoidal natural transformation $\theta\colon J'F\Rightarrow GJ$ as on the left, with $F,J,J'$ strong monoidal and $G$ lax (resp.\ strong) monoidal, induces a lax (resp.\ strong) monoidal action square as on the right:
\[
\begin{tikzcd}[ampersand replacement=\&]
\cat A\ar[r,"J"]\ar[d,"F"', ""{name=left}]\& \cat D\ar[d,"G", ""' {name=right}]\\
\cat A'\ar[r,"J'"']\& \cat D'
\arrow[Rightarrow, from=left, to=right, shorten <=3mm, shorten >=3mm, "\theta"]
\end{tikzcd}
\quad\leadsto\quad
\begin{tikzcd}[ampersand replacement=\&]
\cat A\times\cat D\ar[r,"(\cdot)"]\ar[d,"F\times G"', ""{name=left}]\& \cat D\ar[d,"G", ""' {name=right}]\\
\cat A'\times\cat D'\ar[r,"(\cdot')"']\& \cat D'.
\arrow[Rightarrow, from=left, to=right, shorten <=6mm, shorten >=6mm]
\end{tikzcd}
\]
\end{proposition}

\begin{proposition}[Para functoriality from action squares {\cite[\S 3.6, eq.~(3.59)]{capucci2024actegories}}]\label{prop.para_square}
Suppose $\cat A,\cat A'$ are symmetric monoidal categories acting on $\cat D,\cat D'$, that $F\colon\cat A\to\cat A'$ is a strong monoidal functor, $G\colon\cat D\to\cat D'$ is a functor, and $\theta\colon F(a)\cdot G(x)\to G(a\cdot x)$ is a natural transformation compatible with the action unitors and associators:
\begin{equation}\label{eqn.para_square}
\begin{tikzcd}[ampersand replacement=\&]
\cat A\times\cat D\ar[r,"(\cdot)"]\ar[d,"F\times G"', ""{name=left}]\& \cat D\ar[d,"G", ""' {name=right}]\\
\cat A'\times\cat D'\ar[r,"(\cdot')"']\& \cat D'.
\arrow[Rightarrow, from=left, to=right, shorten <=6mm, shorten >=6mm, "\theta"]
\end{tikzcd}
\end{equation}
Then there is an induced functor
\[
\para{F}{G}\colon\para{\cat A}{\cat D}\;\longrightarrow\;\para{\cat A'}{\cat D'}.
\]
When $\cat D,\cat D'$ are symmetric monoidal, the actions are monoidal, $G$ is strong (resp.\ lax) symmetric monoidal, and $\theta$ is monoidal, then $\para{F}{G}$ extends to a strong (resp.\ lax) monoidal functor.
\end{proposition}

\begin{proof}
The hypothesis is a 1-cell in the 2-category $\Cat{Act}$ of actegories \cite[\S 3.6, eq.~(3.59)]{capucci2024actegories}; the induced functor records the action of $\Cat{Para}$ on it.\qedhere
\end{proof}

\subsection{Coalgebras as Para over the store action}\label{sec.org_as_para}

The Para construction over the store action of $\smset_\cong$ on $\poly$ (\cref{def.store_action}) repackages the truncation $\org_{\cong}$ (with only coalgebra isomorphisms as 2-cells) exactly:

\begin{proposition}\label{prop.org_as_para}
There is an isomorphism of categories
\[
\para{\Cat{Set}_{\cong}}{\poly}\To{\cong}\org_\cong
\]
\end{proposition}

\begin{proof}
For each category---$\para{\Cat{Set}_{\cong}}{\poly}$ and $\org_\cong$---the object set is $\ob\poly$ and the mapping groupoids $p\to q$ agree
\begin{align*}
	\para{\Cat{Set}_{\cong}}{\poly}(p,q)&=
  \sum_{S:\smset_\cong}\poly(S\yon^S\otimes p, q)\\&\cong
  \sum_{S:\smset_\cong}\poly(S\yon^S,[p,q])\\&\cong
  \sum_{S:\smset_\cong}\poly(S,[p,q]\tri S)=\org_\cong(p,q)
\end{align*}
by \cref{eqn.dirichlet_hom,prop.coalg_as_poly_map}. It remains to show that the identities and composites agree. 

In both cases the identity is indexed by the monoidal unit $1$, and $1\To{\id}\poly(p,p)=[p,p](1)$ and $\yon\To{\id}[p,p]$ are identified as above. Given composable maps $p\To{(S,\varphi)}q\To{(T,\psi)}r$, in both cases the composite is of the form $(TS,\psi\circ\phi)$, and it is easy to see too are identified as above.
\end{proof}

\begin{corollary}[Cotangent learners]\label{cor.cotangent_learners}
There is a strong symmetric monoidal functor
\begin{equation}\label{eqn.cotangent_learners}
\Psi\colon\para{\pvect}{\mfd}\to\org.
\end{equation}
It sends $M\mapsto\cot{M}$ and a morphism $(V, \sharp,f\colon V\times M\to N)$ to a coalgebra
\[
V\To{\Psi(V,f)}[\cot{M},\cot{N}]\tri V
\]
which at $v\in V$ sends a point $m\in M$ to  $n\coloneqq f(v,m)\in N$ and for a covector $\xi_N\in T^*_nN$ if we let $(\xi_V,\xi_M)\coloneqq(T_{(v,m)} f)^\top(\xi_N)\in T_v^*V\times T_m^*M$, returns $\xi_M\in T^*_mM$ and updates to
\begin{equation}\label{eqn.cotangent_learner_update}
\upd_v(m,\xi_N)\coloneqq v+\sharp_V(\xi_V).
\end{equation}
\end{corollary}

\begin{proof}
By \cref{prop.pnla_polynomial,prop.para_square,prop.para_strong_induced,prop.org_as_para} we obtain a functor $\para{\pvect}{\mfd}\to\org$. On a map $f\colon V\times M\to N$, one obtains a polynomial map
\[
V\yon^V\otimes\cot{M}\;\To{\theta_V\otimes\id}\;\cot{V}\otimes\cot{M}\;\To{\cong}\;\cot{V\times M}\;\To{\cot{f}}\;\cot{N},
\]
and the rest of the statement is proved simply by tracing through the above propositions. In particular, \eqref{eqn.cotangent_learner_update} comes from \eqref{eqn.directions_sharp}.
\end{proof}

\subsection{Example: deep learning as a dynamic monoidal category}\label{sec.deep_learning}

We recast the deep-learning example of~\cite{shapiro2022dynamic} as the restriction of $\Psi$ \eqref{eqn.cotangent_learners} to Euclidean spaces. Take the object set $O\coloneqq\nn$ and assign to each $n\in\nn$ the polynomial
\[
\Psi(\rr^n)=\cot{\rr^n}\cong{\cot{\rr}}\tpow n.
\]
A layer $m\to n$ is a morphism $\rr^m\to\rr^n$ in $\para{\pvect}{\mfd}$: a parameter space $V=\rr^k$, regarded as a paired vector space with $\sharp_V=-\eta\cdot\id$ for a learning rate $\eta>0$, together with a smooth architecture
\[
f\colon V\times\rr^m\to\rr^n.
\]

Applying $\Psi$ gives an $\ihom{\cot{\rr^m},\cot{\rr^n}}$-coalgebra with state set $V$. At $v\in V$, its action is
\[
\cot{f_v}\colon\cot{\rr^m}\to\cot{\rr^n},
\]
whose forward part is $x\mapsto f_v(x)$ and whose backward part is $(T_x f_v)^\top$.

The update at $v$ on a direction $(x,\xi)\in\rr^m\times T^*_{f_v(x)}\rr^n$ is\textbf{[decide on notation: plain $\xi$ here vs $\xi_N$ for consistency with \cref{cor.cotangent_learners}]}
\[
\upd_v(x,\xi)=v+\sharp_V\bigl((\partial_Vf)_{(v,x)}^\top\xi\bigr)=v-\eta\cdot(\partial_Vf)_{(v,x)}^\top\xi,
\]
by \eqref{eqn.cotangent_learner_update}. When $\xi=dL|_{f_v(x)}$ for a loss $L\colon\rr^n\to\rr$, the chain rule gives $(\partial_Vf)_{(v,x)}^\top\xi=d_v(L\circ f(\blank,x))$; under the Euclidean identification, $\upd_v$ is one step of stochastic gradient descent.

Since $\Psi$ is a functor, composition in $\org$ reproduces backpropagation. Since $\Psi$ is strong monoidal, parallel composition of layers is transported by the comparison $\cot{\rr^m}\otimes\cot{\rr^{m'}}\cong\cot{\rr^{m+m'}}$. Thus the Euclidean restriction of $\Psi$ assembles into an operad functor
\[\cat{L}\colon\cat{W}_{\nn\tn{-MnCat}}\to\org.\]

\chapter{Potentialized lenses}\label{ch.potentials}

In this chapter we define a symmetric monoidal category $\potlens$ of potentialized lenses. \Cref{sec.lenses_general} develops lenses, their wiring-diagram operations, and the backward-monad construction. \Cref{sec.abstract_representation} packages the representation functors. \Cref{sec.potlens_manifolds} specializes the source-side constructions to manifolds: manifold lenses, the potential monad associated to any Lie monoid, and the PN Lie Algebra parameterization, synthesizing them as $\potlens$.

\section{The abstract construction}\label{sec.lenses_general}

\subsection{Lenses}\label{subsec.lenses}
We first recall the category $\Lens{\cat{C}}$ of lenses in a symmetric monoidal category $(\cat{C},I,\otimes)$; see \cite{hedges2017coherence,spivak2019generalized}. An object is a pair $\lensob{c}$ in which $\outp c$ carries a commutative comonoid structure $(\delta_{\outp c},\varepsilon_{\outp c})$. A morphism $f\colon\lensob{c}\to\lensob{d}$ consists of a pair of maps $f=(\outp f,\inpt f)$ where
\begin{equation}\label{eqn.lens_def}
\outp f\colon \outp c\to\outp d
\qqand
\inpt f\colon \inpt d\otimes\outp c\to\inpt c,
\end{equation}
and $\outp f$ is a comonoid homomorphism.%
\footnote{
When $(\cat{C},1,\times)$ is cartesian monoidal, every object has a unique comonoid structure and every morphism is automatically a comonoid homomorphism~\cite{fox1976coalgebras}.
}
These can be depicted as follows:
\[
\begin{tikzpicture}[oriented WD, bb port length=5pt, bb port sep=1, bb min width=.4cm, bbx=.5cm, bby=.7ex, baseline=(fo)]
  \node[bb={1}{1}] (fo) {$\outp f$};
  \node[bb={1}{2}, above left=6 and 1.5 of fo] (fi) {$\inpt f$};
  \node[circle, draw, fill=black, inner sep=0, minimum size=3pt, left=1 of fo] (dot) {};
  \coordinate (co) at ($(dot)+(-3,0)$);
  \coordinate (do) at ($(fo.east)+(1,0)$);
  \coordinate (ci) at (co |- fi_in1);
  \coordinate (di) at (do |- fi_out1);
  \node[left=0 of co] {$\outp c$};
  \node[right=0 of do] {$\outp d$};
  \node[left=0 of ci] {$\inpt c$};
  \node[right=0 of di] {$\inpt d$};
  \draw[ar] (co) to (dot);
  \draw[ar] (dot) to (fo_in1);
  \draw[ar] (fo_out1) to (do);
  \draw[ar] (fi_in1) to (ci);
  \draw[ar] (di) to (fi_out1);
  \draw[ar] (dot) to[out=30, in=0] (fi_out2);
\end{tikzpicture}
\]
Morphisms compose as follows: given also $g\colon \lensob{d}\to\lensob{e}$, the composite $f\then g$ is given by
\[
\outp{(f\then g)}=\outp f\then\outp g
\qqand
\inpt{(f\then g)}=
	(\inpt{e}\otimes\delta_{\outp{c}})\then
	(\inpt{e}\otimes\outp{f}\otimes\outp{c})\then
	(\inpt{g}\otimes\outp{c})\then
	\inpt{f}.
\]
This dense notation can be depicted in a string-diagram as follows:
\begin{equation}\label{eqn.lens_pic}
\begin{tikzpicture}[oriented WD, bb port length=5pt, bb port sep=1, bb min width=.4cm, bbx=.5cm, bby=.7ex, baseline=(fo)]
  \node[bb={1}{1}] (fo) {$\outp f$};
  \node[bb={1}{1}, right=5 of fo] (go) {$\outp g$\vphantom{$\inpt f$}};
  \node[bb={1}{2}, above left=6 and 1.5 of fo] (fi) {$\inpt f$};
  \node[bb={1}{2}, above left=6 and 1.5 of go] (gi) {$\inpt g$\vphantom{$\inpt f$}};
  \node[circle, draw, fill=black, inner sep=0, minimum size=3pt, left=1 of fo] (dot1) {};
  \node[circle, draw, fill=black, inner sep=0, minimum size=3pt] at ($(fo.east)!.75!(go.west)$) (dot2) {};
  \coordinate (co) at ($(dot1)+(-3,0)$);
  \coordinate (eo) at ($(go.east)+(1,0)$);
  \coordinate (ci) at (co |- fi_in1);
  \coordinate (ei) at (eo |- gi_out1);
  \node[left=0 of co] {$\outp c$};
  \node[right=0 of eo] {$\outp e$};
  \node[left=0 of ci] {$\inpt c$};
  \node[right=0 of ei] {$\inpt e$};
  \draw[ar] (co) to (dot1);
  \draw[ar] (dot1) to (fo_in1);
  \draw[ar] (fo_out1) to node[above, font=\scriptsize] {$\outp d$} (dot2);
  \draw[ar] (dot2) to (go_in1);
  \draw[ar] (go_out1) to (eo);
  \draw[ar] (fi_in1) to (ci);
  \draw[ar] (ei) to (gi_out1);
  \draw[ar] (gi_in1) to[out=180, in=0] node[above, font=\scriptsize] {$\inpt d$} (fi_out1);
  \draw[ar] (dot1) to[out=30, in=0] (fi_out2);
  \draw[ar] (dot2) to[out=30, in=0] (gi_out2);
\end{tikzpicture}
\end{equation}
The identity on $\lensob{c}$ is $(\id_{\outp c},\id_{\inpt c}\otimes\varepsilon_{\outp c})$, i.e.\ counit in the second argument. The category $\Lens{\cat{C}}$ inherits a monoidal structure with unit $\binom{I}{I}$ and
\[
\lensob{c}\otimes\lensob{d}\coloneqq\lensobs{c}{d}.
\]

Our main example will be $\cat{C}=\mfd$ with its cartesian product. But the following is simple and instructive.
\begin{example}\label{ex.lens_finsetop}
The category $\finset\op$, opposite to the category of finite sets, is cartesian monoidal where the product is given by the disjoint union coproduct $(0,+)$, i.e.\ the coproduct in $\smset$. A comonoid in $\smset\op$ is a monoid in $\smset$, so an object in $\Lens{\finset\op}$ is a pair of sets $\binom{\inpt A}{\outp A}$ where $\outp A$ is a monoid. 

A lens $\binom{\inpt A}{\outp A}\to\binom{\inpt B}{\outp B}$ consists of two maps
\begin{equation}\label{eqn.lens_finsetop}
\outp A\From{\outp f} \outp B
\qqand
\outp A+\inpt B\From{\inpt f} \inpt A,
\end{equation}
where $\outp f\colon\outp B\to\outp A$ is a monoid homomorphism in $\finset$.
\end{example}

\begin{proposition}
The $\Lens{-}$ construction is functorial in symmetric monoidal categories and strong monoidal functors.
\end{proposition}
\begin{proof}
If $F\colon\cat{C}\to\cat{D}$ is colax monoidal, it preserves comonoid structures, so we can define $\Lens F$ on objects. If $F$ is lax monoidal, it has a productor $F_2\colon F(A)\otimes F(B)\to F(A\otimes B)$ and we can define $F$ on morphisms: given a lens morphism $(\outp f,\inpt f)$ we set $\Lens F(\outp f,\inpt f)\coloneqq(F(\outp f),\,F(\inpt f)\circ F_2)$. For functoriality the lax and colax structures must cohere, i.e. the lax structure must be inverse to the colax structure; we leave the details to the reader.
\qedhere
\end{proof}


\subsection{Lenses represent serial, parallel, feedback, and splitting}\label{sec.plenspot_ops}

Recall from \cref{sec.dynamic_structures} that operads---and hence monoidal categories---can be used to encode compositional syntax. In fact, the monoidal category $\Lens{\cat{C}}$ of lenses in a symmetric monoidal category $\cat{C}$  does just that for ``traced symmetric monoidal semi-categories'', as we now explain.

By a \emph{semi-category} we mean almost a category, except no identities (like a semigroup is a group without identities). So our goal is to show that $\Lens{C}$ interprets diagrams that look like this:
\begin{equation}\label{eqn.arb_wd}
\varphi\coloneqq
\begin{tikzpicture}[oriented WD, bb min width=.6cm, bb port sep=1, bbx=.6cm, bby=1ex, bb port length=2.5pt, baseline=(X2)]
  \node[bb port sep=2, bb={2}{2}] (X1) {};
  \node[bb={1}{1}, right=.7 of X1_out1] (X2) {};
  \node[bb={0}{0}, fit={(X1) (X2) ($(X1.north)+(0,2.5)$)}] (Y) {};
  \node[coordinate] at (Y.west|-X1_in2) (Y_in1) {};
  \node[coordinate] at (Y.east|-X2_out1) (Y_out1) {};
  \node[coordinate] at (Y.east|-X1_out2) (Y_out2) {};
  \draw[ar] (Y_in1) -- node[above, font=\scriptsize] {$A$} (X1_in2);
  \draw[ar] (X1_out1) -- node[above, font=\scriptsize] {$B$} (X2_in1);
  \draw[ar] (X2_out1) -- node[above, font=\scriptsize] {$C$} (Y_out1);
  \draw[ar] (X1_out2) -- node[below, font=\scriptsize] {$D$} (Y_out2);
  \draw[ar] let \p1=(X2.north east), \p2=(X1.north west), \n1={\y2+\bby}, \n2=\bbportlen in
          (X2_out1) to[in=0,in looseness=.9] (\x1+.7*\n2,\n1) -- node[above, font=\scriptsize] {$C$} (\x2-.7*\n2,\n1) to[out=180] (X1_in1);
  \end{tikzpicture}
 \end{equation}

Looking at the two inner boxes and the outer box, we see that this diagram should be represented in $\Lens{\cat{C}}$ by a map of the form
\[
\varphi\colon\binom{C\otimes A}{B\otimes D}\otimes\binom{B}{C}\to\binom{A}{C\otimes D}
\]
where $B,D,C,E$ are comonoids. Unfolding, we need two maps in $\cat{C}$:
\begin{equation}\label{eqn.wd_as_maps}
B\otimes D\otimes C\to C\otimes D
\qqand
(B\otimes D\otimes C)\otimes A\to C\otimes A\otimes B
\end{equation}
But these maps are obvious: they are just counit projections and braidings.

One could take $\cat{C}\coloneqq\mfd$. Then $A,\ldots, E$ are manifolds (all of which are naturally comonoids) and the wiring diagram represents information flow, shuffling coordinates around. But there are many more smooth functions of the form \eqref{eqn.wd_as_maps} than there are wiring diagrams, e.g.\ there is more than one smooth function $\rr\to\rr$.

One can also interpret \eqref{eqn.arb_wd} in $\finset\op$. Taking $A\cong B\cong C\cong D\cong E\cong 1$ for simplicity, each has a unique monoid structure. Now \eqref{eqn.wd_as_maps} consists of two maps
\[
\{b,d,c\}\from\{c,d\}
\qqand
\{b,d,c\}+\{a\}\from\{c,a,b\}.
\]
These two functions entirely determine the wiring diagram, saying how each outer outgoing port $c,d$ is ``fed by'' some inner outgoing port (in $\{b,d,c\}$), and how each inner ingoing port $c,a,b$ is ``fed by'' either an inner outgoing port (in $\{b,d,c\}$) or the outer ingoing port, $a$. 

\begin{lemma}\label{lemma.lens_rr}
The functor $\rr^-\colon\finset\op\to\mfd$ is strong monoidal, hence induces a strong monoidal functor
\[
\Lens{\rr^-}\colon\Lens{\finset\op}\to\Lens{\mfd}
\]
\end{lemma}

\subsection{Backward monads on lenses}\label{subsec.backwards}

We will add a notion of potentials to our lenses by proving a more general result: any strong monad on $\cat{C}$ induces a comonad on $\Lens{\cat{C}}$ by applying it to the input / backward-pass part. We pun on the word ``backwards'' to refer to the fact that it applies to the backward pass and the fact that it turns a monad into a comonad.

Recall that a monad $\Fun{T}$ on a monoidal category is \emph{strong} if there are natural maps $\sigma\colon X\otimes\Fun{T}Y\to\Fun{T}(X\otimes Y)$, satisfying four diagrams~\cite{kock1972strong}. If moreover $X$ carries a counit $\varepsilon_X\colon X\to I$, then naturality of $\sigma$ in its first argument together with the unit-coherence axiom of strength gives a \emph{counital projection} $\sigma_{X,Y}\then\Fun{T}(\varepsilon_X\otimes\id_Y)=\varepsilon_X\otimes\id_{\Fun{T}Y}$:
\begin{equation}\label{eqn.counital_projection}
\begin{tikzcd}
X\otimes\Fun{T}Y\ar[r, "\sigma"]\ar[d, "\varepsilon_X\otimes\id"']&
\Fun{T}(X\otimes Y)\ar[d, "\Fun{T}(\varepsilon_X\otimes\id)"]\\
\Fun{T}Y\ar[r, equal]&
\Fun{T}Y
\end{tikzcd}
\end{equation}
When $\cat{C}$ is cartesian, this specializes to the familiar projection formula $\sigma\then\Fun{T}\pi_Y=\pi_{\Fun{T}Y}$.


\begin{proposition}\label{prop.backward_comonad}
Let $(\Fun{T},\eta,\mu,\sigma)$ be a strong monad on a symmetric monoidal category $(\cat{C},I,\otimes)$. Then there is a comonad $\Lens{\Fun{T}}$ on $\Lens{\cat{C}}$ sending $\lensob{c}\mapsto \binom{\Fun{T}\inpt c}{\outp c}$. We denote its coKleisli category by $\Lens{\cat{C}}^{\Fun{T}}\coloneqq(\Lens{\cat{C}})_{\Lens{\Fun T}}$. 

If moreover $(\Fun{T},\tau)$ is a monoidal monad, then $\Lens{\Fun T}$ is a colax monoidal comonad, thus the coKleisli category $(\Lens{\cat{C}}^{\Fun{T}},I,\otimes)$ is naturally monoidal.
\end{proposition}
\begin{proof}
On a lens map $(\outp f,\inpt f)$ with $\inpt f\colon\outp c\otimes\inpt d\to\inpt c$, define $\Lens{\Fun T}(\outp f,\inpt f)$ to be the following composite:
\begin{equation}\label{eqn.lens_T_morphism}
\outp c\otimes\Fun{T}\inpt d\To{\sigma}\Fun{T}(\outp c\otimes\inpt d)\To{\Fun{T}\inpt f}\Fun{T}\inpt c.
\end{equation}
The counit and comultiplication of $\Lens{\Fun{T}}$ are given by the unit and multiplication of the monad, together with the counit $\varepsilon_{\outp c}$:
\begin{equation}\label{eqn.lens_T_counit}
\outp c\otimes\inpt c\To{\varepsilon_{\outp c}\otimes\id}\inpt c\To{\eta}\Fun{T}\inpt c,
\end{equation}
\begin{equation}\label{eqn.lens_T_comult}
\outp c\otimes\Fun{T}\Fun{T}\inpt c\To{\varepsilon_{\outp c}\otimes\id}\Fun{T}\Fun{T}\inpt c\To{\mu}\Fun{T}\inpt c.
\end{equation}
We need to prove that $\Lens{\Fun{T}}$ preserves identity and composition (is a functor) and that the functor is counital and coassociative. These four axioms correspond precisely to the four axioms of a monad strength.
%The unit axiom~\cite[(1.5)]{kock1972strong} (left unitor coherence) gives preservation of identity, via $\sigma_{X,Y}\then\Fun{T}\pi_Y=\pi_{\Fun{T}Y}$. The $\eta$-coherence axiom of the strong monad gives the two counit laws. The $\mu$-coherence axiom, together with associativity of $\mu$, gives coassociativity.

We include only the least trivial here: preservation of composition. Writing $f_o,g_o$ for forward components and $f_i\colon\outp c\otimes\inpt d\to\inpt c$, $g_i\colon\outp d\otimes\inpt e\to\inpt d$ for backward components, we must show that the following commutes (after appending $\Fun{T}f_i\colon\Fun{T}(\outp c\otimes\inpt d)\to\Fun{T}\inpt c$ on the right):
\[
\begin{tikzcd}[column sep=1.6em, row sep=3em, every label/.append style={font=\scriptsize}, font=\small, ampersand replacement=\&]
\outp c\otimes\Fun{T}\inpt e
  \ar[r, "\delta\otimes\id"]
  \ar[d, "\sigma"']
\& (\outp c)^{\otimes 2}\otimes\Fun{T}\inpt e
  \ar[r, "\id\otimes f_o\otimes\id"]
  \ar[d, "\sigma"]
\&[5pt] \outp c\otimes\outp d\otimes\Fun{T}\inpt e
  \ar[r, "\id\otimes\sigma"]
  \ar[dr, "\sigma" description]
\& \outp c\otimes\Fun{T}(\outp d\otimes\inpt e)
  \ar[r, "\id\otimes\Fun{T}g_i"]
  \ar[d, "\sigma"]
\& \outp c\otimes\Fun{T}\inpt d
  \ar[d, "\sigma"]
\\
\Fun{T}(\outp c\otimes\inpt e)
  \ar[r, "\Fun{T}(\delta\otimes\id)"']
\& \Fun{T}((\outp c)^{\otimes 2}\otimes\inpt e)
  \ar[rr, "\Fun{T}(\id\otimes f_o\otimes\id)"']
\&
\& \Fun{T}(\outp c\otimes\outp d\otimes\inpt e)
  \ar[r, "\Fun{T}(\id\otimes g_i)"']
\& \Fun{T}(\outp c\otimes\inpt d)
\end{tikzcd}
\]
The first two squares are naturality of $\sigma$ in its first argument; the middle triangle is the associator axiom of strength~\cite[(1.6)]{kock1972strong}, which merges the nested $\sigma$'s; the last square is naturality of $\sigma$ in its second argument.

When $(\Fun T,\tau_0,\tau)$ is moreover a monoidal monad, we obtain an colax monoidal structure on $\Lens{\Fun{T}}$: the counitor $\binom{I}{I}\to\binom{\Fun T(I)}{I}$ amounts to $\tau_0\colon I\to\Fun T I$ and the coproductor map $\binom{\Fun{T}(\inpt c\otimes\inpt d)}{\outp c\otimes\outp d}\to\binom{\Fun{T}(\inpt c)\otimes\Fun T(\inpt d)}{\outp c\otimes\outp d}$ amounts to $\tau\colon \Fun{T}(\inpt c)\otimes\Fun T(\inpt d)\to \Fun{T}(\inpt c\otimes\inpt d)$.
\end{proof}

A map in $\Lens{\cat{C}}^{\Fun{T}}$ is a map of the form $\binom{\Fun{T}\inpt c}{\outp c}\to\lensob{d}$ in $\Lens{\cat{C}}$, i.e.
\[
\outp f\colon \outp c\to\outp d
\qqand
\inpt f\colon \inpt d\otimes\outp c\to\Fun{T}\inpt c
\]
where $\outp f$ is a comonoid homomorphism.

\begin{example}[Monoidal monads on $\finset\op$]\label{ex.mms_finsetop}
Any finite set $S$ induces a functor
\[(S\times -)\colon\finset\op\to\finset\op,\]
which is strong monoidal by the usual distributive law
\[
S\times(B+C)\From{\cong}(S\times B)+(S\times C).
\]
It is a (comonad on $\finset$ and hence) a monad on $\finset\op$, easily checked to be monoidal with respect to $(0,+)$.%
\footnote{In fact the $S\times-$, for any $S:\finset$, are the only monoidal monads on $(\finset\op,0,+)$.
}

Thus it induces a comonad $\Lens{S\times-}$ on $\Lens{\finset\op}$. One may think of an element $s:S$ as a \emph{hyperparameter} that decides a sort of internal wiring. Indeed, the data of a morphism in the coKleisli category $\Lens{\finset\op}^{S\times-}$ is similar to that in \cref{eqn.lens_finsetop}, except that we have a function $\inpt f_s\colon\outp A+\inpt B\from\inpt A$ for each $s:S$.
\end{example}

\begin{lemma}\label{lem.lens_T_functoriality}
The construction of \cref{prop.backward_comonad} is functorial in $(\cat{C},\Fun{T})$: a strong symmetric monoidal $F\colon\cat C\to\cat C'$ and strong monad morphism $\theta\colon F\Fun T\Rightarrow\Fun T'F$ induce
\[
\Lens F\colon\Lens{\cat C}^{\Fun T}\to\Lens{\cat C'}^{\Fun T'},\qquad
\lensob c\mapsto\binom{F(\inpt c)}{F(\outp c)},\qquad
(\outp f,\inpt f)\mapsto(F(\outp f),\theta\circ F(\inpt f)\circ F_2),
\]
where $F_2$ is $F$'s productor. If $\Fun T,\Fun T'$ are monoidal monads and $\theta$ is a morphism of monoidal monads, $\Lens F$ is strong symmetric monoidal.
\end{lemma}

\begin{example}
For any set $S$ there is a monoidal monad $M\mapsto M^S$ on $\mfd$. In fact, this monad distributes under any other strong monad $T$,
\[
T(M^S)\to T(M)^S.
\]
The functor $\rr^-\colon\finset\op\to\mfd$ satisfies the conditions of \cref{lem.lens_T_functoriality}, so the functor $\Lens{\rr^-}$ from \cref{lemma.lens_rr} is strong monoidal. We will not pursue this example further in the paper.
\end{example}

\begin{lemma}\label{lem.parameter_lens_action}
Let $\cat V$ be symmetric monoidal, let $\cat C$ be symmetric monoidal with a homomorphic supply of commutative comonoids, let $\Fun T$ be a strong monad on $\cat C$, and let $J\colon\cat V\to\cat C$ be strong monoidal. Then we have an action of $\cat V$ on $\Lens{\cat C}^{\Fun T}$ given by
\[
v\cdot\lensob c\coloneqq\binom{\inpt c}{J(v)\otimes\outp c}.
\]
If moreover $\Fun T$ is a monoidal monad, this action is symmetric monoidal.
\end{lemma}

\begin{proof}
Use the supplied commutative comonoid on $J(v)\in\cat C$; by homomorphicity, $J(u)$ is a comonoid homomorphism for any $u\colon v\to w$ in $\cat V$. Given $f\colon\lensob c\to\lensob d$ in $\Lens{\cat C}^{\Fun T}$ and $u$ as above, set
\[
u\cdot f\coloneqq\binom{\inpt f\circ\epsilon_{J(v)}}{J(u)\otimes\outp f}\colon
\binom{\inpt c}{J(v)\otimes\outp c}
\to
\binom{\inpt d}{J(w)\otimes\outp d}.
\]
Checking functoriality reduces immediately to the counital projection \eqref{eqn.counital_projection} and the homomorphic supply assumption: $J(v)\To{J(u)}J(w)\To{\epsilon_{J(w)}} I$ equals $\epsilon_{J(v)}$.

When $\Fun T$ is monoidal, $\Lens{\cat C}^{\Fun T}$ is symmetric monoidal by \cref{prop.backward_comonad}, and the comparisons $J(I)\cong I$ and $(v\otimes w)\cdot(\lensob c\otimes\lensob d)\cong(v\cdot\lensob c)\otimes(w\cdot\lensob d)$ come directly from the strong monoidality of $J$.
\end{proof}

\section{Internal representation of lenses}\label{sec.abstract_representation}

The previous section gives the category of lenses, and discusses the comonad and coKleisli category associated to a monad on the base category. The following proposition says that we can interpret the category of lenses inside the base category.

\begin{proposition}\label{prop.Theta}
For any symmetric monoidal closed category $(\cat{C},I,\otimes, [-,-])$, there is a normal lax monoidal functor
\[\Theta\colon\Lens{\cat{C}}\to\cat{C},\qquad
\lensob{c}\mapsto\ihom{\inpt c,I}\otimes\outp c.\]
\end{proposition}
\begin{proof}
This is an immediate corollary of \cref{prop.Theta_T_alpha} with $\Fun{T}=\id$. 
%We need to prove the equality $\Theta(g\circ f)=\Theta(g)\circ\Theta(f)\colon\outp c\,H_c\to\outp e\,H_e$ for $f\colon\lensob c\to\lensob d$ and $g\colon\lensob d\to\lensob e$. It amounts to the commutativity of the outer square below, where we write $H_x\coloneqq[\inpt x,I]$, and for space reasons write $c,d,e$ for $\inpt c,\inpt d,\inpt e$ and elide the $\otimes$-symbol, so e.g.\ $\outp d\outp d[d,c H_c]$ means $\outp d\otimes\outp d\otimes[\inpt d,\inpt c\otimes [c,I]]$. 
%\[
%\begin{tikzcd}[column sep=1em, row sep=1.4em, every label/.append style={font=\tiny}, font=\scriptsize, ampersand replacement=\&]
%\outp c H_c
%  \ar[r, "{\delta}"]
%  \ar[d, "{\delta}"']
%\& \outp c\outp c H_c \ar[r, "{\outp f}"] \ar[d, "{\delta}"']
%\& \outp d\outp c H_c \ar[r, "{\coev_d}"] \ar[d, "{\delta}"']
%\& \outp d[d,d\outp c H_c] \ar[r, "{f}"] \ar[d, "{\delta}"']
%\& \outp d[d,c H_c] \ar[r, "{\ev}"] \ar[d, "{\delta}"']
%\& \outp d H_d \ar[d, "{\delta}"] \\
%\outp c\outp c H_c \ar[d, "{\outp f}"'] \ar[r, "{\delta}"]
%\& \outp c\outp c\outp c H_c \ar[r, "{\outp f\outp f}"] \ar[d, "{\outp f}"']
%\& \outp d\outp d\outp c H_c \ar[r, "{\coev_d}"] \ar[d, "{\outp g}"']
%\& \outp d\outp d[d,d\outp c H_c] \ar[r, "{f}"] \ar[d, "{\outp g}"']
%\& \outp d\outp d[d,c H_c] \ar[r, "{\ev}"] \ar[d, "{\outp g}"']
%\& \outp d\outp d H_d \ar[d, "{\outp g}"] \\
%\outp d\outp c H_c \ar[d, "{\outp g}"'] \ar[r, "{\delta}"]
%\& \outp d\outp c\outp c H_c \ar[r, "{\outp g\outp f}"]
%\& \outp e\outp d\outp c H_c \ar[r, "{\coev_d}"] \ar[d, "{\coev_e}"']
%\& \outp e\outp d[d,d\outp c H_c] \ar[r, "{f}"] \ar[d, "{\coev_e}"']
%\& \outp e\outp d[d,c H_c] \ar[r, "{\ev}"] \ar[d, "{\coev_e}"']
%\& \outp e\outp d H_d \ar[d, "{\coev_e}"] \\
%\outp e\outp c H_c \ar[d, "{\coev_e}"']
%\& \outp d[e,e\outp cH_c]\ar[from=ul, "{\coev_e}"]\ar[dl, "\outp g"]\ar[d,"\delta"]
%\& \outp e[e,e\outp d\outp c H_c] \ar[r, "{\coev_d}"] \ar[d, "{g}"'] \ar[dd, equal, bend right=60pt]
%\& \outp e[e,e\outp d[d,d\outp c H_c]] \ar[r, "{f}"] \ar[d, "{g}"']
%\& \outp e[e,e\outp d[d,c H_c]] \ar[r, "{\ev}"] \ar[d, "{g}"']
%\& \outp e[e,e\outp d H_d] \ar[d, "{g}"] \\
%\outp e[e,e\outp c H_c] \ar[d, "{\delta}"']
%\& \outp d[e,e\outp c\outp cH_c]\ar[ld, "{\outp g}"]\ar[ur, "\outp g\outp f"', pos=.3]\ar[from=uu, , bend left=60pt, "{\coev_e}", pos=.3]
%\& \outp e[e,d\outp c H_c] \ar[r, "{\coev_d}"] \ar[dr, equal]
%\& \outp e[e,d[d,d\outp c H_c]] \ar[r, "{f}"] \ar[d, "{\ev}"']
%\& \outp e[e,d[d,c H_c]] \ar[r, "{\ev}"] \ar[d, "{\ev}"']
%\& \outp e[e,d H_d] \ar[d, "{\ev}"] \\
%\outp e[e,e\outp c\outp c H_c] \ar[rr, "{\outp f}"]
%\&\& \outp e[e,e\outp d\outp c H_c] \ar[r, "{g}"]
%\& \outp e[e,d\outp c H_c] \ar[r, "{f}"]
%\& \outp e[e,c H_c] \ar[r, "{\ev}"]
%\& \outp e H_e
%\end{tikzcd}
%\]
\end{proof}


Recall that monoidal (equivalently, commutative) monads $(\Fun{T},\tau)$ have a natural strength
\[
\sigma_{X,Y}\colon X\otimes\Fun{T}Y\To{\eta_X\otimes\id}\Fun{T}X\otimes\Fun{T}Y\To{\tau_{X,Y}}\Fun{T}(X\otimes Y).
\]
and that monoidal monads $\Fun{T}$ preserve monoids.

\begin{definition}\label{def.T_monoid}
Let $(\Fun{T},\eta,\mu,\tau)$ be a monoidal monad on a symmetric monoidal category $\cat{C}$. A \emph{$\Fun{T}$-monoid in $\cat{C}$} is a $\otimes$-monoid $(z,e_z,m_z)$ in $\cat{C}$
\[
\begin{tikzcd}[ampersand replacement=\&]
z\ar[r,"\eta_z"]\ar[dr,equal]\& \Fun{T}z\ar[d,"\alpha"]\\
\& z
\end{tikzcd}
\hspace{1in}
\begin{tikzcd}[ampersand replacement=\&]
\Fun{T}\Fun{T}z\ar[r,"\mu_z"]\ar[d,"\Fun{T}\alpha"']\& \Fun{T}z\ar[d,"\alpha"]\\
\Fun{T}z\ar[r,"\alpha"']\& z
\end{tikzcd}
\]
together with a $\Fun{T}$-algebra structure $\alpha\colon\Fun{T}z\to z$ such that $\alpha$ is a monoid homomorphism
\begin{equation}\label{eqn.T_monoid}
\begin{tikzcd}[ampersand replacement=\&]
I\ar[r,"\tau_0"]\ar[d,"e_z"']\& \Fun{T}I\ar[r,"\Fun{T}e_z"]\& \Fun{T}z\ar[d,"\alpha"]\\
z\ar[r,"\eta_z"']\& \Fun{T}z\ar[r,"\alpha"']\& z
\end{tikzcd}
\hspace{1in}
\begin{tikzcd}[ampersand replacement=\&]
\Fun{T}z\otimes\Fun{T}z\ar[r,"\tau_{z,z}"]\ar[d,"\alpha\otimes\alpha"']\& \Fun{T}(z\otimes z)\ar[r,"\Fun{T}m_z"]\& \Fun{T}z\ar[d,"\alpha"]\\
z\otimes z\ar[rr,"m_z"']\&\& z
\end{tikzcd}
\end{equation}
We call $z$ the \emph{value object}.
\end{definition}

\begin{remark}\label{rmk.T_monoid_EM}
A $\Fun{T}$-monoid in $\cat{C}$ is the same data as a monoid in the Eilenberg--Moore category $\cat{C}^{\Fun{T}}$, equivalently a $\Fun{T}$-algebra in $\Cat{Mon}(\cat{C})$. The two presentations agree because $\Fun{T}$ being monoidal is exactly the condition that the underlying-functor of $\Fun{T}$ lifts to a monad on $\Cat{Mon}(\cat{C})$; see \cite{kock1972strong}.
\end{remark}

\begin{proposition}\label{prop.Theta_T_alpha}
Let $(\Fun{T},\tau)$ be a monoidal monad on a symmetric monoidal closed category $(\cat{C},I,\otimes, [-,-])$, and let $(z,e,m,\alpha)$ be a $\Fun{T}$-monoid in $\cat{C}$ (\cref{def.T_monoid}). Then there is a lax monoidal functor that acts on objects as follows:
\[
\Theta_z\colon\Lens{\cat{C}}^\Fun{T}\to\cat{C},\qquad
\lensob{c}\mapsto\ihom{\inpt c,z}\otimes\outp c.
\]
\end{proposition}
\begin{proof}
Given a coKleisli morphism $f\colon\lensob{c}\to\lensob{d}$, define $\Theta_z(f)$ as the composite
\[
\begin{tikzcd}[column sep=45pt, row sep=22pt, ampersand replacement=\&]
[\inpt c,z]\otimes\outp c\ar[r,"{\id\otimes\delta_{\outp c}}"]\ar[dd, dashed, "\Theta_z(f)"']
\& {[\inpt c,z]\otimes\outp c\otimes\outp c}\ar[r,"\coev\otimes\outp f"]
\& {[\inpt d,\inpt d\otimes\outp c\otimes[\inpt c,z]]\otimes\outp d}\ar[d,"{[\inpt d,\inpt f\otimes\id]\otimes\id}"]
\\
\&
\& {[\inpt d,\Fun{T}\inpt c\otimes[\inpt c,z]]\otimes\outp d}\ar[d,"{[\inpt d,\sigma]\otimes\id}"]
\\
{[\inpt d,z]\otimes\outp d}
\& {[\inpt d,\Fun{T}z]\otimes\outp d}\ar[l,"{[\inpt d,\alpha]\otimes\id}"]
\& {[\inpt d,\Fun{T}(\inpt c\otimes[\inpt c,z])]\otimes\outp d}\ar[l,"{[\inpt d,\Fun{T}\ev]\otimes\id}"]
\end{tikzcd}
\]
Proving that $\Theta_z$ is functorial a long calculation but routine. We leave preservation of identities to the reader and
%IDENTITIES
%\[
%\begin{tikzcd}[column sep=2.6em, row sep=1.85em, every label/.append style={font=\scriptsize}, font=\small, ampersand replacement=\&]
%|[alias=tl]| \outp c\otimes H
%  \ar[rr, "{\delta}"]
%  \ar[ddd, equal]
%  \ar[drr, equal]
%\&
%\& \outp c\otimes\outp c\otimes H
%  \ar[r, "{\mathbf{coev}}", ""{description, name=A1, anchor=center, inner sep=0}]
%  \ar[d, "{\varepsilon}" description]
%\& \outp c\otimes[\inpt c,\inpt c\otimes\outp c\otimes H]
%  \ar[d, "{\varepsilon}"]
%\\
%\&
%\& |[alias=ce]| \outp c\otimes H
%  \ar[r, "{\mathbf{coev}}", ""{description, name=A2, anchor=center, inner sep=0}]
%  \ar[d, "{\eta_H}"']
%  \ar[dl, "{\mathbf{coev}}"' description]
%  \ar[ddll, equal, out=180, in=60]
%\& \outp c\otimes[\inpt c,\inpt c\otimes H]
%  \ar[d, "{\eta_{\inpt c}}"]
%  \ar[ddl, "{\eta_H}"' description]
%\\
%\& \outp c\otimes[\inpt c,\inpt c\otimes H]
%  \ar[dl, "{\ev}"]
%  \ar[dr, "{\eta}"]
%  \ar[ddr, "{\eta_{\inpt c}}"' description]
%  \ar[dddr, bend right=20, "{\eta_{\inpt c\otimes H}}"' description]
%\& \outp c\otimes\Fun{T}H
%  \ar[d, "{\mathbf{coev}}"]
%\& \outp c\otimes[\inpt c,\Fun{T}\inpt c\otimes H]
%  \ar[d, "{\eta_H}", ""{description, name=DR, anchor=center, inner sep=0}]
%\\
%\outp c\otimes[\inpt c,z]
%  \ar[ddr, "{\eta_z}"]\ar[ddd, equal]
%\&\& \outp c\otimes[\inpt c,\inpt c\otimes\Fun{T}H]
%  \ar[r, "{\sigma^{L}}", ""{description, name=DL, anchor=center, inner sep=0}]
%\& \outp c\otimes[\inpt c,\Fun{T}\inpt c\otimes\Fun{T}H]
%  \ar[ddl, "{\tau}"]
%\\
%\&
%\& \outp c\otimes[\inpt c,\Fun{T}\inpt c\otimes H]
%  \ar[ur, "{\eta_H}"]
%\&
%\\
%\& |[alias=mz]| \outp c\otimes[\inpt c,\Fun{T}z]
%  \ar[dl, "{\alpha}"' description]
%\& \outp c\otimes[\inpt c,\Fun{T}(\inpt c\otimes H)]
%  \ar[l, "{\Fun{T}\ev}"' description]
%\&
%\\
%\outp c\otimes H
%\&
%\&
%\&
%\end{tikzcd}
%\]
%Each square is clear.
show composition on the following page. We need to prove that two morphisms $\outp c\otimes H_c\to\outp e\otimes H_e$ (writing $H_X\coloneqq[\inpt X,z]$) are equal: $\Theta_z$ of the coKleisli composite $g\circ f$ and the coKleisli composite $\Theta_z(g)\circ\Theta_z(f)$ of the $\Theta_z$'s. 

On the following page, each ``tile'' is of one of three types: functoriality of $\otimes$; naturality of $\eta$, $\mu$, and $\tau$; and the $\Fun T$-monoid axioms of \cref{def.T_monoid}. 

The unitor is $e_z\colon I\to z=\Theta_z\binom{I}{I}$, and the productor is the canonical map
\[
([\inpt c,z]\otimes\outp c)\otimes([\inpt d,z]\otimes\outp d)
\to
[\inpt c\otimes\inpt d,z\otimes z]\otimes\outp c\otimes\outp d\To{[\id,m_z]\otimes\id}
[\inpt c\otimes\inpt d,z]\otimes\outp c\otimes\outp d
\]
Its naturality, unitality, and associativity amount to the monoidality of $\Fun{T}$, the assumption that $\alpha$ is a monoid homomorphism, and that $z$ is a monoid.


\begin{landscape}
For space reasons we use juxtaposition to denote $\otimes$, write $H_X\coloneqq[\inpt X,z]$, and write $x$ for $\inpt{x}$. So for example, $[d,\Fun T(c)H_c]\outp d$ means $[\inpt d,\Fun T(c)\otimes H_c]\otimes\outp d$. The strength $\sigma$ has been opened as $\eta$ then $\tau$.
\bigskip
\thispagestyle{empty}
\[\hspace{-.5in}
\begin{tikzcd}[column sep=0.5em, row sep=0.6em, every label/.append style={font=\tiny}, font=\tiny, ampersand replacement=\&]
H_c\outp c
  \ar[r, "{\delta}"]
  \ar[d, "{\delta}"']
\& H_c\outp c\outp c \ar[r, "{\outp f}"] \ar[d, "{\delta}"']
\& H_c\outp c\outp d \ar[r, "{\coev_d}"] \ar[d, "{\delta}"']
\& {[d,d\outp c H_c]}\outp d \ar[r, "{\inpt f}"] \ar[d, "{\delta}"']
\& {[d,\Fun T(c)H_c]}\outp d \ar[r, "{\eta}"] \ar[d, "{\delta}"']
\& {[d,\Fun T(c)\Fun T(H_c)]}\outp d \ar[r, "{\tau}"] \ar[d, "{\delta}"']
\& {[d,\Fun T(cH_c)]}\outp d \ar[r, "{\Fun T(\ev)}"] \ar[d, "{\delta}"']
\& {[d,\Fun T(z)]}\outp d \ar[r, "{\alpha}"] \ar[d, "{\delta}"']
\& H_d\outp d \ar[d, "{\delta}"] \\
H_c\outp c\outp c \ar[d, "{\outp f}"'] \ar[r, "{\delta}"]
\& H_c\outp c\outp c\outp c \ar[r, "{\outp f\outp f}"] \ar[d, "{\outp f}"']
\& H_c\outp c\outp d\outp d \ar[r, "{\coev_d}"] \ar[d, "{\outp g}"']
\& {[d,d\outp c H_c]}\outp d\outp d \ar[r, "{\inpt f}"] \ar[d, "{\outp g}"']
\& {[d,\Fun T(c)H_c]}\outp d\outp d \ar[r, "{\eta}"] \ar[d, "{\outp g}"']
\& {[d,\Fun T(c)\Fun T(H_c)]}\outp d\outp d \ar[r, "{\tau}"] \ar[d, "{\outp g}"']
\& {[d,\Fun T(cH_c)]}\outp d\outp d \ar[r, "{\Fun T(\ev)}"] \ar[d, "{\outp g}"']
\& {[d,\Fun T(z)]}\outp d\outp d \ar[r, "{\alpha}"] \ar[d, "{\outp g}"']
\& H_d\outp d\outp d \ar[d, "{\outp g}"] \\
H_c\outp c\outp d \ar[d, "{\outp g}"'] \ar[r, "{\delta}"]
\& H_c\outp c\outp c\outp d \ar[r, "{\outp g\outp f}"]
\& H_c\outp c\outp d\outp e \ar[r, "{\coev_d}"] \ar[d, "{\coev_e}"']
\& {[d,d\outp c H_c]}\outp d\outp e \ar[r, "{\inpt f}"] \ar[d, "{\coev_e}"']
\& {[d,\Fun T(c)H_c]}\outp d\outp e \ar[r, "{\eta}"] \ar[d, "{\coev_e}"']
\& {[d,\Fun T(c)\Fun T(H_c)]}\outp d\outp e \ar[r, "{\tau}"] \ar[d, "{\coev_e}"']
\& {[d,\Fun T(cH_c)]}\outp d\outp e \ar[r, "{\Fun T(\ev)}"] \ar[d, "{\coev_e}"']
\& {[d,\Fun T(z)]}\outp d\outp e \ar[r, "{\alpha}"] \ar[d, "{\coev_e}"']
\& H_d\outp d\outp e \ar[d, "{\coev_e}"] \\
H_c\outp c\outp e \ar[dd, "{\coev_e}"']
\&\& {[e,e\outp d\outp c H_c]}\outp e \ar[r, "{\coev_d}"] \ar[d, "{\inpt g}"']
\& {[e,e\outp d[d,d\outp c H_c]]}\outp e \ar[r, "{\inpt f}"] \ar[d, "{\inpt g}"']
\& {[e,e\outp d[d,\Fun T(c)H_c]]}\outp e \ar[r, "{\eta}"] \ar[d, "{\inpt g}"']
\& {[e,e\outp d[d,\Fun T(c)\Fun T(H_c)]]}\outp e \ar[r, "{\tau}"] \ar[d, "{\inpt g}"']
\& {[e,e\outp d[d,\Fun T(cH_c)]]}\outp e \ar[r, "{\Fun T(\ev)}"] \ar[d, "{\inpt g}"']
\& {[e,e\outp d[d,\Fun T(z)]]}\outp e \ar[r, "{\alpha}"] \ar[d, "{\inpt g}"']
\& {[e,e\outp d H_d]}\outp e \ar[d, "{\inpt g}"] \\
\& {[e,\Fun T(d)\Fun T(\outp c)H_c]}\outp e \ar[d, "{\eta}"'] \ar[dddd, out=200, in=150, equal]
\& {[e,\Fun T(d)\outp c H_c]}\outp e \ar[r, "{\coev_d}"] \ar[d, "{\eta}"] \ar[l, "{\eta}"]
\& {[e,\Fun T(d)[d,d\outp c H_c]]}\outp e \ar[r, "{\inpt f}"] \ar[d, "{\eta}"']
\& {[e,\Fun T(d)[d,\Fun T(c)H_c]]}\outp e \ar[r, "{\eta}"] \ar[d, "{\eta}"']
\& {[e,\Fun T(d)[d,\Fun T(c)\Fun T(H_c)]]}\outp e \ar[r, "{\tau}"] \ar[d, "{\eta}"']
\& {[e,\Fun T(d)[d,\Fun T(cH_c)]]}\outp e \ar[r, "{\Fun T(\ev)}"] \ar[d, "{\eta}"']
\& {[e,\Fun T(d)[d,\Fun T(z)]]}\outp e \ar[r, "{\alpha}"] \ar[d, "{\eta}"']
\& {[e,\Fun T(d)H_d]}\outp e \ar[d, "{\eta}"] \\
{[e,e\outp c H_c]}\outp e \ar[d, "{\delta}"']
\& {[e,\Fun T(d)\Fun T(\outp c)\Fun T(H_c)]}\outp e \ar[r, "{\tau}"]
\& {[e,\Fun T(d)\Fun T(\outp c H_c)]}\outp e \ar[r, "{\Fun T(\coev_d)}"] \ar[d, "{\tau}"']
\& {[e,\Fun T(d)\Fun T([d,d\outp c H_c])]}\outp e \ar[r, "{\Fun T(\inpt f)}"] \ar[d, "{\tau}"']
\& {[e,\Fun T(d)\Fun T([d,\Fun T(c)H_c])]}\outp e \ar[r, "{\Fun T(\eta)}"] \ar[d, "{\tau}"']
\& {[e,\Fun T(d)\Fun T([d,\Fun T(c)\Fun T(H_c)])]}\outp e \ar[r, "{\Fun T(\tau)}"] \ar[d, "{\tau}"']
\& {[e,\Fun T(d)\Fun T([d,\Fun T(cH_c)])]}\outp e \ar[r, "{\Fun T(\ev)}"] \ar[d, "{\tau}"']
\& {[e,\Fun T(d)\Fun T([d,\Fun T(z)])]}\outp e \ar[r, "{\Fun T(\alpha)}"] \ar[d, "{\tau}"']
\& {[e,\Fun T(d)\Fun T(H_d)]}\outp e \ar[d, "{\tau}"] \\
{[e,e\outp c\outp c H_c]}\outp e \ar[d, "{\outp f}"']
\&\& {[e,\Fun T(d\outp c H_c)]}\outp e \ar[r, "{\Fun T(\coev_d)}"] \ar[dr, equal]
\& {[e,\Fun T(d[d,d\outp c H_c])]}\outp e \ar[r, "{\Fun T(\inpt f)}"] \ar[d, "{\Fun T(\ev)}"']
\& {[e,\Fun T(d[d,\Fun T(c)H_c])]}\outp e \ar[r, "{\Fun T(\eta)}"] \ar[d, "{\Fun T(\ev)}"']
\& {[e,\Fun T(d[d,\Fun T(c)\Fun T(H_c)])]}\outp e \ar[r, "{\Fun T(\tau)}"] \ar[d, "{\Fun T(\ev)}"']
\& {[e,\Fun T(d[d,\Fun T(cH_c)])]}\outp e \ar[r, "{\Fun T(\ev)}"] \ar[d, "{\Fun T(\ev)}"']
\& {[e,\Fun T(d[d,\Fun T(z)])]}\outp e \ar[r, "{\Fun T(\alpha)}"] \ar[d, "{\Fun T(\ev)}"']
\& {[e,\Fun T(dH_d)]}\outp e \ar[d, "{\Fun T(\ev)}"] \\
{[e,e\outp d\outp c H_c]}\outp e \ar[d, "{\inpt g}"']
\&\&
\& {[e,\Fun T(d\outp c H_c)]}\outp e \ar[r, "{\Fun T(\inpt f)}"]
\& {[e,\Fun T(\Fun T(c)H_c)]}\outp e \ar[r, "{\Fun T(\eta)}"]
\& {[e,\Fun T(\Fun T(c)\Fun T(H_c))]}\outp e \ar[r, "{\Fun T(\tau)}"]
\& {[e,\Fun T(\Fun T(cH_c))]}\outp e \ar[r, "{\Fun T(\ev)}"] \ar[d, "{\mu}"']
\& {[e,\Fun T(\Fun T(z))]}\outp e \ar[r, "{\Fun T(\alpha)}"] \ar[d, "{\mu}"']
\& {[e,\Fun T(z)]}\outp e \ar[d, "{\alpha}"] \\
{[e,\Fun T(d)\outp c H_c]}\outp e \ar[r, "{\eta}"]
\& {[e,\Fun T(d)\Fun T(\outp c)H_c]}\outp e \ar[r, "{\tau}"]
\& {[e,\Fun T(d\outp c)H_c]}\outp e \ar[r, "{\Fun T(\inpt f)}"]
\& {[e,\Fun T(\Fun T(c))H_c]}\outp e \ar[r, "{\mu}"]
\& {[e,\Fun T(c)H_c]}\outp e \ar[r, "{\eta}"]
\& {[e,\Fun T(c)\Fun T(H_c)]}\outp e \ar[r, "{\tau}"]
\& {[e,\Fun T(cH_c)]}\outp e \ar[r, "{\Fun T(\ev)}"]
\& {[e,\Fun T(z)]}\outp e \ar[r, "{\alpha}"]
\& H_e\outp e
\end{tikzcd}
\]

\medskip

The open hole splits into two regions, which commute as shown below.

\noindent\begin{minipage}{0.48\linewidth}
\[
\begin{tikzcd}[column sep=1em, row sep=0.9em, every label/.append style={font=\tiny}, font=\scriptsize, ampersand replacement=\&]
H_c\outp c\outp d
  \ar[r, "{\delta}"]
  \ar[d, "{\outp g}"']
\& H_c\outp c\outp c\outp d \ar[r, "{\outp g}"] \ar[d, "{\outp g}"']
\& H_c\outp c\outp c\outp e \ar[r, "{\outp f}"] \ar[d, "{\outp f}"']
\& H_c\outp c\outp d\outp e \ar[d, "{\coev_e}"] \\
H_c\outp c\outp e \ar[r, "{\delta}"] \ar[d, "{\coev_e}"']
\& H_c\outp c\outp c\outp e \ar[r, "{\outp f}"] \ar[d, "{\coev_e}"']
\& H_c\outp c\outp d\outp e \ar[r, "{\coev_e}"] \ar[d, "{\coev_e}"']
\& {[e,e\outp d\outp c H_c]}\outp e \ar[d, "{\inpt g}"] \\
{[e,e\outp c H_c]}\outp e \ar[r, "{\delta}"] \ar[d, "{\delta}"']
\& {[e,e\outp c\outp c H_c]}\outp e \ar[r, "{\outp f}"] \ar[d, "{\outp f}"']
\& {[e,e\outp d\outp c H_c]}\outp e \ar[r, "{\inpt g}"] \ar[d, "{\inpt g}"']
\& {[e,\Fun T(d)\outp c H_c]}\outp e \ar[d, "{\eta}"] \\
{[e,e\outp c\outp c H_c]}\outp e \ar[r, "{\outp f}"]
\& {[e,e\outp d\outp c H_c]}\outp e \ar[r, "{\inpt g}"]
\& {[e,\Fun T(d)\outp c H_c]}\outp e \ar[r, "{\eta}"]
\& {[e,\Fun T(d)\Fun T(\outp c)H_c]}\outp e
\end{tikzcd}
\]
\end{minipage}\hfill
\begin{minipage}{0.5\linewidth}
\[
\begin{tikzcd}[column sep=1em, row sep=0.9em, every label/.append style={font=\tiny}, font=\scriptsize, ampersand replacement=\&]
{[e,\Fun T(d)\Fun T(\outp c)H_c]}\outp e
  \ar[r, "{\eta}"]
  \ar[d, "{\tau}"']
\& {[e,\Fun T(d)\Fun T(\outp c)\Fun T(H_c)]}\outp e \ar[r, "{\tau}"] \ar[d, "{\tau}"']
\& {[e,\Fun T(d)\Fun T(\outp c H_c)]}\outp e \ar[r, "{\tau}"]
\& {[e,\Fun T(d\outp c H_c)]}\outp e \ar[d, "{\Fun T(\inpt f)}"] \\
{[e,\Fun T(d\outp c)H_c]}\outp e \ar[r, "{\eta}"] \ar[dd, "{\Fun T(\inpt f)}"']
\& {[e,\Fun T(d\outp c)\Fun T(H_c)]}\outp e \ar[urr, "{\tau}"] \ar[d, "{\Fun T(\inpt f)}"']
\&\& {[e,\Fun T(\Fun T(c)H_c)]}\outp e \ar[d, "{\Fun T(\eta)}"] \\
\& {[e,\Fun T(\Fun T(c))\Fun T(H_c)]}\outp e \ar[urr, "{\tau}"] \ar[dd, "{\mu}"'] \ar[r, "{\Fun T(\eta)}"']
\& {[e,\Fun T(\Fun T(c))\Fun T(\Fun T(H_c))]}\outp e \ar[r, "{\tau}"] \ar[d, "{\mu}"']
\& {[e,\Fun T(\Fun T(c)\Fun T(H_c))]}\outp e \ar[d, "{\Fun T(\tau)}"] \\
{[e,\Fun T(\Fun T(c))H_c]}\outp e \ar[ur, "{\eta}"] \ar[d, "{\mu}"']
\&\& {[e,\Fun T(c)\Fun T(\Fun T(H_c))]}\outp e \ar[dl, "{\mu}"]
\& {[e,\Fun T(\Fun T(cH_c))]}\outp e \ar[d, "{\mu}"] \\
{[e,\Fun T(c)H_c]}\outp e \ar[r, "{\eta}"]
\& {[e,\Fun T(c)\Fun T(H_c)]}\outp e \ar[rr, "{\tau}"]
\&\& {[e,\Fun T(cH_c)]}\outp e
\end{tikzcd}
\qedhere\]
\end{minipage}
\end{landscape}
\end{proof}

\begin{corollary}\label{cor.parameterized_representation}
Let $(\Fun T,\tau)$, $\cat C$, $(z,e,m,\alpha)$, and the resulting $\Theta_z\colon\Lens{\cat{C}}^\Fun{T}\to\cat{C}$ be as in \cref{prop.Theta_T_alpha}, with $\cat C$ carrying a homomorphic supply of commutative comonoids. Suppose $\cat V$ is symmetric monoidal and $J\colon\cat V\to\cat C$ is strong monoidal. Then with the $\cat V$-actions
\[
v\cdot\lensob c\coloneqq\binom{\inpt c}{J(v)\otimes\outp c}
\qqand
v\cdot x\coloneqq J(v)\otimes x
\]
on $\Lens{\cat C}^{\Fun T}$ and on $\cat C$ respectively, $\Theta_z$ extends to a lax symmetric monoidal functor
\begin{equation}\label{eqn.para_theta_z}
\para{\cat V}{\Theta_z}\colon\para{\cat V}{\Lens{\cat C}^{\Fun T}}\to\para{\cat V}{\cat C}.
\end{equation}
\end{corollary}
%It sends a parameterized coKleisli lens $(v,f\colon v\cdot\lensob c\to\lensob d)$ to the parameterized $\cat C$-morphism
%\[
%J(v)\otimes\outp{c}\otimes[\inpt{c},z]\to\outp{d}\otimes[\inpt{d},z].
%\]

\begin{proof}
Use the action of $\cat{V}$ on $\Lens{\cat C}^{\Fun T}$ from \cref{lem.parameter_lens_action} and the action $v\cdot x\coloneqq J(v)\otimes x$ on $\cat C$. The functor $\Theta_z$ is $\cat V$-equivariant:
\[
\begin{tikzcd}[ampersand replacement=\&]
\cat V\times\Lens{\cat C}^{\Fun T}
  \ar[r,"(\cdot)"]
  \ar[d,"\id_{\cat V}\times\Theta_z"']
  \ar[dr,phantom,"\cong"]
\& \Lens{\cat C}^{\Fun T}\ar[d,"\Theta_z"]\\
\cat V\times\cat C
  \ar[r,"(\cdot)"']
\& \cat C,
\end{tikzcd}
\]
where the isomorphism is
\[
\theta_{v,\lensob c}\colon
J(v)\otimes\Theta_z\lensob c
\cong J(v)\otimes\ihom{\inpt c,z}\otimes\outp c
=\Theta_z\left(v\cdot\lensob c\right).
\]
We can thus apply \cref{prop.para_square,prop.Theta_T_alpha} to obtain \eqref{eqn.para_theta_z} as desired.
\end{proof}

\section{The dynamics of potentialized lenses}\label{sec.potential_lenses_to_dynamics}

\subsection{Potentialized lenses}\label{subsec.potentialized_lenses}

\begin{lemma}\label{lemma.monoid_to_monad}
If $\cat{C}$ is a symmetric monoidal category and $(M,e,(*))$ is a commutative monoid object in it, then $M\otimes\blank\colon\cat{C}\to\cat{C}$ carries a monoidal monad structure, functorially in $(\cat C,M)$.
\end{lemma}

\begin{proof}
The monad unit and multiplication are $e\otimes\id$ and $(*)\otimes\id$. The monoidal structure has unitor $e\colon I\to M\cong M\otimes I$ and productor
\begin{equation}\label{eqn.monoid_productor}
(M\otimes X)\otimes(M\otimes Y)\cong M\otimes M\otimes X\otimes Y\To{(*)\otimes\id}M\otimes X\otimes Y.
\end{equation}
The coherence diagrams arise from the monoid laws and commutativity of $M$. Functoriality in $(\cat C,M)$ is straightforward.
\end{proof}

Since $(\rr,0,+)$ is a commutative Lie group, \cref{prop.cot_hopf} gives a commutative monoid---in fact Hopf monoid---object $\cot{\rr}$ in $\poly$. By \cref{lemma.monoid_to_monad} there are corresponding monoidal monads $\rr\times\blank$ and $\cot{\rr}\otimes\blank$ on $\mfd$ and $\poly$, and by \cref{prop.backward_comonad} there are corresponding colax monoidal comonads on $\Lens{\mfd}$ and $\Lens{\poly}$. We let $(\Lens{\mfd}^\rr,\rr^0,\times)$ and $(\Lens{\poly}^\rr,\yon,\otimes)$ denote the corresponding cokleisli categories.

A map in $\lmfd^\rr$ is a tuple $f=(\outp f,\inpt f,U)\colon\lensob M\to\lensob N$ with
\[
\outp f\colon\outp M\to\outp N,\qquad
\inpt f\colon\outp M\times\inpt N\to\inpt M,\qquad
U\colon\outp M\times\inpt N\to\rr.
\]
We call such maps \emph{potentialized manifold lenses}.
A map $\lensob p\to\lensob q$ in $\Lens{\poly}^\rr$ is a pair
\[
\outp f\colon\outp p\to\outp q,\qquad
\inpt f\colon\outp p\otimes\inpt q\to\cot{\rr}\otimes\inpt p.
\]
We call such maps \emph{potentialized polynomial lenses}.

\begin{lemma}\label{lem.cot_lifts_lens_potential}
The strong symmetric monoidal functor $\cot\colon\mfd\to\poly$ lifts to a strong symmetric monoidal functor
\[
\Lens\cot\colon\lmfd^\rr\to\Lens{\poly}^\rr.
\]
\end{lemma}

\begin{proof}
Apply \cref{lem.lens_T_functoriality} with $F=\cot$ (strong symmetric monoidal by \cref{prop.cot_monoidal}) and the monad morphism $\cot{\rr\times\blank}\Rightarrow\cot{\rr}\otimes\cot{\blank}$ supplied by \cref{lemma.monoid_to_monad}.
\end{proof}

In \cref{lem.Theta_poly_potential} we will apply the construction of \cref{prop.Theta_T_alpha}, which works for any $\Fun T$-monoid $(z,e,m,\alpha)$ in $\poly$. For simplicity of presentation, and with a view toward our physics application, we take $z\coloneqq\yon$. It has a unique monoid structure, but $\alpha\colon\cot{\rr}\to\yon$ is extra data.

\begin{lemma}\label{lem.alpha_constant}
The $(\cot{\rr}\otimes\blank)$-monoid structures $(z,e,m,\alpha)$ for which $z=\yon$ are in bijection with real numbers $r:\rr$.
\end{lemma}

\begin{proof}
There is a unique $\otimes$-monoid structure on $\yon$. Maps $\alpha\colon\cot{\rr}\to\yon$ are in bijection with (not-necessarily continuous) covector fields $\omega\colon\rr\to T^*\rr$ on $\rr$, sending $\alpha$ to $\omega\coloneqq\bk{\alpha}$. The first condition of \eqref{eqn.T_monoid} is vacuous. Thus the only condition on $\alpha$ is that the following equation holds for all $x_1,x_2\in\rr$
\[
\left(\omega(x_1),\omega(x_2)\right)=
\left(\omega(x_1+x_2),\omega(x_1+x_2)\right)
\]
Setting $x_2\coloneqq 0$, we find $\omega(0)=\omega(x_1)$ for all $x_1$. In other words, $\omega=r$ is constant at some $r:\rr$. And conversely, any constant covector field satisfies the condition.
 \end{proof}

The role of the Lie group $\rr$ in this section is to act as the recipient of \emph{potentials}. By \cref{lem.alpha_constant}, once we set $z\coloneqq\yon$, the remaining choice of $(\cot{\rr}\otimes\blank)$-monoid structure is a real number $r:\rr$, equivalently a constant covector field on $\rr$. We choose $r\coloneqq+1$, and write the corresponding action map as
\begin{equation}\label{eqn.d_potential}
\potd\colon\cot{\rr}\to\yon,\qquad\bk{\potd}(x)\coloneqq +1\in T^*_x\rr.
\end{equation}
We call it $\potd$ because this is the choice for which potentials contribute their usual differentials. Indeed, for a smooth map $U\colon X\to\rr$, the composite $\potd\circ\cot{U}\colon\cot{X}\to\yon$ corresponds to the covector field $dU$ on $X$: tracing on directions,
\[
\bk{\potd\circ\cot{U}}|_x(*)=(T_x U)^\top(+1)=dU|_x
\]
exactly as in \eqref{eqn.differential}.

\begin{lemma}\label{lem.Theta_poly_potential}
There is a lax monoidal functor
\[
\Theta_{\potd}\colon\Lens{\poly}^\rr\to\poly,\qquad
\Theta_{\potd}\lensob p\coloneqq\ihom{\inpt p,\yon}\otimes\outp p.
\]
%It sends a map $f\colon\lensob p\to\lensob q$, i.e.\ $\outp f\colon \outp p\to\outp q$ and $\inpt f\colon\inpt q\otimes\outp p\to\cot{\rr}\otimes\inpt p$, to the composite
%\[
%\begin{tikzcd}[column sep=45pt, row sep=22pt, ampersand replacement=\&]
%\outp p\otimes\ihom{\inpt p,\yon}\ar[r,"{\delta_{\outp p}\otimes\id}"]\ar[dd, dashed, "\Theta_{\potd}(f)"']
%\& \outp p\otimes\outp p\otimes\ihom{\inpt p,\yon}\ar[r,"\outp f\otimes\coev"]
%\& \outp q\otimes\ihom{\inpt q,\inpt q\otimes\outp p\otimes\ihom{\inpt p,\yon}}\ar[d,"{\id\otimes\ihom{\inpt q,\inpt f\otimes\id}}"]
%\\
%\&
%\& \outp q\otimes\ihom{\inpt q,\Fun T\inpt p\otimes\ihom{\inpt p,\yon}}\ar[d,"{\id\otimes\ihom{\inpt q,\sigma}}"]
%\\
%\outp q\otimes\ihom{\inpt q,\yon}
%\& \outp q\otimes\ihom{\inpt q,\Fun T\yon}\ar[l,"{\id\otimes\ihom{\inpt q,\potd}}"]
%\& \outp q\otimes\ihom{\inpt q,\Fun T(\inpt p\otimes\ihom{\inpt p,\yon})}\ar[l,"{\id\otimes\ihom{\inpt q,\Fun T\ev}}"]
%\end{tikzcd}
%\]
\end{lemma}

\begin{proof}
Follows from \cref{lemma.monoid_to_monad,lem.alpha_constant,prop.Theta_T_alpha}.
\end{proof}

\subsection{The dynamics functor}\label{subsec.dynamics_functor}

By \cref{prop.TT_monoidal,prop.pnla_smc_action,prop.cot_monoidal}, we have two parallel strong symmetric monoidal functors $\pvect\to\poly$,%
\footnote{The vector spaces can be replaced by paired nilpotent Lie algebras; see \cref{rmk.pnla_generalization}.}
connected by the monoidal natural transformation $\rho\colon\cot{T^*}\Rightarrow\cot$ of \cref{lem.rho_natural}. By \cref{lem.parameter_lens_action}, we obtain an action of $\pvect$ on $\lmfd^\rr$ and two actions of $\pvect$ on $\poly$ by
\begin{align*}
\pvect\times\lmfd^\rr&\to\lmfd^\rr,&&V\cdot\lensob M\coloneqq\binom{\inpt M}{V\times\outp M},\\
\pvect\times\poly&\to\poly,&&V\cdot p\coloneqq\cot{V}\otimes p,\\
\pvect\times\poly&\to\poly,&&V\cdot p\coloneqq\cot{T^*V}\otimes p.
\end{align*}
We thus form $\Para$ constructions (\cref{sec.para_general}), which we denote
\[
\para{\pvect}{\lmfd^\rr}\qqand
\para{\cot}{\poly}\qqand
\para{\cot{T^*}}{\poly}.
\]

\begin{definition}\label{def.potlens}
The symmetric monoidal category of \emph{potentialized lenses} is denoted and defined by
\[
\potlens\coloneqq\para{\pvect}{\lmfd^\rr}.
\qedhere
\]
\end{definition}

\begin{lemma}\label{lem.potlens_to_para_poly}
There is a lax symmetric monoidal functor
\[
\cint\colon\potlens\to\para{\cot}{\poly}.
\]
\end{lemma}

\begin{proof}
We apply \cref{prop.para_square} to $\Lens\cot$ from \cref{lem.cot_lifts_lens_potential} and apply \cref{cor.parameterized_representation} to $\Theta_{\potd}\colon\Lens{\poly}^\rr\to\poly$ from
\cref{lem.Theta_poly_potential} to obtain the maps below; the desired functor is their composite:
\[
\potlens\To{\para\pvect{\Lens\cot}}
\para{\pvect}{\Lens\poly^\rr}\To{\para\pvect{\Theta_{\potd}}}
\para{\cot}{\poly}.
\qedhere
\]
\end{proof}
We call $\cint$ the \emph{cotangent internalization} functor: it first applies $\cot$ to a potentialized manifold lens, then uses $\Theta_{\potd}$ to internalize the resulting potentialized polynomial lens as a parameterized map in $\poly$.


\begin{lemma}\label{lem.para_rho}
The monoidal natural transformation $\rho\colon\cot{T^*}\Rightarrow\cot$ of \cref{lem.rho_natural} induces a strong symmetric monoidal functor
\[
\leg\colon\para{\cot}{\poly}\to\para{\cot{T^*}}{\poly}.
\]
We call it the \emph{Legendre refinement} functor.
\end{lemma}

\begin{proof}
Apply \cref{prop.para_strong_induced,prop.para_square} to $\rho$.
\end{proof}

\begin{lemma}\label{lem.poly_to_org}
With the action $V\cdot p=\cot{T^*V}\otimes p$ of $\pvect$ on $\poly$, there is an identity-on-objects lax symmetric monoidal functor
\[
\dyn\colon\para{\cot{T^*}}{\poly}\to\org.
\]
We call it the \emph{dynamical realization} functor.
\end{lemma}

\begin{proof}
Let $F\colon\pvect\to\smset_\cong$ be the strong monoidal composite $\pvect\To{T^*}\pvect\to\smset_\cong$, where the second functor forgets the paired-vector-space structure. By \cref{prop.canonical_symplectic_pairing,prop.pnla_polynomial}, the isomorphism $\cot{T^*V}\cong F(V)\yon^{F(V)}$ fills the action square
\[
\begin{tikzcd}[ampersand replacement=\&, column sep=60pt]
\pvect\times\poly\ar[r,"\cot{T^*\!\blank}\otimes\blank"]\ar[d,"F\times\poly"', ""{name=L}]\& \poly\ar[d, equal, ""' {name=R}]\\
\smset_\cong\times\poly\ar[r,"\Fun{Store}(\blank)\otimes\blank"']\& \poly.
\arrow[Rightarrow, from=L, to=R, shorten <=1.5cm, shorten >=1.5cm, "\cong"]
\end{tikzcd}
\]
\Cref{prop.para_square} gives $\para{\cot{T^*}}\poly\to\para{\smset_\cong}\poly$, and composing with $\para{\smset_\cong}\poly\cong\org_\cong$ from \cref{prop.org_as_para} yields the claim.
\end{proof}

\begin{theorem}\label{thm.functor}
We have a lax symmetric monoidal functor
\[
\Phi\colon\potlens\to\org.
\]
\end{theorem}

\begin{proof}
Define $\Phi$ to be the composite of the maps below, given by \cref{lem.potlens_to_para_poly}, \cref{lem.para_rho}, and \cref{lem.poly_to_org}:
\[
\potlens\To{\cint}
\para{\cot}{\poly}\To{\leg}
\para{\cot{T^*}}{\poly}\To{\dyn}
\org.
\qedhere
\]
\end{proof}

In words, $\Phi$ first internalizes the cotangent lenses, then adds a ``momentum term'' to the state spaces using the Legendre refinement, and finally realizes the result as dynamic algebra. We now unpack $\Phi$ at the ``nuts and bolts'' level, tracing each term and sign in the resulting formulas back to the definitions that produced it.

Let $\inpt M,\outp M:\mfd$ be manifolds, and write
\[\Omega(M)\coloneqq\{\omega\colon M\to T^*M\mid\pi\circ\omega=\id\}\]
for the set of (not-necessarily continuous) \emph{covector fields} on $M$. A smooth map $f\colon M\to N$ induces a \emph{pullback} $f^*\colon\Omega(N)\to\Omega(M)$ defined by $(f^*\omega)(m)\coloneqq(T_m f)^\top\omega(f(m))$.

\paragraph{The action of $\Phi$ on objects.}
The functor $\Phi$ acts on objects as follows:
\begin{equation}\label{eqn.Phi_on_obs}
\Phi\lensob M=
\ihom{\cot{\inpt M},\yon}\otimes\cot{\outp M}\cong
\sum_{(\omega,\outp m):\Omega(\inpt M)\times\outp M}\yon^{\inpt M\times\left(T^*_{\outp{m}}\outp M\right)}.
\end{equation}
A position in $\Phi\lensob M$ is a pair $(\omega,\outp m)$ where $\omega\in\Omega(\inpt M)$ is a covector field and $\outp m\in\outp{M}$ is a point. A direction at $(\omega,\outp m)$ is a pair $(\inpt m,\xi)$ where $\inpt m\in\inpt M$ is a point and $\xi:T^*_{\outp{m}}\outp M$ is a covector.

\paragraph{The action of $\Phi$ on morphisms.}
Describing the action of $\Phi$ on morphisms will extend to the end of this section, \cref{subsec.dynamics_functor}. Let $(V,\sharp_V):\pvect$ and let
$f\colon V\cdot\lensob M\to\lensob N$
be a morphism in $\potlens$, i.e.\ a tuple $(\outp f,\inpt f,U)$ as in \eqref{eqn.para_potential_lens_maps}:
\begin{equation}\label{eqn.para_potential_lens_maps}
\outp f\colon V\times\outp M\to\outp N,\qquad
\inpt f\colon V\times\outp M\times\inpt N\to\inpt M,\qquad
U\colon V\times\outp M\times\inpt N\to\rr.
\end{equation}
It is sent by $\Phi$ to a $[\Phi\lensob M,\Phi\lensob N]$-coalgebra with state set
\begin{equation}\label{eqn.state_space}
X\coloneqq T^*V\cong V\oplus V^*;
\end{equation}
we write $x=(v,\xi)\in X$ for position $v\in V$ and momentum $\xi\in V^*$, and equip $X$ with the canonical symplectic pairing;%
\footnote{This comes from \cref{def.legendre_projection,prop.canonical_symplectic_pairing}.}
its sharp map is
\begin{equation}\label{eqn.sharp_X}
\sharp_X\colon X^*\to X,\qquad (\xi',v')\mapsto(v',-\xi'),
\end{equation}
specializing \eqref{eqn.canonical_sharp}.

We now unpack the coalgebra function itself,
\[
X\To{\Phi(V,\sharp_V,f)}
\left[\Phi\lensob M,\Phi\lensob N\right]\tri X.
\]
Given a state $x\coloneqq(v,\xi)\in X$ and a position $a\coloneqq(\omega_M,\outp m):\Omega(\inpt M)\times\outp M$, define the output point by
\begin{equation}\label{eqn.outpn}
\outp n\coloneqq\outp f(v,\outp m)
\end{equation}
and the output covector field $\omega_N\colon\inpt N\to T^*\inpt N$ by
\begin{equation}\label{eqn.omegaprime}
\omega_N(\inpt n)\;=\;(\inpt f(v,\outp m,\blank))^*\omega_M\big|_{\inpt n}\;+\;d(U(v,\outp m,\blank))\big|_{\inpt n}.
\end{equation}
The first term is the pullback of $\omega_M$ along the partial map $\inpt f(v,\outp m,\blank)\colon\inpt N\to\inpt M$; the second is the differential of the potential $U(v,\outp m,\blank)\colon\inpt N\to\rr$ at $\inpt n$, in the sense of \eqref{eqn.differential}.%
\footnote{The $+1$ that turns the second summand into $dU$ comes from \eqref{eqn.d_potential}; the $+$ between summands is the cotangent splitting $T^*(\rr\times\inpt M)\cong T^*\rr\oplus T^*\inpt M$ (\cref{sec.manifolds_notation}).} Then the readout position is
\[
b\coloneqq\bigl(\omega_N,\,\outp n\bigr)\;\in\;\Omega(\inpt N)\times\outp N.
\]


Now given a direction $(\inpt n,\xi_N):\inpt N\times T^*_{\outp n}\outp N=\Phi\lensob N[b]$, the coalgebra $\Phi(f)$ at $(x,a)$ produces a direction in $\inpt M\times T^*_{\outp m}\outp M$ as follows. Let
\begin{equation}\label{eqn.inptm}
\inpt m\coloneqq\inpt f(v,\outp m,\inpt n)
\end{equation}
and define
\begin{equation}\label{eqn.bigtheta}
(\xi_V,\,\xi_M,\,\xi_{\inpt N})\coloneqq (T_{(v,\outp m)}\outp f)^\top\xi_N \;+\; (T_{(v,\outp m,\inpt n)}\inpt f)^\top\omega_M(\inpt m) \;+\; dU\big|_{(v,\outp m,\inpt n)}\;
\end{equation}
where $(\xi_V,\,\xi_M,\,\xi_{\inpt N}):V^*\oplus T^*_{\outp m}\outp M\oplus T^*_{\inpt n}\inpt N$.%
\footnote{Each summand is a cotangent pullback: $\xi_N$ along $\outp f$, $\omega_M(\inpt m)$ along $\inpt f$, and $+1$ along $U$ (via $\potd$, \eqref{eqn.d_potential}); the $+$'s are the cotangent splitting of $V\times\outp M\times\inpt N$ (\cref{sec.manifolds_notation}), morally the domain of the maps in \eqref{eqn.para_potential_lens_maps}. Throughout, we silently use the canonical chart identification \eqref{eqn.tangent_vec} to treat covectors at points of vector spaces as elements of the dual.}
Note that the first summand naturally lands in $V^*\oplus T^*_{\outp m}\outp M$ and has been extended by $0$ in the remaining $\inpt N$-factor, so by \eqref{eqn.omegaprime} we have $\omega_N(\inpt n)=\xi_{\inpt N}$.

All of the work has now been done: the coalgebra at $(x,a)$ on $(\inpt n,\xi_N)$ returns the direction
\[
(\inpt m,\,\xi_M):\inpt M\times T^*_{\outp m}\outp M=\Phi\lensob M[a]
\]
and updates the state (from its original $x=(v,\xi)$) to
\begin{equation}\label{eqn.state_update}
x+\sharp_X(\xi_V,\sharp_V(\xi))=\bigl(v+\sharp_V(\xi),\,\xi-\xi_V\bigr).
\end{equation}
The sum $x+\sharp_X(\blank)$ is the exponential of the paired vector space $X$, \eqref{eqn.flow_cot}; the use of $\sharp_X$ comes from the canonical symplectic pairing on $X$, \eqref{eqn.canonical_sharp}; and the pair $(\xi_V,\sharp_V(\xi))\in T^*_xX$ comes from the Legendre projection, \cref{def.legendre_projection}. The right-hand side is the computation: \eqref{eqn.canonical_sharp} gives $\sharp_X(\xi_V,\sharp_V(\xi))=(\sharp_V(\xi),-\xi_V)$, then componentwise addition with $x=(v,\xi)$ produces $(v+\sharp_V(\xi),\,\xi-\xi_V)$.

All in all, one could tersely denote the coalgebra map as follows:
\[
\Phi(V,\sharp_V,f)\colon 
(v,\xi)\mapsto
(\omega_M,\outp m)\mapsto
\Big(
	(\omega_N,\outp n),
	(\inpt n,\xi_N)\mapsto
	\big(
		(\inpt m,\xi_M),
		(v+\sharp_V(\xi),\,\xi-\xi_V)
	\big)
\Big),
\]
where $\outp n$ is defined in \eqref{eqn.outpn}, $\omega_N$ in \eqref{eqn.omegaprime}, $\xi_M$ and $\xi_V$ in \eqref{eqn.bigtheta}, and $\inpt m$ in \eqref{eqn.inptm}.

\chapter{Example: the wave equation}\label{sec.spring}

We illustrate $\Phi\colon\potlens\to\org$ by deriving the wave equation. 

That is, we begin with a single particle with position and momentum, which we put in a ``box'' that in some sense delivers the particle's position to its neighbors and updates its state based on some fixed parameters $\kappa,m>0$ and the locations of its neighboring particles. We wire these boxes together in series
\begin{equation}\label{eqn.wave_wd}
\begin{tikzpicture}[oriented WD, bb port length=4pt, bb port sep=1.5, bb min width=.35cm, bbx=.4cm, bby=.8ex, baseline=(lbox1), particle/.style={path picture={\fill (path picture bounding box.center) circle (2.4pt);}}]
  \def\N{7}
  \node[bb={1}{1}, particle] (lbox1) {};
  \foreach \i [evaluate=\i as \prev using int(\i-1)] in {2,...,\N} {
    \node[bb={1}{1}, particle, right=.7 of lbox\prev] (lbox\i) {};
    \draw (lbox\prev_out1) -- (lbox\i_in1);
  }
  \node[right=.4 of lbox\N] (dots) {$\cdots$};
  \draw (lbox\N_out1) -- (dots.west);
  \node[bb={1}{1}, particle, right=.4 of dots] (rbox1) {};
  \draw (dots.east) -- (rbox1_in1);
  \foreach \i [evaluate=\i as \prev using int(\i-1)] in {2,...,\N} {
    \node[bb={1}{1}, particle, right=.7 of rbox\prev] (rbox\i) {};
    \draw (rbox\prev_out1) -- (rbox\i_in1);
  }
  \node[bb={0}{0}, fit={(lbox1) (rbox\N)}] (Y) {};
  \node[coordinate] at (Y.west|-lbox1_in1) (Y_in1) {};
  \node[coordinate] at (Y.east|-rbox\N_out1) (Y_out1) {};
  \draw (Y_in1) -- (lbox1_in1);
  \draw (rbox\N_out1) -- (Y_out1);
\end{tikzpicture}
\end{equation}
and calculate that the resulting state dynamics evolves according to the discrete wave equation:%
\footnote{Taking the continuum limit gives the usual wave equation.}
\begin{equation}\label{eqn.discrete_wave}
m\ddot v_i=\kappa(v_{i-1}-2v_i+v_{i+1}),
\end{equation}
where $\ddot v_i(t)=v_i(t+2)-2v_i(t+1)+v_i(t).$

We will work in the full suboperad of $\potlens$ (in the sense of \cref{sec.wd_operads}) spanned by a single object $\boxob\coloneqq\binom{\rr}{\rr}$, corresponding to a box whose input port and output port are both $\rr$. That is, we'll consider morphisms of the form
\[
\varphi\colon\potlens(\boxob\tpow{n},\boxob)
\]
for varying $n$. Note that $\boxob\tpow{0}=\binom{\rr^0}{\rr^0}$ is the monoidal unit.

First we give a $0$-ary morphism $\Part\colon\binom{\rr^0}{\rr^0}\to\binom{\rr}{\rr}$, for the dynamical system that goes in each little box, representing the particle of mass $m>0$. Then we give a $K$-ary morphism for each $K>0$ that describes the (static) wiring diagram shown in \eqref{eqn.wave_wd}. Then we calculate what happens when we compose.

\paragraph{The particle.}

To define a particle $\Part\colon\binom{\rr^0}{\rr^0}\to\binom{\rr}{\rr}$ in $\potlens$ \eqref{eqn.para_potential_lens_maps}, we need a paired vector space $(V,\sharp_V)$, a smooth ``forward'' map $\outp f\colon V\to\rr$, and a smooth ``potential'' $U\colon V\times\rr\to\rr$. We take the parameter space to be $\rr\in\pvect$ with sharp $\sharp_\rr(p)\coloneqq p/m$. We take the forward map to be the identity, $\outp f(x)\coloneqq x$, and the potential to be $U(x,y)\coloneqq\tfrac{\kappa}{2}(x-y)^2$. This is called the \emph{harmonic} potential.

\paragraph{The wiring diagram.}
For the wiring diagram, we start in $\Lens{\finset\op}$; see \cref{ex.lens_finsetop}. The $K$-ary series composition shown in \eqref{eqn.wave_wd} corresponds to a lens $\phi_K\colon\binom{K}{K}=\binom{1}{1}+\cdots+\binom{1}{1}\to\binom{1}{1}$ given by
\[
K\From{\outp f}1
\qqand
K+1\From{\inpt f} K,
\]
where we identify $K+1$ with the disjoint union of the big-box input and the $K$ small-box outputs, the big-box input first. Then $\outp f$ picks out the last wire, $\outp{f}(*)=K$, and $\inpt f$ picks out the first $K$ wires, $\inpt{f}(n)=n$.

Applying $\rr^-$ from \cref{lemma.lens_rr}, we obtain the lens $\binom{\rr^K}{\rr^K}\to\binom{\rr}{\rr}$ given by
\[
\outp f(v_1,\ldots,v_K)\coloneqq v_K
\qqand
\inpt f\bigl((v_1,\ldots,v_K),\,v_0\bigr)\coloneqq(v_0,\,v_1,\,\ldots,\,v_{K-1}).
\]
Lenses without potential and without parameterization are special cases of $\potlens$-morphisms, i.e.\ taking $V'\coloneqq\rr^0$ and the $0$-potential $U'\coloneqq0\colon\rr^K\times\rr\to\rr$, we have the wiring as a $\potlens$-morphism
\[
\fun{wire}_K\colon\boxob\tpow{K}\to\boxob,\qquad (\outp f,\,\inpt f,\,0)
\]

Composing, we obtain a $0$-ary map, $\fun{wire}_K(\Part,\ldots,\Part)$ in $\potlens$, with parameter $V\coloneqq V'\oplus\rr^{\oplus K}=\rr^K$ (one $\rr$ per particle) and sharp
\begin{equation}\label{eqn.sharp_chain}
\sharp_V\colon V^*\to V,\qquad p\mapsto p/m,
\end{equation}
acting per-particle on each $\rr$-summand (the inverse-mass pairing). The defining maps are
\[
\outp f(v_1,\ldots,v_K)\coloneqq v_K
\qqand
U\bigl((v_1,\ldots,v_K),\,v_0\bigr)\coloneqq\sum_{i=1}^{K}\tfrac{\kappa}{2}(v_i-v_{i-1})^2,
\]
with $\inpt f$ vacuous. The sum in $U$ comes from the productor \eqref{eqn.monoid_productor} of the monoidal monad $\rr\times\blank$ on $\Lens\mfd$, where $(\rr,0,+)$ is the additive monoid. Composition with $\fun{wire}_K$ then substitutes $y_i\mapsto v_{i-1}$ via $\inpt f$ (with $v_0$ the external input).

\paragraph{The discrete dynamics.}

The discrete dynamics appear when we apply the functor $\Phi\colon\potlens\to\org$. The object $\boxob$ is sent to
\[
\Phi(\boxob)=\sum_{(\omega,v):\Omega(\rr)\times\rr}\yon^{\rr\times T^*_v\rr}.
\]
The map $\fun{wire}_K(\Part,\ldots,\Part)$ is sent to a coalgebra
\[
\Phi(\fun{wire}_K(\Part,\ldots,\Part))\colon X\To{}\Phi(\boxob)\tri X
\]
with state space $X\coloneqq T^*\rr^K$, with elements written $x=((v_1,\ldots,v_K),(p_1,\ldots,p_K))$ for positions and momenta (so $\xi=\vec p$ in terms of \eqref{eqn.state_update}), and with the function given by
\begin{equation}\label{eqn.wave_coalg}
(\vec v,\vec p)\mapsto\Bigl(\bigl(\omega_\boxob,\,v_K\bigr),\;(v_0,\xi_N)\mapsto\bigl(\vec v+\tfrac{1}{m}\vec p,\,\vec p-\xi_V\bigr)\Bigr),
\end{equation}
where each piece is read off from the unpacking in \cref{subsec.dynamics_functor}:
\begin{itemize}
\item The output position $v_K$ comes from $\outp f$ via \eqref{eqn.outpn}.
\item The \emph{harmonic spring} covector field on the lefthand side,
\[\omega_\boxob(v_0)=d\bigl[\tfrac{\kappa}{2}(v_1-v_0)^2\bigr]\big|_{v_0}=-\kappa(v_1-v_0),\]
comes from \eqref{eqn.omegaprime} (with $\inpt f$ and $\omega_M$ vacuous, and $dv_0=+1$ via \eqref{eqn.d_potential}).
\item The pullback covector $\xi_V\in T^*_{\vec v}V$ comes from \eqref{eqn.bigtheta}. It is given component-wise (via \eqref{eqn.tangent_vec} identifying $T^*_{\vec v}V\cong V^*=(\rr^K)^*$) by
\begin{align*}
(\xi_V)_i&=
\bigl((T_{\vec v}\outp f)^\top\xi_N\bigr)_i+\partial_{v_i}U(\vec v,v_0)\\&=
\begin{cases}\partial_{v_i}U(\vec v,v_0)+0&1\leq i<K\\ \partial_{v_K}U(\vec v,v_0)+\xi_N&i=K\end{cases}\\&=
\begin{cases}\kappa(2v_i-v_{i-1}-v_{i+1})&1\leq i<K\\ \kappa(v_K-v_{K-1})+\xi_N&i=K\end{cases}
\end{align*}
\item The state update applied to $x=(\vec v,\vec p)$ comes from \eqref{eqn.state_update}: using our sharp $\sharp_V(\vec p)=\vec p/m$ from \eqref{eqn.sharp_chain} and the symplectic sharp $\sharp_X\colon(\xi,v)\mapsto(v,-\xi)$ from \eqref{eqn.sharp_X}, we get
\[
x'\coloneqq x+\sharp_X(\xi_V,\sharp_V(\vec p))=(\vec v+\vec p/m,\,\vec p-\xi_V).
\]
In other words, writing $x'=(v_1',\ldots,v_K',p_1',\ldots,p_K')$, we have
\[
v_i'\coloneqq v_i+p_i/m
\qqand
p_i'\coloneqq\begin{cases}p_i+\kappa(v_{i-1}+v_{i+1}-2v_i)&1\leq i<K\\ p_i+\kappa(v_{K-1}-v_K)-\xi_N&i=K.\end{cases}
\]
\end{itemize}

\paragraph{Recovering the wave equation.}
To recover the discrete wave equation everywhere, we pin the chain between two walls: choose left wall $v_0\coloneqq 0$ and right wall $v_{K+1}\coloneqq L(K+1)$, where $L$ is the equilibrium spacing between adjacent particles. Supplying these via the external direction-input means setting
\begin{equation}\label{eqn.wave_bc}
v_0\coloneqq 0\qqand\xi_N\coloneqq \kappa\bigl(v_K-L(K+1)\bigr).
\end{equation}
With these choices, the boundary case at $i=K$ becomes $p_K'=p_K+\kappa\bigl(v_{K-1}+L(K+1)-2v_K\bigr)$, so the above state update formula reduces to this:
\begin{equation}\label{eqn.state_update_explicit}
v_i'\coloneqq v_i+p_i/m
\qqand
p_i'\coloneqq p_i+\kappa(v_{i-1}+v_{i+1}-2v_i).
\end{equation}

We then think of the state $x$ as a function of time: fix any initial $x(0)\coloneqq(\vec{v}(0),\vec{p}(0))$ and set $x(t+1)\coloneqq x(t)'$ as in \eqref{eqn.state_update_explicit}. We then rewrite the first equation as $p_i(t)=mv_i(t+1)-mv_i(t)$ and substitute it into the second:
\[
mv_i(t+2)-mv_i(t+1)=mv_i(t+1)-mv_i(t)+\kappa(v_{i-1}(t)+v_{i+1}(t)-2v_i(t))
\]
Rearranging, we obtain the wave equation formula \eqref{eqn.discrete_wave} with the same notation as used there.

\section{Functoriality}

We have now recovered the wave equation by connecting $K$-many particles together in $\potlens$, sending the result to $\org$ and investigating the discrete dynamical system (the $(\cot{\rr}\otimes[\cot{\rr},\yon])$-coalgebra) as a function of time. This is the lefthand side below, which by functoriality of $\Phi$ is equal to the righthand side:
\[
\Phi\bigl(\fun{wire}_K(\Part,\ldots,\Part)\bigr)=\Phi(\fun{wire}_K)\bigl(\Phi(\Part),\ldots,\Phi(\Part)\bigr),
\]
which is the $K$-ary composite, taken in $\org$, of $K$ copies of the single-particle dynamical system $\Phi(\Part)$ wired together via $\Phi(\fun{wire}_K)$. We work this out as well.

\paragraph{The particle dynamics.}

We first apply $\Phi$ to $\Part\colon\binom{\rr^0}{\rr^0}\to\boxob$ to obtain a coalgebra on
\[
  \left[\Phi\binom{\rr^0}{\rr^0},\Phi(\boxob)\right]=
  \left[\yon,\Phi(\boxob)\right]=\boxob=
  \sum_{(\omega,v):\Omega(\rr)\times\rr}\yon^{\rr\times T^*_v\rr}
\]
We took the parameter space to be $\rr$, so by \eqref{eqn.state_space} the state space of $\Phi(\Part)$ is $X\coloneqq T^*\rr$, with elements written $x=(v,p)$ for position and momentum. The map
\[
X\To{\Phi(\Part)}
\sum_{(\omega,v):\Omega(\rr)\times\rr}X^{\rr\times T^*_v\rr}
\]
is given by
\[
(v,p)\mapsto\Bigl(\bigl(v_{\fun{prev}}\mapsto\kappa(v_{\fun{prev}}-v),\,v\bigr),\;(v_{\fun{prev}},\xi_{\fun{next}})\mapsto\bigl(v+p/m,\,p-\xi_{\fun{next}}+\kappa(v_{\fun{prev}}-v)\bigr)\Bigr).
\]
In other words, starting at position $v$ and momentum $p$, it outputs $v$ to the right and the covector field $v_{\fun{prev}}\mapsto\kappa(v_{\fun{prev}}-v)$, which ``pulls'' $v_{\fun{prev}}$ toward $v$. Then, receiving a covector $\xi_{\fun{next}}$ from the right and a point $v_{\fun{prev}}$ from the left, it updates its position by adding velocity (momentum divided by mass) and updates its momentum by adding the spring force $\kappa(v_{\fun{prev}}-v)$, which ``pulls'' $v$ toward $v_{\fun{prev}}$, and subtracting the received covector $\xi_{\fun{next}}$.
 
\paragraph{The static wiring.}

Apply $\Phi$ to $\fun{wire}_K\colon\boxob^{\otimes K}\to\boxob$. With trivial parameter $\rr^0$, the state space $T^*\rr^0$ is a single point, so $\Phi(\fun{wire}_K)$ is \emph{static}: a function $1\to[\Phi(\boxob)^{\otimes K},\Phi(\boxob)](1)$ is equivalently (see \eqref{eqn.dirichlet_hom}) a polynomial map $\Phi(\boxob)^{\otimes K}\to\Phi(\boxob)$, i.e.\
\[\sum_{(\vec\omega,\vec v):\Omega(\rr)^K\times\rr^K}\yon^{\rr^K\times T^*_v\rr^K}\To{\Phi(\fun{wire}_K)}\sum_{(\omega,v):\Omega(\rr)\times\rr}\yon^{\rr\times T^*_v\rr}.\]
To each $(\vec \omega,\vec v)$, the vector of outputs of the internal boxes, we assign the external output $(\omega_1,v_K)$, and to each external input $(v_0,\omega_K)$; in our concise format,
\[
(\vec \omega,\vec v)\mapsto\Big((\omega_1,v_K),\;(v_0,\omega_K)\mapsto\bigl((v_0,0),\,(v_1,0),\,\ldots,\,(v_{K-2},0),\,(v_{K-1},\omega_K)\bigr)\Big).
\]


\paragraph{Composition.}

The $\org$-composite $\Phi(\fun{wire}_K)\bigl(\Phi(\Part),\ldots,\Phi(\Part)\bigr)$ stitches together the $K$ single-particle states into $T^*\rr\oplus\cdots\oplus T^*\rr=T^*\rr^K$. The wiring's static action substitutes each particle's neighbor input $y_i\mapsto v_{i-1}$ (so the $i$-th $\Phi(\Part)$ sees its left neighbor's position) and pipes $\xi_N$ into box $K$, exactly reproducing \eqref{eqn.wave_coalg}.

DISCARD AFTER THIS:

no internal dynamics, only a $\poly$-level transformation between $\Phi(\boxob)^{\otimes K}$ and $\Phi(\boxob)$. It implements the wiring rules from \eqref{eqn.wave_wd}:
\begin{itemize}
\item The big-box output position is box $K$'s output position; the output covector field is $\omega^{\mathrm{out}}(v_0)\coloneqq\omega_1(v_0)$, the leftmost box's covector field at $v_0$ (\eqref{eqn.omegaprime} with $U=0$, since only box 1 sees the external input).
\item For each box $i$, the input position $y_i\coloneqq v_{i-1}$ is the left neighbor's output (with $v_0$ the big-box input), from $\inpt f$.
\item The external direction $\xi_N$ feeds into box $K$'s $\xi$-slot, since $\outp f(\vec v)=v_K$ (\eqref{eqn.bigtheta}); all other boxes' $\xi$-slots are zero.
\end{itemize}



The box has input and output $M$, i.e.\ it comes from the object $\binom{\rr}{\rr}$. 
wired together  with a one-dimensional spring cell: a mass coupled to one neighbor by a Hooke spring. The state space comes out to $T^*\rr$, and the dynamics are Hamilton's equations for the harmonic oscillator.

\begin{definition}\label{def.spring_cell}
Fix $\kappa,\mu>0$. The \emph{spring cell} is the $\potlens$-morphism
\[
f_\mathrm{spring}\colon\binom{\rr}{\rr}\to\binom{\rr}{\rr}
\]
with parameter $V\coloneqq(\rr,0,+)$ paired by $\sharp_V(p)\coloneqq p/\mu$, and data
\[
\outp f(v,\outp m)\coloneqq v,
\quad
\inpt f(v,\outp m,\inpt n)\coloneqq\inpt n,
\quad
U(v,\outp m,\inpt n)\coloneqq\tfrac{\kappa}{2}(v-\inpt n)^2.
\]
The output reports the cell's own position $v$; the input passes the neighbor's position $\inpt n$ through; $U$ is the Hooke potential between cell and neighbor.
\end{definition}

\begin{remark}\label{rmk.spring_pairing}
The pairing $\sharp_V(p)=p/\mu$ encodes inverse mass; the standard Euclidean pairing $\sharp_V(p)=p$ corresponds to $\mu=1$.
\end{remark}

By \cref{def.spring_cell}, the state set of $\Phi(f_\mathrm{spring})$ is $T^*V\cong V\oplus V^*=\rr\oplus\rr$, written $(v,p)$ for position $v\in V$ and momentum $p\in V^*$.

\begin{proposition}\label{prop.spring_dynamics}
At state $(v,p)\in T^*V$ on a position $a=(\omega,\outp m)$ with $\omega=0$ and input direction $(\inpt n,\xi_N)$ with $\xi_N=0$, the coalgebra $\Phi(f_\mathrm{spring})$ updates the state to
\[
\dot v = p/\mu,
\qquad
\dot p = \kappa(\inpt n-v).
\]
\end{proposition}

\begin{proof}
Apply \eqref{eqn.bigtheta} at $(v,\outp m,\inpt n)$. With $\xi_N=0$ and $\omega=0$, the first two summands of $\Xi$ vanish, leaving
\[
\Xi=dU\big|_{(v,\outp m,\inpt n)}
\;=\;\bigl(\partial_V U,\;\partial_{\outp M}U,\;\partial_{\inpt N}U\bigr)
\;=\;\bigl(\kappa(v-\inpt n),\;0,\;-\kappa(v-\inpt n)\bigr).
\]
Hence $\dot v=\sharp_V(p)=p/\mu$ and $\dot p=-\Xi_V=-\kappa(v-\inpt n)=\kappa(\inpt n-v)$.\qedhere
\end{proof}

\begin{remark}[Harmonic oscillator]\label{rmk.spring_hamiltonian}
The update of \cref{prop.spring_dynamics} is Hamilton's equations $\dot v=\partial_p H$, $\dot p=-\partial_V H$ for the harmonic-oscillator Hamiltonian
\[
H(v,p)=\frac{p^2}{2\mu}+\tfrac{\kappa}{2}(v-\inpt n)^2,
\]
namely a mass $\mu$ in a Hooke potential well centered at the neighbor's position $\inpt n$. Differentiating $\dot v=p/\mu$ recovers Newton's second law $\mu\ddot v=\kappa(\inpt n-v)$.
\end{remark}

\begin{proposition}\label{prop.spring_wave}
Wiring a $\zz$-indexed chain of spring cells, with each cell's input $\inpt n_i$ identified by the wiring with its right neighbor's output $v_{i+1}$ and feedback in the backward direction giving the left-neighbor coupling, yields the discrete wave equation
\[
\mu\ddot v_i=\kappa(v_{i-1}-2v_i+v_{i+1}).
\]
In the continuum limit $v_i\to u(i\Delta x,t)$ with density $\rho\coloneqq\mu/\Delta x$ and tension $T\coloneqq\kappa\Delta x$, this approaches the wave equation $\rho\,u_{tt}=T\,u_{xx}$.
\end{proposition}

\begin{proof}[Sketch]
Each cell $i$'s own potential contributes $\kappa(\inpt n_i-v_i)=\kappa(v_{i+1}-v_i)$ to $\dot p_i$ (right-side spring). The wiring's backward direction propagates cell $i$'s output covector field $\omega'_i$ (= $\Xi_{\inpt N}$ at cell $i-1$'s coupling) into cell $i-1$'s $\xi_V$, contributing $\kappa(v_{i-1}-v_i)$ (left-side spring). Summing, $\dot p_i=\kappa(v_{i-1}+v_{i+1}-2v_i)$; with $\dot v_i=p_i/\mu$, differentiating gives $\mu\ddot v_i=\kappa(v_{i-1}-2v_i+v_{i+1})$. Pass to the continuum limit.
\end{proof}

\clearpage
\clearpage

\appendix

\section{Potentialized manifold lenses}\label{sec.potlens_manifolds}

We now specialize the above constructions to the category of manifolds, acted on by paired vector spaces.%
\footnote{As mentioned in \cref{rmk.pnla_generalization}, everything about paired vector spaces generalizes to paired nilpotent Lie algebras.}

\begin{proposition}\label{prop.lie_monoid_strong_monad}
For any Lie monoid $G$, there is a strong monad
\[
\Fun T_G\colon\mfd\to\mfd,\qquad \Fun T_G(X)\coloneqq X\times G.
\]
\end{proposition}

\begin{corollary}\label{cor.lmfd_G}
For any Lie monoid $G$, there is a coKleisli category
\[
\Lens{\mfd}^{\Fun T_G}.
\]
\end{corollary}

\begin{corollary}\label{cor.lmfd_potential}
For $G=(\rr,0,+)$, there is a coKleisli category
\[
\lmfd^\rr\coloneqq\Lens{\mfd}^{\Fun T_\rr}.
\]
\end{corollary}

\begin{proposition}\label{prop.pvect_action_lmfd_potential}
There is a symmetric monoidal action
\[
\pvect\times\lmfd^\rr\to\lmfd^\rr,\qquad V\cdot\lensob M\coloneqq\binom{\inpt M}{V\times\outp M}.
\]
\end{proposition}

\begin{definition}\label{def.potlens}
\[
\potlens\coloneqq\para{\pvect}{\lmfd^\rr}.
\]
\end{definition}

\begin{proposition}\label{prop.potlens_maps}
A morphism $f\colon\lensob L\to\lensob M$ in $\potlens$ is a tuple $(V,\outp f,\inpt f,U)$ with
\[
\outp f\colon V\times\outp L\to\outp M,\qquad
\inpt f\colon V\times\outp L\times\inpt M\to\inpt L,\qquad
U\colon V\times\outp L\times\inpt M\to\rr.
\]
\end{proposition}

\begin{corollary}\label{cor.potlens_parameterized_representation}
Let $\cat C,\Fun T,z$, and $J\colon\pvect\to\cat C$ be as in \cref{cor.parameterized_representation}, with $\cat V=\pvect$. Suppose a lax symmetric monoidal functor
\[
\widehat F\colon\lmfd^\rr\to\Lens{\cat C}^{\Fun T}
\]
is equipped with coherent monoidal isomorphisms
\[
\widehat F(V\cdot\lensob M)\cong J(V)\cdot\widehat F(\lensob M).
\]
Then there is a lax symmetric monoidal functor
\[
\potlens=\para{\pvect}{\lmfd^\rr}\to\para{\pvect}{\cat C}.
\]
It sends a morphism $f\colon\lensob L\to\lensob M$ with parameter $V$ to a map
\[
J(V)\otimes\Xi_z\widehat F(\lensob L)\to\Theta_z\widehat F(\lensob M).
\]
\end{corollary}

\iffalse
Applying the abstract construction of \cref{sec.lenses_general} to $\cat C=\mfd$ (cartesian monoidal), the resulting category $\lmfd\coloneqq\Lens\mfd$ of \emph{manifold lenses} has pairs of manifolds $\lensob M$ as objects and pairs $f=(\outp f,\inpt f)$ of smooth maps as morphisms (\eqref{eqn.lens_def}).

Any Lie monoid $(G,\cdot,e)$ gives a strong monad on $\mfd$:
\[
\Fun T_G(X)\coloneqq X\times G,\qquad\eta_X(x)\coloneqq(x,e),\qquad\mu_X(x,g,g')\coloneqq(x,g\cdot g'),
\]
with strength $X\times(Y\times G)\To{\cong}(X\times Y)\times G$ by reassociation. The backward-monad construction (\cref{subsec.backwards}) produces a comonad $\Lens{\Fun T_G}$ on $\lmfd$; we write $\lmfd^G$ for its coKleisli category. A morphism $\lensob L\to\lensob M$ in $\lmfd^G$ has backward component
\[
\inpt f\colon\outp L\times\inpt M\to\inpt L\times G,
\]
a usual backward map augmented by a smooth $G$-valued map; coKleisli composition multiplies $G$-components.

The case of primary interest is $G=(\rr,0,+)$, the \emph{potential monad}: the $\rr$-component
\[
U\colon\outp L\times\inpt M\to\rr
\]
is the \emph{potential}, and coKleisli composition adds potentials.

We lift the $\pvect$-action on $\mfd$ (\cref{prop.pnla_smc_action}) to $\lmfd$ by acting on the $\outp M$-slot,
\[
(\cdot)\colon\pvect\times\lmfd\to\lmfd,\qquad V\cdot\lensob M\coloneqq\binom{\inpt M}{V\times\outp M},
\]
and define the category of \emph{potentialized lenses}
\[
\potlens\coloneqq\para{\pvect}{\lmfd^\rr}.
\]
A morphism $f\colon\lensob L\to\lensob M$ in $\potlens$ is a tuple $(V,\outp f,\inpt f,U)$ with $V$ a paired vector space,
\begin{equation}\label{eqn.map_paralens}
\outp f\colon V\times\outp L\to\outp M
\qqand
\inpt f\colon V\times\outp L\times\inpt M\to\inpt L,
\end{equation}
and a smooth potential $U\colon V\times\outp L\times\inpt M\to\rr$:
\[
\begin{tikzpicture}[oriented WD, bb port length=5pt, bb port sep=1, bb min width=.4cm, bbx=.5cm, bby=.7ex, baseline=(fo)]
  \node[bb={2}{1}] (fo) {$\outp f$};
  \node[bb={1}{3}, above left=4 and 4 of fo] (fi) {$\inpt f$};
  \node[bb={1}{3}, above right=3 and 0 of fi] (V) {$U$};
  \node[circle, draw, fill=black, inner sep=0, minimum size=3pt, left=1 of fo_in2] (dotL) {};
  \coordinate (Lo) at ($(dotL)+(-4.5,0)$);
  \coordinate (Mo) at ($(fo.east)+(1,0)$);
  \coordinate (Li) at (Lo |- fi_in1);
  \coordinate (Mi) at (Mo |- fi_out3);
  \node[circle, draw, fill=black, inner sep=0, minimum size=3pt] at ($(fi_out3)+(2,0)$) (dotMI) {};
  \node[left=0 of Lo] {$\outp L$};
  \node[right=0 of Mo] {$\outp M$};
  \node[left=0 of Li] {$\inpt L$};
  \node[right=0 of Mi] {$\inpt M$};
  \draw[ar] (Lo) to (dotL);
  \draw[ar] (dotL) to (fo_in2);
  \draw[ar] (fo_out1) to (Mo);
  \draw[ar] (fi_in1) to (Li);
  \draw[ar] (Mi) to (dotMI);
  \draw[ar] (dotMI) to (fi_out3);
  \draw[ar] (dotL) to[out=60, in=0] (fi_out2);
  \draw[ar] (dotMI) to[out=60, in=0] (V_out3);
  \draw[ar] (dotL) to[out=60, in=0] (V_out2);
  %
  \node [above left=1 and 0 of Lo] (P) {$V$};
  \node[circle, draw, fill=black, inner sep=0, minimum size=3pt] at (dotL|-P) (dotP) {};
  \draw[ar] (P) to (dotP);
  \draw[ar] (dotP) to (fo_in1);
  \draw[ar] (dotP) to[out=60, in=0] (fi_out1);
  \draw[ar] (dotP) to[out=60, in=0] (V_out1);
  \coordinate (R) at ($(V.west)+(-2,0)$);
  \node[left=0 of R] {$\rr$};
  \draw[ar] (V_in1) to (R);
\end{tikzpicture}
\]
The identity on $\lensob{L}$ has parameter $\rr^0:\pvect$, backward projection on $\lmfd$, and $U=0$ (the unit $e=0$ of $\rr$). Given also $g=(W,\outp g,\inpt g,U')\colon\lensob{M}\to\lensob{N}$, the composite $f\then g$ has parameter $V\times W$, inherits forward and backward components from $\lmfd$, and its potential adds the two via the coKleisli multiplication $\mu=+$:
\begin{align*}
U\bigl(v,\ell,\,\inpt g(w,\outp f(v,\ell),n')\bigr)+U'\bigl(w,\outp f(v,\ell),n'\bigr).
\end{align*}
The monoidal product is inherited from $\lmfd$. The following depicts composition, though note that we include it only for the algorithmically-minded reader's convenience: \emph{it follows from the abstract framework we already have developed}:
\[
\begin{tikzpicture}[oriented WD, bb port length=5pt, bb port sep=1, bb min width=.4cm, bbx=.5cm, bby=.7ex, baseline=(fo)]
  \node[bb={2}{1}] (fo) {$\outp f$};
  \node[bb={2}{1}, right=7 of fo] (go) {$\outp g$};
  \node[bb={1}{3}, above left=4 and 3 of fo] (fi) {$\inpt f$};
  \node[bb={1}{3}, above left=4 and 3 of go] (gi) {$\inpt g$};
  \node[bb={1}{3}, above right=1 and 0 of fi] (V) {$U$};
  \node[bb={1}{3}, above right=1 and 0 of gi] (W) {$U'$};
  \node[bb={1}{2}, above left=-1 and 2 of V] (plus) {$+$};
  %
  \node[circle, draw, fill=black, inner sep=0, minimum size=3pt, left=1 of fo_in2] (dotL) {};
  \node[circle, draw, fill=black, inner sep=0, minimum size=3pt] at ($(gi.east|-fo.east)+(1,0)$) (dotM) {};
  %
  \coordinate (Lo) at ($(dotL)+(-5,0)$);
  \coordinate (No) at ($(go.east)+(1,0)$);
  \coordinate (Li) at (Lo |- fi_in1);
  \coordinate (Ni) at (No |- gi_out3);
  \node[left=0 of Lo] {$\outp L$};
  \node[right=0 of No] {$\outp N$};
  \node[left=0 of Li] {$\inpt L$};
  \node[right=0 of Ni] {$\inpt N$};
  %
  \draw[ar] (Lo) to (dotL);
  \draw[ar] (dotL) to (fo_in2);
  \draw[ar] (fo_out1) to node[above, font=\scriptsize] {$\outp M$} (dotM);
  \draw[ar] (dotM) to (go_in2);
  \draw[ar] (go_out1) to (No);
  \draw[ar] (fi_in1) to (Li);
  \node[circle, draw, fill=black, inner sep=0, minimum size=3pt] at ($(gi_out3)+(3,0)$) (dotNI) {};
  \draw[ar] (Ni) to (dotNI);
  \draw[ar] (dotNI) to (gi_out3);
  \node[circle, draw, fill=black, inner sep=0, minimum size=3pt] at ($(fi_out3)+(3,0)$) (dotMI) {};
  \draw[ar] (gi_in1) to[out=180, in=0] node[above, font=\scriptsize] {$\inpt M$} (dotMI);
  \draw[ar] (dotMI) to (fi_out3);
  \draw[ar] (dotL) to[out=60, in=0] (fi_out2);
  \draw[ar] (dotM) to[out=60, in=0] (gi_out2);
  %
  \node[above left=1 and 0 of Lo] (P) {$V$};
  \node[above left=4 and 0 of Lo] (Q) {$W$};
  \node[circle, draw, fill=black, inner sep=0, minimum size=3pt] at (dotL|-P) (dotP) {};
  \node[circle, draw, fill=black, inner sep=0, minimum size=3pt] at (dotM|-Q) (dotQ) {};
  \draw[ar] (P) to (dotP);
  \draw[ar] (Q) to (dotQ);
  \draw[ar] (dotP) to (fo_in1);
  \draw[ar] (dotP) to[out=60, in=0] (fi_out1);
  \draw[ar] (dotQ) to (go_in1);
  \draw[ar] (dotQ) to[out=60, in=0] (gi_out1);
  %
  \draw[ar] (dotP) to[out=60, in=0] (V_out1);
  \draw[ar] (dotL) to[out=60, in=0] (V_out2);
  \draw[ar] (dotMI) to[out=120, in=0] (V_out3);
  \draw[ar] (dotQ) to[out=60, in=0] (W_out1);
  \draw[ar] (dotM) to[out=60, in=0] (W_out2);
  \draw[ar] (dotNI) to[out=120, in=0] (W_out3);
  %
  \draw[ar] (V_in1) to[out=180, in=0] (plus_out2);
  \draw[ar] (W_in1) to[out=180, in=0] (plus_out1);
  \coordinate (R) at ($(plus.west)+(-1,0)$);
  \node[left=0 of R] {$\rr$};
  \draw[ar] (plus_in1) to (R);
\end{tikzpicture}
\]
\fi


\chapter{From lenses to dynamics}\label{ch.potlens_to_org}

This chapter defines the functor $\Phi\colon\potlens\to\org$ by applying the abstract representation of \cref{sec.abstract_representation} to polynomial functors. The key infrastructure --- that $\Cat{Para}$ over the store action of $\Cat{Set}_{\cong}$ on $\poly$ presents $\org$ --- was set up in \cref{sec.org_as_para,cor.cotangent_learners}. The remaining inputs are the cotangent functor, the potential monad, and the PN Lie Algebra parameterization.

\section{The cotangent-state refinement}\label{sec.cotangent_state_refinement}

\begin{remark}[Tangent-state alternative]\label{rmk.tangent_state}
Taking $T$ instead of $T^*$ (\cref{prop.TT_endofunctors}) fits the same framework. The analogous \emph{tangent projection} $\tilde\rho_V\colon\cot{TV}\to\cot{V}$ has position-action $(v,X)\mapsto v$ and direction-action $\xi_V\mapsto(\xi_V,\,\flat_V(X))$, using the flat map $\flat_V\colon V\to V^*$ to place $X$ in the second summand of $T^*_{(v,X)}(TV)\cong V^*\oplus V^*$. This produces a functor $\potlens\to\org$ with state set $TV$ (Lagrangian-style, second-order), which we don't develop here; the preference for $T^*$ is physics (cotangent bundles carry the canonical symplectic pairing of \cref{prop.canonical_symplectic_pairing}, recovering Hamilton's equations).
\end{remark}

\begin{remark}[Legendre equivalence of Hamiltonian and Lagrangian $\Phi$]\label{rmk.legendre_equivalence}
By \cref{prop.legendre}, $\sharp\colon T^*\Rightarrow T$ is a monoidal natural isomorphism, so $\rho$ and $\tilde\rho$ correspond via $\cot{\sharp}$. The two functors $\Phi^{T^*},\Phi^{T}\colon\potlens\to\org$ are naturally isomorphic: $\sharp_V\colon T^*V\to TV$ induces a 2-cell in $\org$ identifying $\Phi^{T^*}(f)$ with $\Phi^{T}(f)$. This is the \emph{Legendre equivalence} of Hamiltonian and Lagrangian mechanics in our setting.
\end{remark}

\section{Functoriality and lax monoidality}\label{sec.functoriality}

\begin{proposition}\label{prop.potlens_monoidal}
$\potlens$ inherits from $\lmfd$ and $\pvect$ a symmetric monoidal structure with unit $\binom{\rr^0}{\rr^0}$ and tensor $\lensob L\otimes\lensob{L'}\coloneqq\lensobs{L}{L'}$. On morphisms, $f\otimes f'$ has parameter $V\oplus V'$ (with pairing $\sharp_V\oplus\sharp_{V'}$), independent forward and backward components, and potential $U^\otimes=U\circ\pr+U'\circ\pr'$.
\end{proposition}

\begin{proof}
Para over a symmetric monoidal action of an SMC on an SMC is itself symmetric monoidal. Here $\pvect$ is SMC (\cref{prop.pnla_monoidal}) and $\lmfd^\rr$ is SMC via lensification of $\mfd$ together with the writer monad on $\rr$; the $\pvect$-action on $\lmfd^\rr$ is symmetric monoidal via the forgetful $\pvect\to\mfd$ (\cref{prop.pnla_smc_action}). The morphism formula reads off from the Para tensor: parameter $V\oplus V'$, slotwise forward/backward, additive potential.\qedhere
\end{proof}

\begin{remark}[Lagrangian formulation]\label{rmk.lagrangian}
The Lagrangian formulation is the same dynamics written on $TV$ rather than on $T^*V$. A point of $T^*V$ is a position together with a momentum, while a point of $TV$ is a position together with a velocity. The component $\sharp_V\colon T^*V\To{\cong}TV$ of \cref{prop.legendre} converts momentum to velocity, with inverse $\flat_V$, so $X_H$ transports to a vector field on $TV$:
\[
Y_H\colon
TV\To{\flat_V}T^*V\To{X_H}T(T^*V)\To{T(\sharp_V)}T(TV).
\]
In coordinates $(q,w)\in TV$ with $\xi\coloneqq\flat_V(w)$, the integral curves of $Y_H$ satisfy
\begin{equation}\label{eqn.lagrangian_pushforward}
\dot q=\partial_\xi H|_{(q,\xi)},\qquad \dot w=-\sharp_V\partial_v H|_{(q,\xi)}.
\end{equation}
For an integral curve $\gamma(t)=(q(t),w(t))$ in $TV$, its tangent is a point
\[
\dot\gamma(t)=((q,w),(\dot q,\dot w))\in T(TV)\cong TV\oplus TV.
\]
Thus a vector field on $TV$ is a second-order equation for $q$ exactly when its first tangent component is $w$.

When the pairing on $V$ is symmetric and the Hamiltonian $H\colon T^*V\to\rr$ is of the form
\[
H(v,\xi)=\tfrac{1}{2}\xi(\sharp_V\xi)+U(v),
\]
where $U\colon V\to\rr$ (the \emph{potential}), then\textbf{[decide on notation: $\xi'$ vs $\xi_V'$ vs $\zeta$ for the differential direction; same space as $\xi$, different role]}
\[
d_\xi\left(\tfrac{1}{2}\xi(\sharp_V\xi)\right)(\xi')
=\tfrac{1}{2}\bigl(\xi'(\sharp_V\xi)+\xi(\sharp_V\xi')\bigr)
=\xi'(\sharp_V\xi)
\]
for every $\xi'\in V^*$. Thus, under $(V^*)^*\cong V$, $\partial_\xi H=\sharp_V\xi$, while $\partial_v H=dU$, so \eqref{eqn.lagrangian_pushforward} becomes
\[
\dot q=w,\qquad \dot w=-\sharp_V dU(q).
\]
The first equation identifies $w$ with the velocity of $q$; the second says $\ddot q=\dot w=-\sharp_V dU(q)$, equivalently $\flat_V(\ddot q)=-dU(q)$. This is Newton's second law in this notation: $\flat_V$ converts acceleration to the corresponding force covector, and $-dU(q)$ is the force determined by the potential.

The associated Lagrangian function is obtained from the same Hamiltonian by rewriting momentum as velocity:
\[
L(q,w)\coloneqq \flat_V(w)(w)-H(q,\flat_V(w))
=\tfrac{1}{2}\flat_V(w)(w)-U(q).
\]
The Euler--Lagrange equation for this $L$,
\[
\frac{d}{dt}\partial_w L=\partial_q L,
\]
is exactly $\flat_V(\ddot q)=-dU(q)$, so the Lagrangian has recovered the same second-order equation without using momentum as a state variable.
In this special case one often says ``kinetic plus potential'' for $H$ and ``kinetic minus potential'' for $L$; these names only describe this separable example. The main point is the change of state space: Hamiltonian dynamics uses $(q,\xi)\in T^*V$, while Lagrangian dynamics uses $(q,w)\in TV$, and $\sharp_V$ identifies the two descriptions.\qedhere
\end{remark}



\printbibliography

\end{document}